

\medterm{21-Hydroxylase Deficiency} An inherited condition in which the adrenal glands cannot make enough of certain hormones. It can cause salt loss, low blood pressure, and excess male-type hormones, and is the most common cause of congenital adrenal hyperplasia.

\medterm{Acute Pulmonary Histoplasmosis} An acute lower respiratory tract infection caused by inhaling microconidia of the dimorphic fungus \textit{Histoplasma capsulatum}, commonly found in soil enriched with bird or bat guano. Manifestations range from mild flu-like symptoms to severe pneumonia with mediastinal lymphadenopathy.

\synonyms
Primary Pulmonary Histoplasmosis; Acute Histoplasmosis Pneumonia; Cave Disease

\medterm{Adult Hypertension} A chronic cardiovascular condition defined by sustained elevation of systemic arterial blood pressure ($\ge 130/80\text{ mmHg}$ in adults). A major modifiable risk factor for myocardial infarction, heart failure, stroke, and chronic kidney disease.

\synonyms
Essential Hypertension; Primary Hypertension; Systemic Hypertension

\medterm{Biological Half-Life} The time needed for the amount of a medicine in the blood to fall by half. A medicine usually takes several half-lives to reach a stable level after regular dosing or to leave the body after it is stopped.

\synonyms
T1/2

\medterm{Cardiac Electrophysiology Study} An electrophysiology study in which thin tubes are placed through a blood vessel into the heart to record and stimulate its electrical activity. It helps find the cause of some abnormal rhythms or unexplained fainting.

\medterm{Combined Electrolyte Imbalances} Abnormal blood levels of calcium, phosphate, and potassium. This combination does not have one single cause; each result must be confirmed and interpreted with kidney function, medicines, symptoms, and the rest of the electrolyte panel.

\medterm{Congenital Hemolytic Jaundice} Jaundice present from or caused by an inherited condition that increases red-blood-cell destruction, producing excess unconjugated bilirubin; examples include hereditary spherocytosis and some enzymatic defects.

\medterm{Developmental History} A record of a child's developmental milestones and abilities, including motor, language, cognitive, social, and adaptive skills, used to identify possible delays or regression.

\medterm{Dorsocervical Fat Pad} A visible accumulation of fat at the back of the neck and upper shoulders, associated with obesity, corticosteroid excess, some medicines, or lipodystrophy; it differs from spinal kyphosis.

\medterm{Extremity Hardware Removal} Surgical extraction of orthopedic implants such as plates, screws, intramedullary rods, or wires from an arm or leg. Typically indicated for persistent localized pain, implant infection, mechanical irritation, soft-tissue impingement, nonunion, or planned subsequent procedures.

\synonyms
Orthopedic Hardware Removal; Implant Removal; Plate and Screw Removal

\medterm{Geriatric Hospitalization} Hospital care for an older adult, with attention to medical illness as well as mobility, medicines, nutrition, hearing, vision, and cognition. Preventing falls, delirium, pressure injuries, and loss of independence supports safer recovery.

\synonyms
Inpatient Geriatric Care; Acute Care for Elders (ACE)

\medterm{H And E Staining} A standard laboratory stain using hematoxylin and eosin to show cell nuclei, cytoplasm, and tissue structure under a microscope. It is widely used to examine biopsies and identify inflammation, infection, or tumors.

\medterm{H And H} Abbreviation for hemoglobin and hematocrit, two blood measurements commonly used to assess red-blood-cell levels and anemia.

\medterm{H And P} Abbreviation for a patient’s history and physical examination. The history covers symptoms, past health, medicines, and relevant background; the examination checks the body for signs of illness.

\medterm{H Influenzae Meningitis} A serious infection of the coverings of the brain and spinal cord caused by Haemophilus influenzae bacteria. It may cause fever, headache, stiff neck, confusion, seizures, or brain injury, especially in unvaccinated children.

\medterm{H1N1 Flu} Alternate name; see \textit{H1N1 Influenza}.

\medterm{H1N1 Influenza} Influenza A infection caused by an H1N1 virus. It spreads through respiratory droplets and commonly causes sudden fever, cough, sore throat, muscle aches, headache, and tiredness; some people develop pneumonia or severe breathing problems.

\medterm{H1N2} A subtype of influenza A virus that can infect humans and pigs. Human infections are uncommon and usually cause an influenza-like respiratory illness; surveillance and testing help distinguish it from other influenza viruses.

\medterm{H2 Blockers} Medicines that reduce stomach-acid production. They can relieve heartburn, indigestion, and some ulcers; possible adverse effects include headache, diarrhea, constipation, and interactions with other medicines.

\medterm{H2 Receptor Antagonists Overdose} Taking too much of an H2 blocker medicine. It may cause drowsiness, confusion, agitation, abnormal heartbeat, or rarely seizures; urgent poison-control or emergency advice is needed after a significant overdose.

\medterm{H5N1} A subtype of avian influenza A virus that mainly infects birds but can occasionally infect mammals and humans after close contact with infected animals or contaminated environments. Human illness can be severe and requires public-health evaluation.

\medterm{Habenula} A small brain structure helping process reward, unpleasant experiences, stress, motivation, and learned behavior. It connects several brain regions and may influence mood, sleep, pain, and responses to negative outcomes.

\medterm{Habit} A repeated behavior or pattern of action that may be ordinary, helpful, or harmful; in clinical assessment, its frequency, control, and effect on functioning determine whether it is a concern.

\medterm{Habituation} A gradual reduction in response to a repeated, harmless stimulus. For example, a person may stop noticing a constant sound; habituation differs from addiction or drug tolerance.

\medterm{Habitus} A person’s general body build and appearance as observed during a medical examination. Clinicians may describe body habitus when noting nutritional state, muscle mass, or features relevant to a condition.

\medterm{Haem} The iron-containing part of hemoglobin that binds oxygen inside red blood cells. Its structure allows blood to carry oxygen from the lungs to tissues and return carbon dioxide.

\medterm{Haemangioma} A usually harmless growth made of extra blood vessels. It is common in babies and may appear as a red or purple mark or lump.

\medterm{Haemarthrosis} Bleeding into a joint, causing swelling, pain, and reduced movement. It may follow an injury or occur with hemophilia, blood-thinning medicines, surgery, or certain joint disorders.

\medterm{Haematemesis} Vomiting blood or material that looks like coffee grounds, usually from bleeding in the esophagus, stomach, or upper small intestine. It can be life-threatening and requires emergency medical assessment.

\medterm{Haematinics} Nutrients or medicines needed to make healthy red blood cells, especially iron, vitamin B12, and folate. Deficiency can cause anemia, although treatment depends on the specific cause.

\medterm{Haematocele} A collection of blood in a body cavity or tissue, often in the scrotum after injury, surgery, or bleeding. It may cause swelling, pain, or pressure and sometimes needs imaging or drainage.

\medterm{Haematocolpos} Accumulation of menstrual blood in the vagina, usually because the vaginal opening is blocked or too narrow. It can cause absent periods, monthly pelvic pain, and a swollen vagina or uterus.

\medterm{Haematocrit} The proportion of blood volume occupied by red blood cells, expressed as a percentage or decimal fraction. The reference range varies with age, sex, altitude, pregnancy, and laboratory method.

\medterm{Haematogenous} Spread through the bloodstream. An infection, cancer, or other material that travels this way can reach organs or tissues far from its original site.

\medterm{Haematologist} A medical doctor who specializes in the diagnosis, treatment, and prevention of diseases
of the blood and blood-forming organs.

\medterm{Haematoma} A collection of blood outside a blood vessel, usually caused by injury, surgery, bleeding disorders, or blood-thinning medicines. A hematoma may cause swelling and pain and can compress nearby organs or nerves.

\medterm{Haematometra} A collection of menstrual blood in the uterus, usually because the cervix or vagina is blocked or absent. It may cause missed periods, recurring pelvic pain, and uterine enlargement.

\medterm{Haematomyelia} Bleeding into the spinal cord, usually after severe injury or because of an abnormal blood vessel or bleeding disorder. It can cause sudden back pain, weakness, numbness, paralysis, or loss of bladder control.

\medterm{Haematopericardium} Blood within the sac around the heart, usually after chest injury, heart surgery, or a tear in the heart or a blood vessel. It can compress the heart and cause life-threatening tamponade.

\medterm{Haematopoiesis} The process of making blood cells, mainly inside the bone marrow. It produces red cells that carry oxygen, white cells that fight infection, and platelets that help stop bleeding.

\medterm{Haematoporphyria} An older term for a disorder involving porphyrins, substances used to make hemoglobin. Porphyrin disorders may cause sensitivity to sunlight, abdominal or nerve symptoms, or red-brown urine, depending on the type.

\medterm{Haematosalpinx} Blood collected inside a fallopian tube, often from ectopic pregnancy, endometriosis, pelvic injury, or blockage. It may cause pelvic pain and can affect fertility or require urgent treatment.

\medterm{Haematuria} Blood in the urine. It may be visible or found only with a urine test and can have many causes, including infection, stones, medicines, or kidney disease.

\medterm{Haemochromatosis} A disorder in which the body absorbs and stores too much iron. Over time, iron can damage the liver, heart, pancreas, joints, skin, and reproductive organs; early treatment can prevent complications.

\medterm{Haemodialysis} A treatment for severe kidney failure in which a machine filters waste products and extra fluid from the blood. It may be needed temporarily or regularly when the kidneys cannot do enough of this work.

\medterm{Haemoflagellate} A parasite with a whip-like structure that lives in blood or tissues. Examples include the parasites that cause sleeping sickness and Chagas disease.

\medterm{Haemoglobin} The iron-containing protein inside red blood cells that carries oxygen from the lungs to tissues and helps return carbon dioxide to the lungs. A low level is anemia; an abnormal form can cause disease.

\medterm{Haemoglobinopathy} An inherited disorder affecting the amount or structure of hemoglobin, the oxygen-carrying protein in red blood cells. Examples include sickle-cell disease and thalassemia, which can cause anemia or other complications.

\medterm{Haemoglobinuria} Free hemoglobin in the urine, often released when red blood cells break apart in the bloodstream. It can make urine red or brown and may occur with infections, medicines, burns, or blood disorders.

\medterm{Haemoglobinuria} Alternate British spelling; see \textit{Hemoglobinuria}.

\medterm{Haemolysis} The breakdown of red blood cells, which releases hemoglobin into the blood. It can happen normally as old cells are removed or too quickly because of disease.

\medterm{Haemolytic Anaemia} Anemia caused by red blood cells being destroyed faster than the bone marrow can replace them. It may cause tiredness, jaundice, dark urine, rapid heartbeat, breathlessness, or an enlarged spleen.

\medterm{Haemolytic Disease Of The Newborn} Anemia and jaundice in a newborn caused by antibodies from the pregnant person crossing the placenta and destroying the baby’s red blood cells. Severe cases can cause brain injury without prompt treatment.

\medterm{Haemoperitoneum} Bleeding into the abdominal cavity, commonly after injury, a ruptured ectopic pregnancy, or a bleeding liver, spleen, artery, or other organ. It can cause shock and requires urgent assessment.

\medterm{Haemophilia} Alternate spelling; see \textit{Hemophilia}.

\medterm{Haemophilia A} Hemophilia A, an inherited bleeding disorder caused by deficiency or dysfunction of clotting factor VIII. It can cause prolonged bleeding, easy bruising, and bleeding into joints or muscles; treatment replaces or stimulates factor VIII activity.

\medterm{Haemophilus} A group of bacteria that may live harmlessly in the nose and throat or cause infection. Depending on the species, illness can affect the ears, sinuses, lungs, blood, joints, or brain coverings.

\medterm{Haemophilus Aegyptius} A bacterium associated with eye infections and, rarely, severe bloodstream disease. It can spread through respiratory secretions or direct contact, and laboratory identification helps guide treatment and public-health precautions.

\medterm{Haemophilus Haemolyticus} A bacterium that commonly lives harmlessly in the upper airway but can occasionally cause respiratory or invasive infection, especially in people with weakened immunity or serious underlying illness.

\medterm{Haemophilus Infections} Infections caused by Haemophilus bacteria, which may affect the ears, sinuses, lungs, bloodstream, joints, or brain coverings. Severity depends on the species, infection site, age, vaccination, and immune health.

\medterm{Haemophilus Influenzae} A type of bacteria that can cause ear infections, pneumonia, meningitis, and other infections. Type b, or Hib, can cause severe disease in children, but vaccination greatly reduces this risk.

\medterm{Haemophilus Parainfluenzae} A small gram-negative bacterium that normally colonizes the upper respiratory tract but can rarely cause respiratory, bloodstream, or other opportunistic infections.

\medterm{Haemoptysis} Coughing up blood or blood-stained mucus from the airways or lungs. Causes include infection, bronchiectasis, a blood clot, cancer, or irritation; coughing up more than a small streak, or having breathlessness or chest pain, needs urgent medical assessment.

\medterm{Haemorrhage} British spelling of Hemorrhage; see Hemorrhage.

\medterm{Haemorrhagic Fever} A group of severe infections caused by certain viruses that can damage blood vessels and disrupt clotting. Fever, weakness, bleeding, shock, and organ failure may occur; urgent isolation and specialist care may be needed.

\medterm{Haemorrhagic Fever With Renal Syndrome} A hantavirus infection usually acquired by inhaling dust contaminated with infected rodent urine or droppings. It causes fever, headache, bleeding, low blood pressure, and temporary kidney failure and may be severe.

\medterm{Haemorrhagic Stroke} A stroke caused by a burst blood vessel in the brain or in the space around it. The bleeding injures brain tissue and may cause sudden weakness, speech trouble, severe headache, vomiting, or loss of consciousness.

\medterm{Haemorrhoids} Alternate spelling; see \textit{Hemorrhoids}.

\medterm{Haemosiderin} A brownish iron-storage material that forms inside cells when iron accumulates, often after repeated bleeding, transfusions, or excess iron absorption. Large deposits can contribute to tissue damage.

\medterm{Haemosiderosis} Excess storage of iron-containing haemosiderin in tissues, often after repeated transfusions, bleeding into the lungs, or increased iron absorption. Effects depend on the organs involved and may include organ damage.

\medterm{Haemostasis} The body’s normal process for stopping bleeding. Blood vessels narrow, platelets form a temporary plug, and clotting proteins strengthen it with a stable clot.

\medterm{Haemostat} A substance, device, or instrument used to control bleeding. It may work by pressing on a vessel, sealing tissue, promoting clotting, or temporarily blocking blood flow during an operation.

\medterm{Haemothorax} Blood collecting between the lung and chest wall, usually after chest injury, surgery, or a medical procedure. It can compress the lung, cause breathing difficulty or shock, and may require drainage.

\medterm{Haff Disease} Sudden muscle breakdown and severe muscle pain after eating certain contaminated fish or seafood. It can release muscle proteins that damage the kidneys; urgent treatment focuses on fluids and monitoring for complications.

\medterm{Hafnia} A genus of gram-negative intestinal bacteria; Hafnia alvei is the species most often identified in human specimens.

\medterm{Hafnia Alvei} A Gram-negative bacterium found in the intestine and environment that rarely causes opportunistic infection, usually in people with serious illness or weakened immunity.

\medterm{Hageman Factor} An older name for factor XII, a blood protein that helps start one pathway of clotting. Factor XII deficiency can prolong a laboratory clotting test but usually does not cause abnormal bleeding.

\medterm{Hailey-Hailey Disease} An inherited skin disorder that causes repeated blisters, raw areas, and painful cracks, especially in skin folds such as the neck, armpits, and groin. Heat, sweating, friction, and infection can trigger flares.

\medterm{Hair} Keratinized filaments produced by follicles in the skin. Hair protects the scalp, contributes to temperature regulation and sensation, and grows through anagen, catagen, and telogen phases.

\medterm{Hair Bleach Poisoning} Illness after swallowing or inhaling hair-bleaching chemicals, which may contain peroxide or persulfates. They can irritate or burn the mouth, throat, skin, eyes, or lungs; contact poison-control services promptly.

\medterm{Hair Bulb} The rounded base of a hair follicle where new hair is produced. It contains growing cells and a small blood-supplied area that provides nutrients and helps determine hair growth and color.

\medterm{Hair Cells} Special cells in the inner ear that change sound vibrations and head movement into nerve signals for hearing and balance.

\medterm{Hair Disease} A disorder affecting hair growth, the hair shaft, or the scalp; causes include genetics, inflammation, infection, medicines, and nutritional problems.

\medterm{Hair Dye Poisoning} Harm caused by swallowing or inhaling hair-dye chemicals. Depending on the ingredients and amount, symptoms may include vomiting, swelling, breathing difficulty, allergic reaction, kidney injury, or muscle damage and require urgent advice.

\medterm{Hair Follicle} A small tube in the skin that produces a hair. Each follicle has a growth phase, a brief transition phase, and a resting phase.

\medterm{Hair Loss} Loss or thinning of hair on the scalp or body. It may be temporary or permanent, patchy or widespread, and caused by genetics, illness, medicines, autoimmune disease, infection, or scarring of the scalp.

\synonyms
Alopecia; Effluvium; Baldness

\medterm{Hair Loss} The partial or complete loss of hair from areas where it normally grows, most commonly
the scalp. Causes include genetic predisposition (androgenetic alopecia), autoimmune
conditions (alopecia areata), hormonal changes, medications, stress, and nutritional
deficiencies.

\synonyms
Alopecia; Baldness; Effluvium

\medterm{Hair Loss Classification} Hair loss is classified as scarring or nonscarring. Scarring hair loss destroys follicles and is usually permanent; nonscarring loss leaves follicles intact and may improve when its cause, such as illness, medicines, stress, or autoimmune disease, is treated.

\medterm{Hair Removal} Removal of unwanted hair by shaving, trimming, waxing, depilatory products, electrolysis, or laser treatment. Possible problems include irritation, burns, ingrown hairs, pigment changes, and infection.

\medterm{Hair Shaft} The part of a hair that extends above the skin. It is made of dead, hardened cells produced by the follicle and can be affected by cutting, heat, chemicals, friction, or damage to the hair root.

\medterm{Hair Spray Poisoning} Illness after swallowing or inhaling hair spray. Small accidental exposures may irritate the mouth or lungs, while larger exposures can cause vomiting, drowsiness, breathing problems, or chemical injury; contact poison-control services.

\medterm{Hair Straightener Poisoning} Harm after swallowing a hair-straightening product, which may contain corrosive chemicals. It can burn the mouth, throat, esophagus, or stomach and cause breathing or swallowing problems; do not induce vomiting and seek emergency advice.

\medterm{Hair Tonic Poisoning} Illness after swallowing or inhaling a hair-tonic product. Ingredients such as alcohols or solvents may cause vomiting, drowsiness, low blood sugar, or breathing problems; the product and amount determine emergency treatment.

\medterm{Hair Transplant} A procedure that moves hair follicles, usually from a donor area of the scalp, to areas affected by permanent hair loss; results depend on diagnosis, donor supply, technique, and postoperative care.

\medterm{Hair-Ball} A hair-ball, or trichobezoar, is a mass of swallowed hair in the stomach or intestine, usually associated with repeated hair eating.

\medterm{Hairy Cell Leukaemia} Hairy cell leukaemia is a rare slow-growing blood cancer in which abnormal B lymphocytes accumulate in the bone marrow, spleen, and blood.

\medterm{Hairy Cell Leukemia} A rare, usually slow-growing cancer of a type of white blood cell. Abnormal cells build up in the bone marrow, blood, and spleen, causing low blood counts, infections, tiredness, or an enlarged spleen.

\synonyms
HCL; Leukemic Reticuloendotheliosis

\medterm{Hairy Tongue} A usually harmless overgrowth and discoloration of the tongue's papillae, often appearing black, brown,
or white. It may be associated with smoking, poor oral hygiene, dry mouth, or antibiotics;
gentle tongue cleaning and removal of contributing factors usually help.

\medterm{Hajdu-Cheney Syndrome} A rare inherited bone disorder causing fragile bones, short stature, abnormal skull and facial development, and sometimes loose joints or dental problems. Severity varies, and specialist monitoring helps prevent complications.

\medterm{Halcinonide} A potent topical corticosteroid used to reduce inflammation and itching in steroid-responsive dermatitis and other inflammatory skin disorders; prolonged or extensive use can cause skin atrophy and systemic effects.

\medterm{Haldane Effect} The effect of oxygen binding to hemoglobin in the lungs, which helps release carbon dioxide for exhalation. In body tissues, hemoglobin without oxygen carries more carbon dioxide back toward the lungs.

\medterm{Half-Life} The time it takes for the amount of a medicine in the body, or its blood level, to decrease by half. It helps determine how often a medicine should be taken and how long it remains in the body.

\medterm{Halitosis} Persistent unpleasant breath odor, most often from bacteria on the tongue, teeth, or gums. Dry mouth, smoking, foods, sinus or throat disease, reflux, and some illnesses or medicines can also contribute.

\medterm{Hallermann's Syndrome} A rare congenital disorder characterized by distinctive craniofacial development, sparse hair, dental abnormalities, and sometimes eye or skeletal findings; it is also called Hallermann--Streiff syndrome.

\medterm{Hallervorden-Spatz Disease} Alternate historical term; see \textit{Neurodegeneration With Brain Iron Accumulation}.

\medterm{Hallucination} A sensory experience that feels real but occurs without a matching external stimulus. It may involve seeing, hearing, feeling, smelling, or tasting something that is not present.

\medterm{Hallucinations} Hallucinations are perceptions without an external sensory stimulus and can arise from psychiatric illness, neurologic disease, delirium, substances, medications, sensory loss, or sleep transitions.

\medterm{Hallucinogen} A psychoactive substance that can alter perception, thought, and mood and may produce hallucinations; effects and toxicity vary by substance, dose, co-ingestants, and mental state.

\medterm{Hallux} The great toe, or big toe. It has two bones and helps the foot stay stable and push off during walking. Common problems include hallux valgus (a bunion), hallux rigidus (a stiff, arthritic joint), turf toe (a sprain), and gout causing sudden pain and swelling at the joint where the toe meets the foot.

\medterm{Hallux Rigidus} Arthritis of the joint where the big toe meets the foot, causing pain, stiffness, swelling, and reduced upward movement. A bony bump may develop and make walking or wearing shoes difficult.

\medterm{Hallux Valgus} A deformity in which the big toe angles toward the other toes and the joint at its base becomes prominent on the inner side of the foot. It may cause pain, rubbing, swelling, and difficulty wearing shoes.

\medterm{Hallux Valgus} A bunion: a change in the alignment of the joint at the base of the big toe, with the toe pointing inward and a prominent area on the inner forefoot. It may become painful from pressure or inflammation.

\synonyms
Bunion; Hallux Abductovalgus

\medterm{Hallux Varus} Abnormal medial deviation of the great toe, often following surgery, trauma, neuromuscular imbalance, or congenital deformity; it can cause shoe difficulty and pressure-related pain.

\medterm{Halo Nevus} A usually harmless mole surrounded by a ring of lighter skin. The immune system causes the mole and nearby pigment cells to lose color; it often occurs in children or young adults and may slowly fade.

\medterm{Halogen} An element in group 17 of the periodic table, including fluorine, chlorine, bromine, iodine, and astatine; halogens and their compounds have medical, biologic, and toxicologic importance.

\medterm{Haloperidol} Haloperidol is a first-generation antipsychotic that blocks dopamine D2 receptors and is used for psychosis, agitation, delirium, and tics; extrapyramidal reactions, tardive dyskinesia, hyperprolactinaemia, and QT prolongation are important risks.

\medterm{Halothane} A volatile inhaled anesthetic formerly used to induce or maintain general anesthesia; it can cause hypotension, arrhythmias, and rare severe immune-mediated liver injury.

\medterm{Hamamelis} A botanical preparation from witch hazel used topically as an astringent for minor skin or mucosal irritation; evidence and product composition vary.

\medterm{Hamartoma} A usually harmless growth made of an abnormal arrangement of tissue types normally found in that part of the body. It may occur in organs such as the lung, skin, brain, or intestine and may cause problems if it grows or presses on nearby structures.

\medterm{Hamartomatous Polyposis} A condition in which multiple hamartomatous gastrointestinal polyps develop from disorganized but benign tissue overgrowth, sometimes as part of an inherited cancer-predisposition syndrome.

\medterm{Hamate} A small, wedge-shaped bone on the little-finger side of the wrist. Its hook provides attachment for ligaments and tendons; fractures there may cause wrist pain or pressure on a nearby nerve.

\medterm{Hamate Bone} The hamate is a wedge-shaped wrist bone on the little-finger side that has a hook serving as an attachment and passage point for tendons and nerves.

\medterm{Hamman Sign} A crunching or clicking noise heard with a stethoscope over the chest in time with the heartbeat, caused by air trapped in the center of the chest.

\medterm{Hamman's Sign} A crackling or crunching sound heard over the chest in time with the heartbeat, usually caused by air in the central chest around the heart and major vessels. It can occur with pneumomediastinum and needs clinical assessment.

\medterm{Hammer} The common name for the malleus, one of the three small bones of the middle ear.

\medterm{Hammer Toe} A deformity in which a lesser toe bends downward at its middle joint and may press against the shoe. It can cause pain, stiffness, corns, or calluses and may be flexible at first or become fixed.

\medterm{Hammer Toe Repair} Surgery to correct a fixed or painful flexion deformity of a lesser toe, using soft-tissue release, tendon balancing, bone correction, or temporary fixation as indicated.

\medterm{Hammertoe} A deformity in which a toe bends downward at its middle joint, giving it a claw-like
appearance. It may cause pain, stiffness, corns, or calluses, especially when wearing shoes.

\medterm{Hampton Hump} Wedge-shaped opacification with its base abutting the pleura on chest radiography,
representing pulmonary infarction secondary to Pulmonary Embolism. See Pulmonary
Embolism (PE).

\medterm{Hampton's Hump} Alternate spelling; see \textit{Hampton Hump}.

\medterm{Hamstring} One of three muscles forming the posterior compartment of the thigh: semitendinosus,
semimembranosus, and biceps femoris (long and short heads). Common origin: ischial
tuberosity (except short head of biceps femoris which originates from the linea aspera
of the femur).

\medterm{Hamstring Injury} A strain or tear of one or more muscles or tendons at the back of the thigh.

\medterm{Hamstring Muscles} The muscles at the back of the thigh. They bend the knee and help move the hip, especially during walking, running, and jumping. A sudden stretch or forceful contraction can strain or tear them.

\medterm{Hand} The part of the upper limb beyond the wrist, containing the palm, fingers, bones, joints, muscles, nerves, and blood vessels used for grasping and sensation.

\medterm{Hand Deformity} An abnormal shape, alignment, or loss of motion of the hand or fingers caused by congenital difference, trauma, arthritis, nerve dysfunction, tendon disease, or other conditions.

\medterm{Hand Foot And Mouth Disease} A contagious viral illness, most common in young children, that causes fever, painful mouth sores, and a rash or small blisters on the hands, feet, and sometimes buttocks. It spreads through close contact and contaminated secretions or stool.

\medterm{Hand, Foot, And Mouth Disease} A contagious enterovirus infection, usually in children, causing fever, painful mouth sores, and small spots or blisters on the hands, feet, and sometimes buttocks. It usually improves on its own, but dehydration can occur.

\medterm{Hand Hygiene} Cleaning the hands with soap and water or an alcohol-based hand rub to remove or kill germs. It is especially important before and after patient contact, after touching body fluids, and when hands are visibly dirty.

\medterm{Hand Lotion Poisoning} Illness after swallowing hand lotion. Most products cause mild stomach upset, but some contain alcohols or other chemicals that can cause drowsiness, low blood sugar, or breathing problems; contact poison-control services.

\medterm{Hand Microsurgery} Specialized hand surgery performed with magnification and fine instruments to repair or reconstruct small blood vessels, nerves, and other delicate tissues. It may be used for complex injuries, replantation, and restoration of blood flow to a threatened digit or hand.

\medterm{Hand of Benediction} An abnormal hand posture observed when a patient attempts to make a fist, resulting in extension of the index and middle fingers with flexion of the fourth and fifth fingers. It occurs due to proximal median nerve injury, which paralyses the flexor digitorum superficialis and the lateral half of the flexor digitorum profundus.

\synonyms
Benediction Sign; Preacher's Hand

\medterm{Hand Or Foot Spasms} Sudden, involuntary tightening of muscles in a hand or foot. Causes include overuse, dehydration, abnormal calcium or other salts, nerve disease, medicines, and reduced blood flow; repeated or painful spasms need assessment.

\medterm{Hand Pain} Pain or discomfort in the hand, fingers, or wrist caused by injury, arthritis, nerve compression, tendon problems, infection, or other conditions. Sudden severe pain, deformity, numbness, or loss of circulation needs prompt assessment.

\medterm{Hand Prolapse} A fetal hand or arm lying beside or below the part leading into the birth canal. It can interfere with labor or reduce blood flow to the limb; the position requires prompt obstetric assessment and should not be pushed back at home.

\medterm{Hand X-Ray} Radiographic imaging of the hand used to evaluate fractures, dislocations, arthritis, infection, tumors, bone age, and other abnormalities of the bones and joints.

\medterm{Hand-Arm Vibration Syndrome} An occupational disorder from prolonged exposure to vibrating tools, causing numbness, tingling, reduced grip, and episodic whitening of the fingers from vascular spasm.

\medterm{Hand-Foot Syndrome} A dose-limiting toxic cutaneous reaction triggered by certain systemic antineoplastic and targeted chemotherapy agents (notably capecitabine, 5-fluorouracil, pegylated liposomal doxorubicin, cytarabine, and multi-kinase inhibitors like sorafenib and sunitinib). It is characterized by progressive bilateral burning dysesthesia, tingling, painful symmetric erythema, swelling, and subsequent blistering or desquamation primarily localized to the friction-bearing pressure areas of the palms and soles.
\synonyms{Palmar-Plantar Erythrodysesthesia; PPE; Chemotherapy-Induced Acral Erythema; Burgdorf Reaction}

\medterm{Hand-Foot-And-Mouth Disease} A contagious viral illness causing sores in the mouth and a rash on the hands and feet, usually in children.

\medterm{Hand-Foot-And-Mouth Syndrome} A usually mild, contagious viral illness, commonly caused by enteroviruses, producing fever, mouth sores,
and small spots or blisters on the hands and feet. It mainly affects young children but can occur in adults. Hydration,
pain relief, and hygiene are important; severe dehydration, neurologic symptoms, or persistent fever requires care.

\medterm{Hand-Foot-Mouth Disease} A common enteroviral infection, usually in young children, causing fever, painful oral
ulcers, and a vesicular rash on the hands, feet, and sometimes buttocks. It is highly
contagious and usually self-limited.

\medterm{Handling Sharps And Needles} Handling sharps and needles means using, passing, disposing of, and storing sharp medical instruments safely to prevent needlestick injuries and infection transmission.

\medterm{Hand-SchüLler-Christian Disease} Hand-Schüller-Christian disease is an older name for a multisystem form of Langerhans cell histiocytosis, often involving bone, skin, and the pituitary gland.

\medterm{Handwashing} Handwashing is cleaning the hands with soap and water or an alcohol-based product to remove germs and reduce transmission of infection.

\medterm{Hanging And Strangulation Signs} Injuries or other findings caused by pressure around the neck that restricts breathing or blood flow. Findings may include neck marks, swelling, voice changes, trouble swallowing, breathing difficulty, confusion, or loss of consciousness; urgent medical assessment is required.

\medterm{Hang-Nail} A small torn flap of skin beside a fingernail or toenail, often caused by dryness or trauma and vulnerable to local infection if torn or contaminated.

\medterm{Hangover} A group of symptoms after alcohol use, such as headache, nausea, thirst, fatigue, dizziness, and sensitivity to light; symptoms reflect effects of alcohol and its metabolites, dehydration, sleep disruption, and inflammation.

\medterm{Hansen Disease} A long-lasting infection caused by Mycobacterium leprae that mainly affects the skin and peripheral nerves. Pale or reddish numb patches, weakness, and eye problems may occur; multidrug antibiotics cure infection and prevent disability.

\synonyms
Leprosy; Hanseniasis; Mycobacterium Leprae Infection

\medterm{Hantavirus} A group of viruses spread mainly by breathing dust contaminated with infected rodent urine, droppings, or nesting material. They can cause severe lung disease or a fever that damages the kidneys.

\medterm{Hantavirus Infection} A zoonotic viral infection acquired mainly through inhalation of contaminated rodent excreta. It may cause hemorrhagic fever with renal syndrome or hantavirus pulmonary syndrome, which can progress rapidly to respiratory failure.

\medterm{Hantavirus Pulmonary Syndrome} A severe lung infection acquired mainly by inhaling dust contaminated with infected rodent urine, droppings, or nesting material. Fever, muscle aches, headache, and stomach symptoms can rapidly progress to fluid-filled lungs and respiratory failure.

\synonyms
HPS; Sin Nombre Virus Infection

\medterm{Haploid} Having one complete set of chromosomes, as in a human sperm or egg. Haploid cells contain half the chromosome number of most body cells and restore the full number when they join during fertilization.

\medterm{Haploidentical Donor} A donor who shares one human leukocyte antigen haplotype with the recipient, commonly used for selected hematopoietic stem-cell or organ transplantation when a fully matched donor is unavailable.

\medterm{Haploidy} The state of having one complete set of chromosomes, as in human gametes; haploid cells contain 23 chromosomes.

\medterm{Haploinsufficiency} A genetic state in which one functional copy of a gene does not produce enough gene product for normal function, so loss or inactivation of the other copy can cause disease.

\medterm{Haplotype} A group of nearby gene versions inherited together on one chromosome. Haplotype information can help study inherited disease risk, identify tissue donors, and trace how genetic traits are passed through families.

\medterm{Hapten} A hapten is a small molecule that is not immunogenic by itself but can trigger an immune response after binding to a larger carrier protein, as in some drug allergies.

\medterm{Haptoglobin} A blood protein made mainly by the liver that binds free hemoglobin released when red blood cells break down. Low levels can support evaluation of hemolysis, but liver disease and inflammation also affect the result.

\medterm{Haptoglobin Analyses} Laboratory tests measuring haptoglobin, a plasma protein that binds free hemoglobin; low values can support evaluation of intravascular hemolysis but are also affected by liver disease and inflammation.

\medterm{Haptoglobin Blood Test} A blood test measuring haptoglobin, a protein that binds free hemoglobin; low levels may support intravascular hemolysis but can also reflect liver disease or inflammation.

\medterm{Haptoglobin Decreased} A low blood haptoglobin result, commonly caused by intravascular red-cell destruction, reduced hepatic production, or severe liver disease; it is interpreted with hemoglobin, bilirubin, reticulocytes, and a blood smear.

\medterm{Hard Palate} The bony, mucosa-covered front portion of the roof of the mouth that separates the oral and nasal cavities and provides a surface for chewing and speech.

\medterm{Harlequin Color Change} A temporary, harmless change in a newborn’s skin color in which one side becomes red while the other remains pale. It often appears when the baby lies on one side and usually fades within minutes.

\medterm{Harlequin Foetus} An older term for a baby born with harlequin ichthyosis, a severe inherited skin disorder causing thick, hard plates of skin separated by deep cracks. The condition can impair temperature control, movement, feeding, and protection from infection.

\medterm{Harlequin Ichthyosis} A rare severe inherited skin disorder in which newborns have thick plates of skin separated by deep cracks. It can impair temperature control, movement, feeding, breathing, and protection from infection and requires intensive neonatal care.

\medterm{Harm Reduction} Evidence-based strategies that decrease health, social, and legal harms associated with
substance use without requiring immediate abstinence. Examples include naloxone access,
sterile injecting equipment, drug-checking where available, safer-use education, and
linkage to treatment.

\medterm{Harm Reduction Therapy} Care that reduces illness, injury, infection, overdose, and other harms related to substance use or another risky behavior, even when stopping immediately is not possible. It can include safer-use support and treatment for dependence.

\synonyms
Harm Reduction; Risk Mitigation Strategy

\medterm{Harmless Murmur} An innocent or physiologic heart murmur produced by normal blood flow without structural heart disease; it is diagnosed only after appropriate clinical assessment.

\medterm{Harrison's Sulcus} Harrison's sulcus is a horizontal groove along the lower chest wall caused by inward pulling of the ribs, classically in long-standing childhood breathing difficulty.

\medterm{Hartmann Procedure} An operation that removes the diseased lower colon, brings the remaining colon through the abdominal wall as a colostomy, and closes the downstream bowel. It is often used when infection, blockage, perforation, or severe illness makes immediate reconnection unsafe.

\medterm{Hartnup Disease} A rare inherited problem absorbing and reusing certain amino acids. It may cause a sun-sensitive rash, poor coordination, or mood and thinking changes, although some people have no symptoms.

\medterm{Hartnup Disorder} An inherited disorder of neutral amino-acid transport that can cause pellagra-like skin changes, neurologic symptoms, or no symptoms.

\medterm{HAS-BLED Score} A score estimating bleeding risk in a person with atrial fibrillation who may need a blood thinner. It considers high blood pressure, kidney or liver disease, previous stroke or bleeding, unstable clotting results, age, medicines, and alcohol; a high score prompts closer monitoring and correction of changeable risks.

\synonyms
HAS-BLED; Anticoagulant Bleeding Risk Index

\medterm{Hashimoto Disease} An autoimmune disorder in which antibodies and lymphocytes progressively damage the thyroid, commonly causing hypothyroidism, goiter, fatigue, cold intolerance, constipation, and weight change.

\medterm{Hashimoto Thyroiditis} An autoimmune disorder in which immune-mediated inflammation gradually damages the thyroid gland and may cause
hypothyroidism. The gland may be enlarged or atrophic, and thyroid autoantibodies are
often present.

\medterm{Hashimoto's Disease} An autoimmune disorder in which the immune system gradually damages the thyroid, often causing low thyroid-hormone levels. It may produce tiredness, cold intolerance, weight change, constipation, or an enlarged thyroid.
\synonyms
Hashimoto Thyroiditis; Chronic Autoimmune Thyroiditis

\medterm{Hashimoto's Thyroiditis} A chronic autoimmune inflammation of the thyroid that commonly causes progressive hypothyroidism and may produce a firm, enlarged thyroid gland.

\medterm{Hashish} A concentrated form of cannabis made from resin of the cannabis plant. It contains cannabinoids that can impair memory, coordination, and judgment and may cause anxiety, dependence, or psychosis in susceptible people.

\medterm{Hashitoxicosis} A temporary period of excess thyroid hormone caused by release of stored hormone during Hashimoto thyroiditis.

\medterm{Hassall's Corpuscles} Hassall's corpuscles are rounded structures in the thymus made of specialized epithelial cells and involved in the development of immune tolerance.

\medterm{Haustra} The series of small pouches along the outside of the large intestine. They form as the bowel muscles contract and help the colon mix and move its contents.

\medterm{Haverhill Fever} An illness caused by Streptobacillus bacteria, usually after drinking contaminated milk or water. It may cause fever, vomiting, headache, muscle or joint pain, and a rash and requires antibiotic treatment.

\medterm{Hay Fever} A common name for seasonal allergic rhinitis, an allergic response to airborne outdoor plant pollens and mold spores that causes sneezing, a blocked or runny nose, an itchy palate, and watery, irritated eyes. See also \textit{Seasonal Allergies}.
\synonyms
Seasonal Allergies; Seasonal Allergic Rhinitis; Pollinosis

\medterm{Hazardous Materials} Substances that can cause injury, poisoning, fire, radiation exposure, or environmental harm; safe handling requires labeling, protective equipment, exposure control, and emergency procedures.

\medterm{Hazardous Substance} A chemical, biological agent, or other material that can harm health through swallowing, inhalation, skin contact, injection, radiation, or environmental exposure. Risk depends on its properties, amount, route, and duration of exposure.

\medterm{Hb} Hemoglobin, the iron-containing protein in red blood cells that carries oxygen and contributes to carbon-dioxide transport and blood's acid--base buffering.

\medterm{HbA1c} Abbreviation; see \textit{Glycated Hemoglobin}.
\synonyms
Glycated Hemoglobin; Hemoglobin A1c; A1c Test

\medterm{HCG} Human chorionic gonadotropin is a hormone produced mainly during pregnancy; its measurement is used in pregnancy testing and evaluation of selected tumors.

\medterm{HCG In Urine} A urine test detecting human chorionic gonadotropin, commonly used as a pregnancy test; accuracy depends on timing, urine concentration, and test performance.

\medterm{Hct} The percentage of blood volume occupied by red blood cells, measured as part of a complete blood count.

\medterm{HCV PCR} A polymerase chain reaction test that detects or measures hepatitis C virus RNA in blood, helping identify active infection and monitor treatment response.

\medterm{HCV-Associated Rheumatic Disease} Joint, muscle, blood-vessel, or autoimmune problems associated with hepatitis C virus infection. Possible manifestations include joint pain or inflammation, mixed cryoglobulinemia, vasculitis, fatigue, or muscle symptoms.

\medterm{HDL} High-density lipoprotein, a class of particles that transports cholesterol and other lipids in the blood. HDL
concentration is one component of cardiovascular risk assessment; it should not be interpreted
as a protective or harmful measure by itself.

\medterm{HDL Cholesterol} Cholesterol carried in high-density lipoprotein particles; HDL participates in transporting cholesterol from tissues toward the liver, but its concentration alone does not determine cardiovascular risk.

\medterm{HDL Test} A blood test measuring high-density lipoprotein cholesterol, which is used with the rest of the lipid profile to assess cardiovascular risk.

\medterm{HDR Syndrome} A genetic disorder characterized by hypoparathyroidism, sensorineural deafness, and renal dysplasia or other kidney abnormalities, usually caused by a pathogenic variant in \textit{GATA3}.

\medterm{Head} The body region containing the brain, eyes, ears, nose, mouth, skull, and face, connected to the trunk by the neck and supplied by major cranial nerves and vessels.

\medterm{Head And Face Reconstruction} Reconstructive surgery to restore the form or function of the skull, face, or associated soft tissues after trauma, cancer treatment, congenital difference, or other disease.

\medterm{Head And Neck Cancer} Cancers arising in the oral cavity, pharynx, larynx, sinonasal tract, and salivary
glands, most commonly squamous cell carcinoma. Risk factors include tobacco, alcohol,
and HPV (oropharyngeal). Treatment involves surgery, radiation, and chemotherapy, often
combined to preserve function.

\medterm{Head And Neck Oncologic Surgery} Surgical treatment of cancers of the oral cavity, throat, larynx, nose, sinuses, salivary glands, or neck lymph nodes. Procedures may remove the tumor, treat involved lymph nodes, and reconstruct affected tissues while preserving breathing, swallowing, speech, and appearance when possible.

\synonyms
Head and Neck Surgical Oncology; Maxillofacial Oncologic Surgery

\medterm{Head Circumference} The measurement around the largest part of a baby’s or child’s head, taken above the eyebrows and ears. It is plotted over time to assess brain growth; unusually small, large, or rapidly changing measurements may need evaluation.

\medterm{Head CT Scan} Computed tomography of the head using x-rays and computer reconstruction to evaluate the brain, skull, bleeding, stroke, tumors, trauma, and other abnormalities.

\medterm{Head Injury} Damage to the scalp, skull, or brain caused by a blow, fall, collision, or penetrating object. It may cause a wound, skull fracture, concussion, bleeding, or brain damage. Worsening headache, repeated vomiting, seizure, confusion, weakness, or loss of consciousness requires urgent assessment.

\synonyms
Traumatic Brain Injury (TBI); Craniocerebral Trauma

\medterm{Head Injury First Aid} Immediate emergency first-aid protocols following head trauma, including stabilizing the cervical spine, maintaining an open airway, controlling external cranial bleeding with gentle pressure, and continuously monitoring level of consciousness and vital signs while urgently seeking medical evaluation for warning signs such as loss of consciousness, repeated vomiting, or focal neurologic deficits.

\synonyms
Cranial Trauma First Aid; Head Trauma Immediate Care; Concussion First Aid

\medterm{Head Lice} Infestation of the scalp by the parasitic insect \textit{Pediculus humanus capitis}. It causes itching and visible nits or lice and spreads mainly through close head-to-head contact; treatment requires an effective product and careful follow-up.

\medterm{Head Movement} Motion of the head produced mainly by the cervical spine and neck muscles, including flexion, extension, rotation, and side bending; pain or restricted movement may indicate injury or disease.

\medterm{Head MRI} Magnetic resonance imaging of the head using a magnetic field and radio waves to produce detailed images of the brain, cranial nerves, blood vessels, and surrounding structures.

\medterm{Head Trauma} An injury to the scalp, skull, brain, or other tissues of the head. Concussion, bleeding, skull fracture, and brain injury can result; worsening headache, repeated vomiting, seizure, confusion, weakness, or loss of consciousness requires emergency care.

\medterm{Headache} Pain or discomfort in the head, scalp, or upper neck. Primary headaches include migraine and tension headache; secondary headaches result from another problem such as infection, injury, bleeding, very high blood pressure, or a brain disorder.

\medterm{Headache and Migraine} Alternate name; see \textit{Headache}.

\medterm{Headache Red Flags} Critical clinical warning signs in a patient presenting with headache that indicate potentially life-threatening intracranial pathology (such as subarachnoid hemorrhage, meningitis, intracranial tumor, or temporal arteritis). Warning signs include sudden thunderclap onset, fever with meningismus, new focal neurological deficits, onset after age 50, papilledema, or postural worsening.

\synonyms
Headache Warning Signs; Red Flag Headaches; Danger Signs in Headache

\medterm{Head-Lice} Head lice are small insects that live on the scalp and cause itching; they spread mainly through close head-to-head contact.

\medterm{Healing} The body’s process of repairing damaged tissue. It usually involves stopping bleeding, reducing inflammation, forming new tissue, and remodeling a scar; healing may be delayed by infection, poor blood flow, diabetes, or inadequate nutrition.

\medterm{Health Anxiety Disorder} Excessive worry about having or developing a serious illness, often despite reassurance or limited medical findings. It can cause repeated checking or health-care visits and is treated with supportive care and evidence-based psychological therapy.

\medterm{Health Behavior} An action or pattern of action that affects health, such as eating, physical activity, tobacco use, alcohol use, sexual practices, or taking medicines.

\medterm{Health Behavior Change} Modification of behaviors that affect health, such as tobacco use, diet, physical
activity, and adherence to treatment.

\medterm{Health Behavior Motivation} Factors that influence initiation and maintenance of health-related behaviors, including
adherence, physical activity, and smoking cessation.

\medterm{Health Care} The prevention, diagnosis, treatment, and ongoing support provided to maintain or improve health, delivered by professionals, institutions, public-health systems, and informal caregivers.

\medterm{Health Care Agents} People legally chosen to make healthcare decisions for someone who cannot make or communicate them. Their authority and responsibilities depend on local law and the document that appoints them.

\medterm{Health Care Facility} A place where health services are delivered, such as a clinic, hospital, laboratory, nursing facility, or rehabilitation center; its services and level of care determine appropriate referrals.

\medterm{Health Care Personnel} Professionals and support workers involved in patient care, including clinicians, nurses, therapists, technicians, aides, and public-health staff.

\medterm{Health Care Proxy} A person designated to make healthcare decisions for someone who cannot make or communicate those decisions. The appointment may be recorded in an advance directive or other legal document.

\medterm{Health Checkup} A preventive visit that reviews symptoms, medical history, medicines, lifestyle, vaccinations, examination findings, and age- or risk-appropriate screening. The content should be individualized rather than based on unnecessary tests.

\medterm{Health Disparities} Differences in health outcomes or access to care among groups, often associated with social,
economic, environmental, geographic, or structural factors.

\medterm{Health Disparity} A difference in health outcomes or access to healthcare between population groups. It may be linked to income, housing, education, discrimination, environment, language, disability, geography, or unequal treatment.

\medterm{Health Education} The process of giving people understandable information and skills to prevent illness, recognize symptoms, use treatment safely, and make informed health decisions.

\medterm{Health Equity} The fair opportunity for all people to attain their highest possible level of health, without avoidable differences caused by social or economic disadvantage.

\medterm{Health For All} A public-health goal emphasizing that all people should have access to needed health services and conditions that support health, without financial hardship.

\medterm{Health Inequity} A systematic, avoidable, and unjust difference in health or access to health resources
between social groups. It differs from health variation that is not preventable or not
linked to unfair social conditions.

\medterm{Health Information Technology} The use of computers and digital systems to collect, store, exchange, protect, and use health information. Examples include electronic records, telemedicine, patient portals, decision-support tools, and remote monitoring; safe systems protect privacy and support accurate care.

\medterm{Health Literacy} The ability to find, understand, judge, and use health information and services to make appropriate decisions. It includes communicating with clinicians, following instructions, reading medicine labels, and recognizing when to seek help; clear communication by healthcare staff also improves health literacy.

\medterm{Health Maintenance Organization} A health-insurance arrangement that usually uses a network of clinicians and requires a primary-care clinician to coordinate specialist care. Coverage rules, referrals, costs, and services vary by plan.

\medterm{Health Promotion} Actions that help people improve or protect their health and prevent illness. They may include vaccination, healthy eating, physical activity, sleep, screening, safer environments, social connection, and avoiding tobacco. Advice should reflect a person’s age, abilities, culture, risks, and preferences.

\medterm{Health Risk} The likelihood that a behavior, exposure, condition, or other factor will cause harm to health, considered together with the severity and timing of possible harm.

\medterm{Health Risks Of Alcohol Use} Alcohol use can increase the risk of liver disease, several cancers, high blood pressure, injuries, dependence, and problems affecting mental health and relationships.

\medterm{Health Surveillance} The ongoing collection and analysis of health data to identify patterns, monitor disease, and guide public-health action.

\medterm{Health Survey} A structured set of questions used to collect information about health, symptoms, behaviors, healthcare use, or population needs.

\medterm{Health Visitor} A community health professional, especially in the United Kingdom, who supports families with infants and young children. They provide health advice, developmental monitoring, prevention, and referral to other services.

\medterm{Health Workforce} The people and organizations that provide, support, manage, or study health services, including clinicians, public-health workers, and support staff.

\medterm{Healthcare Associated Infection} An infection acquired while receiving healthcare or as a result of a medical procedure. It may occur in hospitals, clinics, nursing facilities, or at home and can involve a wound, urinary tract, lungs, bloodstream, or other site.

\medterm{Healthcare Provider} A healthcare professional or organization that provides medical care, such as a
physician, nurse, therapist, or clinic.

\medterm{Healthcare Proxy} A person designated to make healthcare decisions for someone who cannot make or communicate those decisions. The appointment may
be recorded in an advance directive or other legal document.

\medterm{Health-Related QoL} A person’s physical, emotional, social, and daily functioning as affected by health and illness. It reflects the patient’s experience and is assessed separately from laboratory results or survival.

\medterm{Healthy Aging} Developing and maintaining the physical and mental abilities that support well-being and independence in later life. It is influenced by health conditions, activity, nutrition, relationships, environment, and access to care.

\medterm{Healthy Years} Years lived without substantial disease, disability, or loss of function. This population measure helps describe quality of life and health needs, but it does not predict an individual’s exact lifespan.

\medterm{Healthy Years Lost} The years lost because of early death or time lived with illness or disability. This population measure estimates the overall effect of disease; it does not predict exactly how long one person will live.

\medterm{Hearing} The ability to detect and understand sound. Sound vibrations are converted into nerve signals in the inner ear and interpreted by the brain. Hearing loss may result from the outer or middle ear, inner ear, hearing nerve, or brain pathways.

\medterm{Hearing Aid} A small electronic device worn in or behind the ear that makes sounds louder or clearer for someone with hearing loss. It contains a microphone, processor, and speaker and is fitted and adjusted according to hearing-test results.

\medterm{Hearing Aids And Devices} Electronic devices worn in or behind the ear that amplify sound for individuals with
hearing loss. Modern hearing aids include digital signal processing, directional
microphones, and connectivity to smartphones and other audio sources.

\medterm{Hearing Disorder} A condition that reduces or distorts hearing through disease of the outer or middle ear, cochlea, auditory nerve, or central pathways; audiometry and examination help distinguish conductive from sensorineural loss.

\medterm{Hearing Impairment} Partial or complete reduction in the ability to detect or understand sound. It may be
conductive, sensorineural, or mixed and can result from congenital disorders, infection,
noise, ototoxic injury, trauma, ageing, or disease of the ear or auditory pathways.

\medterm{Hearing Loss} Reduced ability to hear. Conductive loss results from a problem in the outer or middle ear; sensorineural loss results from damage to the inner ear or hearing nerve; mixed loss has both types. Causes include aging, noise, infection, injury, and some medicines.

\synonyms
Hearing Impairment; Hypoacusis

\medterm{Hearing Loss And Music} Difficulty hearing or distinguishing musical sounds because of reduced hearing, damage to inner-ear cells, or problems processing sound. Music may seem quieter, distorted, or lacking in certain pitches; hearing protection and individualized treatment may help.

\medterm{Hearing Loss: Sudden Deafness} A rapid sensorineural hearing loss, usually developing over hours to a few days, that may affect one or both ears. It is an otologic emergency; urgent hearing testing and evaluation are needed because early treatment may improve recovery.

\medterm{Hearing Screening} A quick test used to identify possible hearing loss before symptoms are obvious. Newborns are commonly screened before leaving hospital; an abnormal result requires repeat testing and a full hearing assessment rather than confirming permanent loss by itself.

\medterm{Hearing Test} An audiologic assessment that measures hearing thresholds and speech perception across frequencies to identify conductive, sensorineural, or mixed hearing loss.

\medterm{Heart} A four-chambered muscular pump (two atria, two ventricles) located in the mediastinum,
responsible for circulating blood throughout the body.

\medterm{Heart And Vascular Services} Heart and vascular services include prevention, diagnosis, and treatment of diseases affecting the heart and blood vessels.

\medterm{Heart Aneurysm} A localized bulging or outpouching of a weakened heart wall, often after myocardial infarction or from cardiomyopathy, that can impair contraction or form clots and arrhythmias.

\medterm{Heart Arrhythmia} An abnormal heart rate or rhythm caused by altered electrical activity in the heart. At rest, the adult heart usually beats regularly at about 60--100 times per minute; an arrhythmia may make it beat too fast, too slowly, or irregularly.

\synonyms
Cardiac Dysrhythmias; Irregular Heart Rhythm

\medterm{Heart Attack} A heart attack occurs when blood flow carrying oxygen to part of the heart muscle suddenly becomes blocked, usually by a blood clot in a narrowed heart artery. The lack of oxygen can injure or destroy heart muscle and may cause chest pressure, shortness of breath, sweating, nausea, or other symptoms.

\medterm{Heart Attack First Aid} Heart attack first aid means calling emergency services immediately, helping the person rest, and following dispatcher instructions; aspirin may be appropriate for some adults who can safely take it.

\medterm{Heart Attack Prevention} Heart-attack prevention targets blood pressure, atherogenic cholesterol, diabetes, tobacco exposure, diet, physical activity, weight, sleep, and adherence to prescribed medicines. Risk-based preventive treatment is more effective than relying on one supplement or symptom.

\medterm{Heart Block} A delay or interruption in the electrical signal traveling from the heart’s upper chambers to its lower chambers. It may cause a slow or irregular heartbeat, dizziness, fainting, or no symptoms.

\medterm{Heart Bypass Surgery} Coronary artery bypass grafting creates new routes around blocked coronary arteries using graft vessels to improve myocardial blood flow and relieve ischemic symptoms or reduce selected risks.

\medterm{Heart Catheterization} A procedure in which a catheter is passed through a blood vessel into the heart to measure pressures, assess coronary arteries, inject contrast, obtain samples, or perform treatment.

\medterm{Heart Contraction} The squeezing of heart muscle that pushes blood from the lower chambers to the lungs and the rest of the body. Electrical signals coordinate the heartbeat, and abnormal contraction can reduce blood flow or cause rhythm problems.

\medterm{Heart CT Scan} Computed tomography of the heart and nearby structures, sometimes with ECG synchronization or contrast, used to evaluate coronary arteries, cardiac anatomy, calcium, or selected masses.

\medterm{Heart Defect} An abnormality of the heart or major blood vessels. It may be present from birth or develop later and can affect blood flow, heart rhythm, valve function, or pumping ability.

\medterm{Heart Disease} A broad term for conditions affecting the heart or blood vessels, including coronary artery disease, heart failure, abnormal rhythms, valve disease, and congenital defects. Symptoms and treatment depend on the part affected and may include chest pain, breathlessness, palpitations, or swelling.

\medterm{Heart Disease And Depression} Depression can occur with heart disease and may worsen quality of life, treatment adherence, and recovery; both conditions should be assessed and treated.

\medterm{Heart Disease And Intimacy} Heart disease or its treatment may affect sexual activity, desire, and confidence; individualized medical advice can help address symptoms and safety.

\medterm{Heart Disease And Women} Heart disease in women may involve different risk factors, symptoms, and timing than in men; chest pressure, shortness of breath, nausea, fatigue, or jaw or back pain can be important symptoms.

\medterm{Heart Failure} A condition in which the heart cannot pump or fill with enough blood to meet the body's needs. It does not mean that the heart has stopped. In systolic heart failure, the heart does not squeeze strongly enough; in diastolic heart failure, the heart is stiff and does not relax or fill properly. Symptoms may include shortness of breath, tiredness, reduced ability to exercise, and swelling caused by fluid buildup.

\medterm{Heart Failure In Children} A condition in which a child's heart cannot pump or fill well enough to meet the body's needs. It may cause fast breathing, sweating or tiring during feeds, poor growth, swelling, or reduced activity.

\medterm{Heart Failure Stages} A framework describing progression from risk without heart disease, to heart changes without symptoms, to heart failure with symptoms, and finally advanced disease with persistent symptoms. Stages help guide prevention and treatment and are different from scales that describe how limited a person feels during activity.

\medterm{Heart Failure With Mildly Reduced Ejection Fraction} Heart failure in which the left ventricle pumps out a somewhat smaller-than-usual proportion of its blood, commonly about 41--49 percent. Symptoms may include breathlessness, tiredness, and fluid-related swelling.

\medterm{Heart Failure With Preserved Ejection Fraction} Heart failure in which the left ventricle still squeezes out a usual proportion of its blood but is too stiff to fill normally. Fluid then backs up, causing breathlessness, tiredness, and swelling.

\medterm{Heart Failure With Reduced Ejection Fraction} Heart failure in which the left ventricle cannot squeeze strongly enough to supply the body. It may cause breathlessness, tiredness, and swelling and can result from blocked arteries, heart-muscle disease, high blood pressure, or valve disease.

\medterm{Heart Injury} Damage to heart muscle or surrounding cardiac structures caused by trauma, ischemia, inflammation, toxins, or procedures; symptoms and tests depend on the injured tissue and severity.

\medterm{Heart MRI} Magnetic resonance imaging of the heart that evaluates cardiac structure, function, tissue characteristics, blood flow, and selected congenital, ischemic, inflammatory, or infiltrative disorders.

\medterm{Heart Murmur} An extra or unusual sound heard when blood flows through the heart or nearby blood vessels. Some murmurs are harmless, while others result from a valve problem, heart defect, or altered blood flow.

\medterm{Heart Murmurs} Heart murmurs are extra or unusual sounds caused by turbulent blood flow through the heart or nearby vessels; some are harmless and others reflect heart disease.

\medterm{Heart Muscle} The myocardium, the specialized cardiac muscle that contracts rhythmically to pump blood through the heart and circulation.

\medterm{Heart Neoplasm} An abnormal growth in the heart or its surrounding tissues. It may be benign or cancerous and can interfere with blood flow, heart rhythm, or pumping; evaluation uses imaging and sometimes biopsy.

\medterm{Heart Pacemaker} A heart pacemaker is a device that sends electrical impulses to maintain an appropriate heartbeat when the heart's natural rhythm is too slow or unreliable.

\medterm{Heart Palpitations} An awareness of the heartbeat as pounding, racing, fluttering, or irregular.
Palpitations may occur with normal rhythm or arrhythmia and can be triggered by stress,
exertion, fever, stimulants, endocrine disease, or structural heart disease.

\synonyms
Cardiac Palpitations; Ectopic Beats

\medterm{Heart PET Scan} Positron-emission tomography of the heart using a radiotracer to assess myocardial perfusion, metabolism, viability, inflammation, or selected cardiac masses.

\medterm{Heart Rate} The number of heartbeats per minute. Resting rates vary with age, fitness, activity, illness, medicines, and emotion; a persistently very fast, slow, or irregular rate may need assessment.

\medterm{Heart Rhythm} The rate and pattern of cardiac electrical activation and contraction, normally initiated by the sinoatrial node; abnormal rhythms include bradycardia, tachycardia, atrial fibrillation, and ventricular arrhythmias.

\medterm{Heart Rhythm Disorders} Arrhythmias, in which the heart beats too fast, too slowly, or irregularly because of abnormal electrical impulse formation or conduction.

\synonyms
Cardiac Arrhythmias; Dysrhythmias

\medterm{Heart Septum} The muscular or membranous wall separating the right and left sides of the heart, including the atrial and ventricular septa; defects can permit abnormal blood flow between chambers.

\medterm{Heart Sounds} The sounds made as the heart valves open and close during each heartbeat. The usual two sounds are called “lub-dub”; extra or unusual sounds may suggest a valve problem, heart failure, or another condition.

\medterm{Heart Surgery} Surgical procedures performed on the heart or great vessels to treat structural or
functional cardiac conditions. Common operations include coronary artery bypass grafting
(CABG), valve repair or replacement, and repair of congenital heart defects.

\medterm{Heart Tamponade} A life-threatening condition in which fluid or blood builds up around the heart and prevents it from filling normally. It can cause breathlessness, faintness, low blood pressure, and shock and requires emergency treatment.

\medterm{Heart Transplant} Replacement of a failing heart with a donor heart, followed by lifelong immunosuppression and surveillance for rejection, infection, coronary allograft disease, kidney injury, and medication complications.

\medterm{Heart Transplantation} Surgical replacement of a failing heart with a donor heart, usually for selected patients with advanced heart failure who meet specialist criteria.

\medterm{Heart Valve} A flap-like structure that directs one-way blood flow through the heart. The four valves are the tricuspid, pulmonary, mitral, and aortic
valves.

\medterm{Heart Valve Disease} Structural or functional abnormality of one or more cardiac valves causing stenosis, regurgitation, or both. Common causes include degeneration, congenital malformation, rheumatic disease, and infective endocarditis.

\medterm{Heart Valve Disorder} Alternate term; see \textit{Heart valve disease}.

\medterm{Heart Valve Prosthesis} An artificial valve used to replace or repair a diseased heart valve. Mechanical valves are durable but usually require long-term blood-thinning medicine. Tissue valves may avoid lifelong blood thinners but can wear out and sometimes need another procedure.

\medterm{Heart Valve Surgery} Repair or replacement of a diseased cardiac valve to correct severe narrowing or leakage and improve blood flow; choice of repair, mechanical valve, or tissue valve depends on anatomy and patient factors.

\medterm{Heart Ventricle} One of the two lower heart chambers that pumps blood to the lungs or systemic circulation; the right ventricle pumps to the pulmonary artery and the left to the aorta.

\medterm{Heart Ventricles} The two lower chambers of the heart. The right ventricle pumps blood to the lungs, and the left ventricle pumps blood to the body.

\medterm{Heartburn} A burning feeling behind the breastbone caused by stomach contents flowing back into the esophagus. It may be accompanied by a sour taste or regurgitation and is not usually caused by the heart. Frequent symptoms may indicate reflux disease.

\medterm{Heartland Virus} A tick-borne virus associated with fever, fatigue, headache, muscle aches, diarrhea, and low white-cell or platelet counts. It can resemble ehrlichiosis and may be severe, especially in older or medically vulnerable people.

\medterm{Heart-Lung Machine} A heart-lung machine temporarily circulates and oxygenates blood during some heart or major blood-vessel operations while the heart is stopped.

\medterm{Heart-Lung Machine} A device that temporarily takes over the pumping and gas-exchange functions of the heart and lungs during selected cardiac operations.

\synonyms
Cardiopulmonary Bypass Pump; Extracorporeal Circuit

\medterm{Heat} Thermal energy that raises tissue or body temperature; excessive heat exposure can cause heat cramps, heat exhaustion, heatstroke, or thermal injury.

\medterm{Heat Cramps} Painful, involuntary muscle spasms occurring during or after strenuous activity in heat, often related to fluid and
electrolyte loss. Resting in a cool place and replacing fluids and electrolytes may help. Confusion, fainting, or
ongoing symptoms may signal heat illness requiring urgent medical attention.

\medterm{Heat Emergencies} Heat emergencies range from heat cramps and exhaustion to heat stroke, in which dangerous core-temperature elevation causes central nervous system dysfunction and organ injury.

\medterm{Heat Exhaustion} A heat-related illness caused by losing too much water and salt through sweating. It may cause heavy sweating, weakness, dizziness, headache, nausea, or fainting. Confusion, seizures, or very hot skin suggest heat stroke and require emergency care.

\medterm{Heat Intolerance} An abnormal inability to tolerate warm environments or increased body heat, often
causing excessive sweating, flushing, fatigue, or dizziness. Common causes include
hyperthyroidism, menopause, fever, medications, autonomic disorders, and dehydration.

\medterm{Heat Pack} A device that applies warmth to the body surface to relieve some muscle pain or stiffness. Excessive heat or prolonged contact can cause burns.

\medterm{Heat Prostration} An older term for severe heat exhaustion, marked by weakness, fatigue, dizziness, headache, nausea, and sweating after heat exposure; confusion or loss of consciousness suggests heat stroke.

\medterm{Heat Rash} An itchy or prickly rash caused by blocked sweat ducts during hot, humid conditions. Small red bumps or clear vesicles
often appear in skin folds or areas covered by clothing. Cooling, keeping the skin dry, and avoiding occlusive products
usually help; infection or persistent symptoms warrants assessment.

\medterm{Heat Shock Protein} A stress-response protein that helps other proteins fold, avoid aggregation, and recover after heat, oxidative, infectious, or other cellular stress; some also influence immunity and cancer biology.

\medterm{Heat Stress} The strain placed on the body when it cannot adequately lose heat, often during hot or humid conditions, exertion, or inadequate hydration. It can progress from heat cramps or exhaustion to life-threatening heat stroke.

\medterm{Heat Stroke} A life-threatening heat-related illness with markedly elevated body temperature and altered mental status. It may cause confusion, seizures, loss of consciousness, and damage to multiple organs.

\medterm{Heavy Menstrual Bleeding} Menstrual bleeding that is unusually heavy, prolonged, or disruptive to daily life. Causes include fibroids, hormone changes, bleeding disorders, medicines, pregnancy-related problems, or cancer and may cause anemia.

\synonyms
Menorrhagia; Hypermenorrhea

\medterm{Heavy Metal Poisoning} Harm from excessive exposure to metals such as lead, mercury, arsenic, cadmium, or thallium. Effects may include abdominal pain, anemia, nerve injury, kidney damage, or confusion; treatment starts with removing exposure and may include specialist antidote therapy.

\medterm{Hebefrenia} Alternate historical term; see \textit{Schizophrenia}.

\medterm{Heberden Node} A firm bony bump at the finger joint nearest the fingertip, usually caused by osteoarthritis. It may be painful or tender while developing and can cause stiffness or reduced finger movement.

\medterm{Heberden Nodes} Firm bony enlargements at the finger joints nearest the fingertips, commonly caused by osteoarthritis. They may be painful or tender while forming, then remain as permanent bumps that limit movement in some people.

\synonyms
Heberden's Arthritic Nodes; DIP Joint Osteophytes

\medterm{Heberden's Node} A hard, bony enlargement of the distal interphalangeal joint caused by osteoarthritis, usually affecting the fingers and sometimes producing pain, stiffness, or reduced dexterity.

\medterm{Heberden's Nodes} Firm bony enlargements of the distal interphalangeal finger joints, usually caused by osteoarthritis.

\medterm{Hecht's Pneumonia} A rare, severe form of measles-related pneumonia, usually affecting people with weakened immunity. It can cause rapidly worsening breathing problems and requires urgent hospital treatment.

\medterm{Hedgehog Proteins} Signaling proteins that help guide the growth and organization of tissues before birth and during repair. Abnormal activity in this pathway can contribute to some birth differences and cancers.

\medterm{Heel} The heel is the back and underside of the foot, formed mainly by the calcaneus and a cushioning fat pad.

\medterm{Heel Pain} Pain beneath or around the heel, commonly caused by plantar fasciitis, Achilles
tendinopathy, bursitis, stress fracture, nerve compression, or inflammatory arthritis.

\synonyms
Plantar Fasciitis Discomfort; Calcaneal Spur Pain

\medterm{Heel Spur} A bony growth arising from the heel bone, often associated with plantar-fascia strain. The spur itself
may cause no symptoms; heel pain is more often related to surrounding soft-tissue inflammation.

\medterm{Heel-to-Shin Test} A neurological test of lower limb coordination in which the patient places the heel of one foot on the opposite knee and slides it smoothly down the crest of the tibia to the ankle. Inability to maintain contact or an irregular, jerky trajectory indicates cerebellar ataxia or sensory ataxia.

\medterm{Hegar's Dilators} Hegar's dilators are graduated metal instruments used to gently widen the cervix or other body openings during selected procedures.

\medterm{Hegar's Sign} Softening of the lower part of the uterus, once used as an early sign of pregnancy. It is not specific and is not sufficient to confirm pregnancy without testing.

\medterm{Height} The vertical distance from the base of a person to the top of the head, measured for growth assessment, nutritional evaluation, drug dosing, and calculation of body-mass index.

\medterm{Heimlich Maneuver} A common name for abdominal thrusts used to help a responsive person with severe choking. First-aid steps differ for adults, children, infants, and pregnant people; emergency services should be called for serious airway obstruction.

\medterm{Heimlich Maneuver On Self} A self-administered abdominal-thrust technique sometimes used for severe choking when a person cannot breathe, speak, or cough effectively; emergency services should be activated immediately.

\medterm{Heimlich Manoeuvre} An emergency procedure for relieving upper airway obstruction caused by a foreign body
(choking). The maneuver is performed on a conscious, choking adult who cannot speak,
cough effectively, or breathe.

\medterm{Heine-Medin Disease} An older name for polio, a viral infection that can inflame the spinal cord and brain. It may cause fever and weakness and sometimes permanent paralysis; vaccination prevents most cases.

\medterm{Heinz Bodies} Clumps of damaged hemoglobin inside red blood cells, usually formed after oxidative injury. They may occur in glucose-6-phosphate dehydrogenase deficiency or after exposure to certain medicines, chemicals, or infections.

\medterm{Hejama} An Arabic term for cupping, a traditional practice in which suction is applied to the skin; wet cupping also involves superficial skin incisions and carries bleeding and infection risks.

\medterm{Helcococcus} A genus of gram-positive cocci that can rarely cause opportunistic skin, wound, bone, or bloodstream infections.

\medterm{Helcococcus Kunzii} A Gram-positive round bacterium usually found on skin that rarely causes opportunistic infection. Reported infections may involve wounds, blood, or implanted devices in people with serious illness or weakened defenses.

\medterm{Helcococcus Pyogenes} A rarely encountered Gram-positive round bacterium that may cause opportunistic human infection. It has been found in wounds or blood, particularly in people with serious illness, chronic wounds, or weakened immunity.

\medterm{Helicase} An enzyme that uses cellular energy to separate the paired strands of DNA or RNA. This allows genetic material to be copied, repaired, or read.

\medterm{Helicobacter} A genus of spiral-shaped bacteria that colonize the stomach or intestine; some species cause ulcers, gastritis, or other disease.

\medterm{Helicobacter Cinaedi} A Helicobacter species that can cause bloodstream, skin, bone, joint, or other invasive infections, especially in people with weakened immunity.

\medterm{Helicobacter Fennelliae} An intestinal Helicobacter species that rarely causes bloodstream or other opportunistic infection. Disease is more likely in people with weakened defenses, and diagnosis requires specialized laboratory identification.

\medterm{Helicobacter Pylori} A gram-negative, spiral-shaped bacterium adapted to colonize the stomach. It can cause chronic
gastritis, peptic ulcer disease, and increase the risk of gastric cancer and gastric MALT
lymphoma.

\medterm{Helicobacter Pylori Infection} Infection of the stomach by H. pylori, a spiral-shaped bacterium associated with chronic
gastritis, peptic ulcer disease, gastric adenocarcinoma, and gastric MALT lymphoma. It
is diagnosed with noninvasive or endoscopic testing and treated with combination
eradication therapy.

\medterm{Helicobacteraceae} A family of spiral-shaped bacteria that includes Helicobacter and related genera, many of which inhabit the digestive tract.

\medterm{Heliotherapy} Treatment using controlled exposure to sunlight or artificial ultraviolet light. It has been used for selected skin disorders, but excessive exposure can cause burns, premature aging, eye injury, and skin cancer.

\medterm{Helium} A chemically inert gas used in respiratory mixtures, cryogenics, imaging, and laboratory applications; inhalation can displace oxygen and cause hypoxia despite its lack of direct toxicity.

\medterm{Helix} The curved outer rim of the external ear. Its shape is formed by cartilage covered with skin; swelling, tenderness, or a skin change on the helix may follow injury, infection, sun damage, or other skin disease.

\medterm{Heller Myotomy} A surgical procedure that cuts the muscles of the lower esophageal sphincter to relieve obstruction in achalasia, usually performed laparoscopically and often combined with an antireflux procedure.

\medterm{Heller's Operation} A surgical procedure that cuts the lower esophageal muscle to relieve achalasia, usually with an anti-reflux procedure.

\medterm{HELLP Syndrome} A dangerous pregnancy complication involving destruction of red blood cells, liver injury, and a low platelet count. It may cause upper-abdominal pain, nausea, headache, or high blood pressure and can occur with preeclampsia. Urgent obstetric assessment is required.

\medterm{Helminth} A parasitic worm, including roundworms, tapeworms, and flukes, that may live in or on a host.

\medterm{Helminthiasis} Infection or infestation caused by parasitic worms, including roundworms, tapeworms, and flukes. Symptoms vary by species and may include abdominal problems, anemia, cough, rash, malnutrition, or organ damage.

\synonyms
Worm Infestation; Helminthic Parasitosis

\medterm{Helminths} Parasitic worms that cause disease in humans. The major groups are nematodes
(roundworms), cestodes (tapeworms), and trematodes (flukes). Helminths have complex life
cycles often involving intermediate hosts. Infection occurs through ingestion of eggs or
larvae, penetration of skin, or insect vectors.

\medterm{Help Prevent Hospital Errors} Actions that reduce hospital errors include confirming identity, checking medicines, communicating clearly, asking questions, and reporting changes or concerns promptly.

\medterm{Helper T Cell} A type of white blood cell that coordinates immune responses by sending chemical signals that activate and guide other immune cells, including antibody-producing cells and cells that destroy infected cells.

\medterm{Helper T Cells} A type of white blood cell that helps organize the immune response. These cells release chemical messages that activate and guide other immune cells, including cells that make antibodies or destroy infected cells.

\synonyms
CD4+ T Lymphocytes; T Helper Cells

\medterm{Helper Virus} A virus that supplies functions needed for another virus or defective viral genome to replicate.

\medterm{HELPERR Mnemonic} A memory aid for managing shoulder dystocia during birth: call for help, consider an episiotomy, flex the birthing person’s legs, apply suprapubic pressure, perform internal maneuvers, remove the baby’s posterior arm, and roll to an all-fours position. It requires trained obstetric staff.

\medterm{Helping A Loved One With A Drinking Problem} Support for someone with alcohol use disorder includes calm nonjudgmental communication, boundaries, encouragement toward medical treatment, attention to withdrawal risk, and emergency action for overdose or dangerous withdrawal.

\medterm{Hemagglutination Test} A test in which antibodies or other agents cause red blood cells to clump; the reaction can help detect antibodies, antigens, or selected infectious agents.

\medterm{Hemagglutinin} A surface glycoprotein, especially on influenza viruses, that binds sialic-acid receptors and helps the virus attach to and enter host cells; it is a major target of protective antibodies.

\medterm{Hemangioblastoma} A usually slow-growing tumor made of blood-vessel-forming cells, most often found in the cerebellum or spinal cord. It may cause headache, balance problems, or weakness and can be associated with von Hippel–Lindau disease.

\medterm{Hemangioendothelioma} A rare tumor arising from cells that line blood vessels. Its behavior ranges from slow-growing to cancerous, depending on the subtype and location, so diagnosis requires tissue examination and specialist assessment.

\medterm{Hemangioma} A benign vascular tumor formed by the abnormal proliferation of blood vessel endothelial cells. Hemangiomas are broadly classified based on vessel caliber and clinical presentation into capillary hemangiomas (such as infantile hemangiomas that proliferate in infancy and spontaneously involute), cavernous hemangiomas (composed of larger, mature blood-filled vascular spaces that do not regress), and lobular capillary hemangiomas (pyogenic granulomas). While many are asymptomatic or undergo spontaneous involution, problematic lesions causing functional impairment, ulceration, or visual compromise (e.g., periocular capillary hemangiomas) are treated medically with beta-blockers (propranolol), laser therapy, or surgical excision.

\synonyms
Vascular Hemangioma; Benign Vascular Tumor

\medterm{Hemangiopericytoma} A rare tumor formerly defined by branching “staghorn” vessels and perivascular cells; many lesions classified this way are now recognized as solitary fibrous tumors.

\medterm{Hemangiosarcoma} A malignant tumor of vascular endothelial cells, an aggressive sarcoma that may arise in
the skin, liver, or bone.

\medterm{Hemarthrosis} Bleeding into a joint, most often the knee, causing rapid swelling, warmth, pain, and limited movement. Injury and bleeding disorders such as hemophilia are common causes; repeated episodes can damage cartilage and lead to long-term joint disease.

\medterm{Hematemesis} Vomiting blood or material that looks like coffee grounds, usually from bleeding in the esophagus, stomach, or first part of the small intestine. It can be life-threatening and requires emergency medical assessment.

\medterm{Hematidrosis} A very rare condition in which blood-tinged fluid appears through apparently intact skin or other surfaces, sometimes during extreme stress. The diagnosis requires careful examination to exclude injury, skin disease, bleeding disorders, or contamination.

\medterm{Hematocele} A collection of blood around a testicle, usually after scrotal injury or surgery and sometimes from a tumor. It can cause a heavy, firm, painful swelling; ultrasound is needed to check the testicle and blood flow, and large collections may need surgery.

\synonyms
Scrotal Hematocele; Tunica Vaginalis Blood Collection

\medterm{Hematochezia} Passing bright-red or maroon blood from the rectum. It often comes from the large intestine, rectum, or anus, but heavy bleeding higher in the digestive tract can look the same and requires urgent assessment.

\medterm{Hematocolpos} Accumulation of menstrual blood in the vagina, usually because an imperforate hymen or another obstructing anomaly prevents outflow; it may cause cyclic pelvic pain and a bulging vaginal membrane.

\medterm{Hematocrit} The proportion of blood volume occupied by packed red blood cells, expressed as a percentage. One of the most commonly ordered laboratory tests in clinical medicine.

\medterm{Hematogenous} Spread through the bloodstream, as when infection, tumor cells, or embolic material travels from one body site to another.

\medterm{Hematologic Agents} Medicines or biologic therapies that act on blood cells, coagulation, marrow, or immune pathways, including agents used for anemia, thrombosis, leukemia, and lymphoma.

\medterm{Hematologic Changes In Pregnancy} Pregnancy normally increases the fluid part of blood more than the red-cell mass, so blood counts may show mild dilutional anemia. White-cell counts also commonly rise, mainly because of more neutrophils. Results must be interpreted using pregnancy-specific ranges.

\medterm{Hematologic Malignancy} Alternate name; see \textit{Blood Cancer}.

\medterm{Hematologic Response} A measurable improvement or change in blood counts or marrow findings after treatment, such as recovery of neutrophils, platelets, or hemoglobin or reduction of malignant cells.

\medterm{Hematological Disease} A disorder of blood cells, bone marrow, lymphatic tissue, or coagulation, including anemia, leukemia, lymphoma, bleeding disorders, and clotting disorders.

\medterm{Hematologist} A physician who diagnoses and treats disorders of blood, bone marrow, lymph nodes, clotting, and hematologic cancers.

\medterm{Hematology} The branch of medicine concerned with the study, diagnosis, treatment, and prevention of
disorders of the blood and blood-forming organs (bone marrow, lymph nodes, spleen).

\medterm{Hematology-Oncology} The medical specialty concerned with blood disorders and cancers, including anemia, leukemia, lymphoma, myeloma, and clotting disorders.

\medterm{Hematoma} A collection of blood outside a blood vessel, usually after injury, surgery, or bleeding caused by illness or medicines. It may cause swelling and pain and can press on nearby nerves or organs.

\medterm{Hematometra} Blood trapped inside the uterus because the normal outlet is blocked or too narrow. It can cause absent periods, monthly pelvic pain, and an enlarged tender uterus; the blockage may be present from birth or develop after scarring, surgery, infection, or radiation.

\medterm{Hematometrocolpos} A buildup of menstrual blood in the uterus and vagina because the normal outlet is blocked, often by tissue present from birth. It may cause absent periods, monthly pelvic pain, and a lower-abdominal mass.

\medterm{Hematomyelia} Bleeding inside the spinal cord, usually after a serious injury or from an abnormal blood vessel or bleeding disorder. Sudden neck or back pain, weakness, numbness, or loss of bladder or bowel control is an emergency requiring immediate imaging and treatment.

\medterm{Hematophobia} An excessive or persistent fear of blood, blood-related sights, or medical procedures involving blood that causes distress or avoidance; some people have a vasovagal fainting response.

\medterm{Hematopoiesis} The process of making and maturing blood cells, mainly in the bone marrow. It produces red blood cells for oxygen transport, white blood cells for immune defense, and platelets for stopping bleeding.

\medterm{Hematopoietic Growth Factor} A cytokine or medicine that stimulates production, maturation, or survival of blood-forming cells in bone marrow, such as erythropoietin or granulocyte colony-stimulating factor.

\medterm{Hematopoietic Neoplasm} A cancer arising from blood-forming cells in the bone marrow or immune tissues. Leukemias, lymphomas, and plasma-cell cancers are examples; symptoms and treatment depend on the exact cell type and extent.

\medterm{Hematopoietic Stem Cell Transplantation} Infusion of blood-forming stem cells to replace damaged or diseased bone marrow. The cells may come from the patient or a donor. Treatment can cause serious infection, bleeding, infertility, or an immune reaction against donor and recipient tissues.

\medterm{Hematopoietic System} The bone marrow and related organs and cells that produce, mature, circulate, and remove blood cells, including red cells, white cells, and platelets.

\medterm{Hematoporphyrin Derivative} A porphyrin-based photosensitizer that accumulates in some tumors and can be activated by light in photodynamic treatment; it may also cause prolonged photosensitivity.

\medterm{Hematosalpinx} Accumulation of blood within a fallopian tube, often caused by ectopic pregnancy, endometriosis, pelvic surgery, or obstruction, and potentially producing pelvic pain or infertility.

\medterm{Hematosis} The exchange of gases in the lungs: oxygen moves from inhaled air into the blood, while carbon dioxide moves from the blood into the air to be breathed out.

\medterm{Hematospermia} Blood in the semen, often temporary and associated with infection, inflammation, trauma, or a recent urologic procedure. Persistent or recurrent bleeding, pain, fever, urinary symptoms, or bleeding in an older adult requires medical evaluation.

\medterm{Hematoxylin} A dye used in histology, usually with eosin, to stain cell nuclei blue-purple and help show tissue structure.

\medterm{Hematuria} Blood in the urine. It may be visible as pink, red, or brown urine or found only under a microscope. Causes include infection, stones, medicines, injury, prostate disease, and kidney or urinary-tract disorders.

\medterm{Heme} Heme is an iron-containing porphyrin group that binds oxygen in haemoglobin and myoglobin and forms the functional centre of several enzymes; abnormal heme metabolism contributes to porphyrias and jaundice.

\medterm{Hemeralopia} Hemeralopia is reduced vision in bright light, sometimes associated with retinal disorders or abnormal cone function.

\medterm{Hemianopia} Loss of vision in half of the area a person can see, affecting one or both eyes. It often results from a problem in the brain’s visual pathways, such as a stroke or tumor, and sudden onset is an emergency.

\medterm{Hemianopsia} Loss of vision in half of the visual field of one or both eyes, usually from a lesion of the optic tract, radiations, or occipital cortex; sudden onset may indicate stroke.

\medterm{Hemiballism} A movement disorder causing sudden, large-amplitude, flinging movements on one side of the body, usually from a lesion involving the subthalamic region.

\medterm{Hemiballismus} Hemiballismus is involuntary, rapid, flinging movement of one side of the body, usually caused by a lesion of the opposite subthalamic region.

\medterm{Hemicolectomy} Surgical removal of the right or left portion of the colon, sometimes with nearby lymph nodes and reconnection of the bowel.

\medterm{Hemidesmosome} A specialized junction that anchors basal epithelial cells to the underlying basement membrane by linking intracellular keratin filaments with extracellular matrix proteins.

\medterm{Hemidiaphragm} One of the two muscular halves of the diaphragm separating the chest from the abdomen and assisting breathing.

\medterm{Hemifacial Atrophy} Progressive wasting of skin, subcutaneous tissue, muscle, and sometimes bone on one side of the face, also called Parry--Romberg syndrome.

\medterm{Hemifacial Microsomia} A birth difference in which one side of the face is smaller or less developed than the other, often affecting the jaw, ear, cheek, and nearby soft tissues. Hearing, vision, and spinal differences may also occur.

\synonyms

\medterm{Hemifacial Spasm} Repeated involuntary twitching on one side of the face, usually beginning around the eye and sometimes spreading to the cheek or mouth. Irritation of the facial nerve is a common cause; persistent spasms need evaluation.

\medterm{Hemihypertrophy} Unequal overgrowth of one side or region of the body, producing asymmetry of limbs, trunk, or face; some childhood cases carry increased risk of embryonal tumors.

\medterm{Hemilaryngectomy} Hemilaryngectomy is surgery to remove part of the larynx, usually to treat a selected laryngeal cancer while preserving as much speech and swallowing function as possible.

\medterm{Hemimandibulectomy} Hemimandibulectomy is surgery to remove one side of the lower jaw, usually for a tumor or severe bone disease, sometimes followed by reconstruction.

\medterm{Hemimegalencephaly} A rare developmental brain malformation in which one cerebral hemisphere is abnormally enlarged and dysplastic, often causing infantile seizures, developmental impairment, and contralateral neurologic signs.

\medterm{Hemimelia} A congenital limb-reduction defect in which part or all of a limb segment is absent, such as fibular or radial hemimelia.

\medterm{Hemiparesis} Weakness on one side of the body, affecting the face, arm, leg, or a combination. It is less severe than paralysis and may result from stroke, brain injury, a tumor, multiple sclerosis, or another nervous-system disorder.

\medterm{Hemipelvectomy} Hemipelvectomy is major surgery that removes part or all of one side of the pelvis, with or without removal of the leg, usually for extensive cancer or severe trauma.

\medterm{Hemiplegia} Complete or near-complete paralysis of one side of the body, including the face, arm, or leg. It usually results from injury to movement pathways in the brain, commonly a stroke, and sudden onset is an emergency.

\medterm{Hemispatial Neglect} Hemispatial neglect is reduced awareness of one side of the body or space, usually after injury to the opposite cerebral hemisphere.

\medterm{Hemispherectomy} Hemispherectomy is surgery that removes or disconnects much of one cerebral hemisphere to treat severe, otherwise uncontrolled epilepsy in selected patients.

\medterm{Hemisternebra} A congenital or developmental partial segment of the sternum, caused by incomplete formation or fusion of sternal elements and sometimes associated with chest-wall anomalies.

\medterm{Hemithorax} One side of the thorax, including the lung, pleural space, ribs, and other structures on that side of the chest.

\medterm{Hemivertebrae} Congenital vertebrae formed from only part of the normal vertebral body, producing wedge-shaped bone and potentially causing congenital scoliosis or spinal-cord abnormalities.

\medterm{Hemizygous} Having only one copy of a gene or chromosome region rather than the usual pair, as occurs for many X-linked genes in people with one X chromosome.

\medterm{Hemochromatosis} A disorder in which too much iron builds up in the body and can damage the liver, pancreas, heart, joints, skin, and hormone-producing organs. It may be inherited or develop from another illness or repeated transfusions.

\medterm{Hemochromatosis Cardiomyopathy} Heart-muscle damage caused by excess iron buildup. The heart may become stiff or weak and cause breathlessness, swelling, abnormal rhythms, or heart failure; treatment aims to remove excess iron and support heart function.

\medterm{Hemocyanin} A copper-containing oxygen-transport protein in many mollusks and arthropods; unlike human hemoglobin, it circulates dissolved in their hemolymph.

\medterm{Hemocytometer} A ruled counting chamber used with a microscope to determine the concentration of cells or other particles in a measured volume of fluid.

\medterm{Hemodialysis} A treatment for kidney failure in which blood passes through a filter outside the body to remove waste and extra fluid before returning to the person. It usually requires repeated sessions and a blood-vessel access; low blood pressure, infection, clotting, and access problems can occur.

\medterm{Hemodialysis Access Care} Daily maintenance protocols and hygiene measures for an arteriovenous fistula, arteriovenous graft, or central venous catheter used for hemodialysis. Includes daily palpation for a vascular thrill, avoiding blood pressure cuffs or constrictive clothing on the access arm, and promptly reporting signs of infection or thrombosis.

\synonyms
Vascular Access Care; Fistula and Graft Maintenance; Dialysis Access Management

\medterm{Hemodialysis Access Procedures} Procedures that create, maintain, repair, or revise vascular access for hemodialysis, including fistula or graft creation, catheter placement, angioplasty, thrombectomy, and treatment of stenosis or infection.

\medterm{Hemodialysis Clearance} The volume of plasma from which hemodialysis removes a solute per unit time; it depends on blood and dialysate flow, membrane properties, treatment time, and the substance’s distribution.

\medterm{Hemodialysis Solution} The dialysate fluid circulated across a semipermeable membrane during hemodialysis to remove waste and help correct electrolyte, acid--base, and fluid abnormalities.

\medterm{Hemodynamics} The movement of blood through the heart and blood vessels, including pressure, flow, heart output, and resistance. Hemodynamic measurements help assess circulation and guide treatment in conditions such as shock, heart failure, and critical illness.

\medterm{Hemofiltration} A blood-cleaning treatment that removes extra water and dissolved waste through a filter. Replacement fluid is added to maintain the right blood volume; it is used in some forms of temporary or long-term kidney support.

\medterm{Hemoglobin} The iron-containing protein in red blood cells that carries oxygen from the lungs to tissues and helps carry carbon dioxide back to the lungs. Low or abnormal hemoglobin can cause anemia or other blood disorders.

\medterm{Hemoglobin A} The main type of hemoglobin in most adults. It is the oxygen-carrying protein in red blood cells and is made from two alpha and two beta protein parts.

\medterm{Hemoglobin A1c} A routine blood test measuring the percentage of red blood cell hemoglobin coated with sugar (glycated hemoglobin). Because red blood cells live for approximately 120 days (3 months), the test reflects average blood sugar levels over the preceding two to three months, unlike daily fingerstick glucose checks that fluctuate from hour to hour. It is the gold standard for diagnosing and monitoring diabetes and prediabetes (normal is below 5.7\%, prediabetes is 5.7\% to 6.4\%, and diabetes is 6.5\% or higher).
\synonyms
HbA1c; Glycated Hemoglobin; Glycosylated Hemoglobin

\medterm{Hemoglobin A1C Variability} Factors that can make hemoglobin A1c less reliable include altered red-cell survival, hemoglobin variants, anemia,
pregnancy, kidney disease, transfusion, and some treatments. An alternative measure of
glucose may be needed when A1c does not reflect recent glycemia.

\medterm{Hemoglobin C} A beta-globin variant caused by a hemoglobin gene mutation; homozygosity can cause mild hemolytic anemia and splenomegaly, while inheritance with hemoglobin S causes a clinically significant sickling disorder.

\medterm{Hemoglobin C Disease} An inherited hemoglobinopathy caused by a beta-globin variant that produces hemoglobin
C. It may cause mild hemolytic anemia, jaundice, splenomegaly, or, when combined with
hemoglobin S, clinically significant sickling disease.

\medterm{Hemoglobin Electrophoresis} A blood test that separates the different types of hemoglobin. It helps identify inherited conditions such as sickle-cell disease, sickle-cell trait, and thalassemia. Results are interpreted with age, pregnancy, blood counts, recent transfusion, and other tests.

\medterm{Hemoglobin H} An unusual form of hemoglobin made from four beta protein parts when the body makes too little alpha hemoglobin. It is unstable and can cause ongoing breakdown of red blood cells in alpha-thalassemia.

\medterm{Hemoglobin Increased} An increased hemoglobin concentration, which may reflect dehydration, chronic hypoxemia, erythrocytosis, or a myeloproliferative neoplasm and must be interpreted with hematocrit and clinical context.

\medterm{Hemoglobin M} A group of abnormal hemoglobins that stabilize iron in the ferric state and can cause congenital methemoglobinemia and cyanosis.

\medterm{Hemoglobin Normal Values} Reference ranges for blood hemoglobin concentration, which vary with age, sex, pregnancy, altitude, laboratory method, and other factors; results should be interpreted using the reporting laboratory's range.

\medterm{Hemoglobin S} An inherited form of hemoglobin that can make red blood cells become stiff and sickle-shaped when oxygen is low. It can cause anemia, painful blockages, and organ damage in sickle-cell disease.

\medterm{Hemoglobin SC Disease} An inherited form of sickle cell disease caused by having one hemoglobin S gene and one hemoglobin C gene. It can cause anemia, pain episodes, and organ complications.
\synonyms
HbSC; Sickle-Hemoglobin C Disease

\medterm{Hemoglobin Trait} The inherited presence of a hemoglobin variant in a heterozygous state, usually causing no symptoms but potentially affecting screening results or reproductive risk.

\medterm{Hemoglobinopathies: Thalassemias} Thalassemias are inherited disorders of globin chain synthesis imbalance:
alpha-thalassemia (HBA1/HBA2 deletions) and beta-thalassemia (HBB mutations).

\medterm{Hemoglobinopathy} An inherited disorder affecting the structure or production of hemoglobin. Examples include sickle-cell disease and thalassemia; effects range from no symptoms to anemia, painful blood-vessel blockages, and organ damage.

\medterm{Hemoglobinuria} Haemoglobinuria is free haemoglobin in urine, usually from intravascular haemolysis; urine may be dipstick-positive for blood with few or no red cells, and causes include transfusion reactions, malaria, and paroxysmal nocturnal haemoglobinuria.

\medterm{Hemoglobinuria Test} A laboratory test evaluating urine for hemoglobin or related findings, used to investigate intravascular hemolysis, hemoglobinuria, transfusion reactions, or other causes of dark urine.

\medterm{Hemolysin} A toxin or immune substance that damages the outer covering of red blood cells and makes them break open. Some bacteria produce hemolysins, and the resulting destruction of red cells is called hemolysis.

\medterm{Hemolysis} The early destruction of red blood cells, releasing hemoglobin into the blood. It can occur in the bloodstream or in the spleen and liver and may cause anemia, jaundice, dark urine, tiredness, or shortness of breath.

\medterm{Hemolytic} Relating to hemolysis, the premature destruction of red blood cells, which can release hemoglobin and produce anemia or jaundice.

\medterm{Hemolytic Anemia} Anemia caused by red blood cells being destroyed faster than the body can replace them. It may result from an inherited cell problem, an immune reaction, infection, medicine, toxin, or another illness and can cause jaundice and dark urine.

\medterm{Hemolytic Anemia Caused By Chemicals And Toxins} Anemia caused by accelerated destruction of red blood cells after exposure to a chemical or toxin; evaluation identifies hemolysis and removes or treats the responsible exposure.

\medterm{Hemolytic Anemias: Autoimmune} Anemias in which the immune system mistakenly destroys red blood cells. Warm-antibody and cold-antibody forms differ in triggers and temperature sensitivity; symptoms may include tiredness, jaundice, dark urine, and breathlessness.

\medterm{Hemolytic Anemias: Hereditary} Inherited conditions in which red blood cells are destroyed too quickly. Hereditary spherocytosis is one example; it can cause anemia, jaundice, an enlarged spleen, and gallstones. Other inherited red-cell or enzyme problems can cause similar anemia.

\medterm{Hemolytic Crisis} A sudden increase in red-cell destruction causing worsening anemia, jaundice, dark urine, pain, or other manifestations of hemolysis.

\medterm{Hemolytic Disease Of The Newborn} Anemia and jaundice caused when antibodies from the pregnant person cross the placenta and destroy the fetus’s or newborn’s red blood cells. Severe disease can cause widespread fluid buildup or brain injury without prompt treatment.

\medterm{Hemolytic Index} A laboratory-quality indicator estimating hemolysis in a blood specimen; substantial in-vitro hemolysis can interfere with measurement of several analytes.

\medterm{Hemolytic Jaundice} Jaundice resulting from excessive hemolysis (breakdown of red blood cells) causing
overproduction of unconjugated bilirubin that exceeds the liver's capacity for
conjugation and excretion. See Jaundice for complete overview.

\medterm{Hemolytic Transfusion Reaction} An immune or nonimmune complication in which transfused red blood cells are destroyed, potentially causing fever, pain, hemoglobinuria, kidney injury, shock, or disseminated intravascular coagulation.

\medterm{Hemolytic Uremic Syndrome} A serious illness in which tiny blood clots damage small blood vessels, causing red blood cells to break apart, platelets to fall, and the kidneys to become injured. It often follows severe diarrheal infection but can have other causes.

\medterm{Hemolytic-Uremic Syndrome} A serious illness in which small blood-vessel clots destroy red blood cells, lower the platelet count, and injure the kidneys. It often follows severe diarrhea caused by certain bacteria, but inherited immune problems and other triggers can also cause it.

\medterm{Hemoperfusion} A treatment in which blood passes through a material such as charcoal or resin that traps selected poisons or medicines. It is used only for certain severe overdoses because effectiveness depends on the substance and timing.

\medterm{Hemopericardium} Accumulation of blood in the pericardial sac around the heart. If pressure rises and prevents the heart from filling, cardiac tamponade can occur and requires emergency treatment.

\medterm{Hemoperitoneum} Blood collecting inside the abdominal cavity, usually from trauma, a ruptured blood vessel or organ, surgery, or a ruptured ectopic pregnancy. It can cause abdominal pain, distention, dizziness, or shock and may require emergency imaging and surgery or another procedure.

\medterm{Hemopexin} A plasma glycoprotein that binds free heme and transports it to the liver, helping limit oxidative injury during hemolysis.

\medterm{Hemophagocytic Lymphohistiocytosis} A rare, life-threatening illness in which the immune system remains excessively active and injures the body’s organs. It may be inherited or triggered by infection, cancer, or immune disease. Persistent fever, enlarged spleen, low blood counts, liver problems, bleeding, or confusion require urgent specialist care.

\medterm{Hemophagocytosis Of Sickle Cells} A finding in which immune cells engulf blood cells in the bone marrow or other tissues during severe inflammation, infection, or sickle-cell complications. It can worsen anemia and low platelet counts and requires specialist evaluation.

\medterm{Hemophilia} An inherited bleeding disorder in which the blood cannot clot properly due to a shortage of essential clotting proteins (most commonly Factor VIII in Hemophilia A or Factor IX in Hemophilia B). Because the condition is carried on the X chromosome, it almost exclusively affects males, while females are usually carriers. Beyond prolonged bleeding from minor cuts, its hallmark is spontaneous, painful bleeding inside deep muscles and large joints (such as the knees, elbows, and ankles), which can lead to chronic joint stiffness and deformity over time. Historically known as the ``royal disease'' and ``bleeder's disease,'' it is managed by replacing the missing clotting factor through intravenous infusions.

\synonyms
Haemophilia; Bleeder's Disease; The Royal Disease; Factor Deficiency Disorder; Classic Hemophilia

\medterm{Hemophilia A} An X-linked recessive deficiency of coagulation factor VIII, causing prolonged bleeding into joints, muscles, and soft tissues. Severity correlates with factor VIII activity levels.

\medterm{Hemophilia B} An X-linked recessive deficiency of coagulation factor IX (Christmas disease), causing prolonged bleeding similar to hemophilia A. Severity correlates with factor IX activity levels.

\medterm{Hemophilus} An older spelling or variant of Haemophilus, a group of small gram-negative bacteria that includes respiratory and genital pathogens.

\medterm{Hemophilus Ducreyi} A gram-negative bacterium that causes chancroid, a sexually transmitted infection with painful genital ulcers and tender lymph nodes.

\medterm{Hemophobia} An excessive or persistent fear of blood, also called hematophobia, that causes significant distress, avoidance, or sometimes a vasovagal faint.

\medterm{Hemopoietic Stem Cell} A blood-forming stem cell found mainly in the bone marrow. It can make copies of itself and develop into red blood cells, white blood cells, or platelets.

\synonyms
Hematopoietic Stem Cell (HSC); Blood-Forming Cell

\medterm{Hemopoietic Tissue} Tissue that produces blood cells, chiefly bone marrow, containing stem cells and supporting cells that generate erythrocytes, leukocytes, and platelets.

\medterm{Hemoprotein} A protein containing a heme group with an iron atom, such as hemoglobin, myoglobin, cytochromes, and many oxidoreductive enzymes.

\medterm{Hemoptysis} Coughing up blood or blood-stained mucus from the airways or lungs. Causes include infection, irritated or widened airways, a blood clot, or cancer. More than a few streaks, breathlessness, or chest pain requires urgent care.

\medterm{Hemorrhage} Abnormal escape of blood from a damaged blood vessel or, less commonly, from the heart. It may be external or internal, acute or chronic, and range from minor bleeding to life-threatening blood loss causing anemia, shock, or organ compression; the cause and site determine management.

\medterm{Hemorrhagic} Pertaining to or characterized by hemorrhage (extravasation of blood from vessels). Used in contexts such as hemorrhagic stroke, hemorrhagic shock, and hemorrhagic fever.

\medterm{Hemorrhagic Anemia} Anemia caused by losing blood faster than the body can replace it. Blood loss may be sudden, such as after an injury, or slow, such as from heavy menstrual or intestinal bleeding.

\medterm{Hemorrhagic Disorder} A condition that causes abnormal bleeding because of platelet number or function, clotting-factor deficiency, excessive fibrinolysis, or blood-vessel fragility.

\medterm{Hemorrhagic Shock} Circulatory failure caused by severe blood loss. It may produce rapid pulse, low blood pressure, cool
skin, reduced urine output, confusion, or collapse. Treatment requires urgent control of
bleeding, restoration of circulating volume, blood products when indicated, and correction
of associated injuries or clotting abnormalities.

\medterm{Hemorrhagic Stroke} A stroke caused by a blood vessel breaking and bleeding in or around the brain. It can cause sudden weakness, speech or vision problems, severe headache, vomiting, confusion, or loss of consciousness. The effects depend on the location and amount of bleeding.

\medterm{Hemorrhagic Thyroid Nodule} A sudden bleed inside a thyroid nodule (a lump in the thyroid gland), usually caused by the rupture of fragile blood vessels within an existing fluid-filled cyst or noncancerous growth. It produces an abrupt, tender swelling in the front of the neck that may radiate pain toward the ear or jaw, sometimes accompanied by throat tightness, hoarseness, or mild difficulty swallowing. Although the sudden enlargement can appear alarming, most hemorrhagic nodules are benign and gradually shrink as the trapped blood is reabsorbed over several weeks. It is diagnosed with a thyroid ultrasound, and needle aspiration can drain the fluid to relieve local pressure and confirm the diagnosis.

\synonyms
Intrathyroidal Hemorrhage; Hemorrhagic Thyroid Cyst; Bleeding Into A Thyroid Nodule; Thyroid Nodule Bleed

\medterm{Hemorrhoid} Singular form; see \textit{Hemorrhoids}.

\medterm{Hemorrhoid Surgery} An operation to remove or reduce symptomatic internal or external hemorrhoids when bleeding, prolapse, thrombosis, or pain persists despite appropriate conservative treatment.

\medterm{Hemorrhoidal Hemorrhage} Rectal bleeding from enlarged or injured hemorrhoidal vessels, usually bright red and associated with defecation; heavy, persistent, painful, or unexplained bleeding requires evaluation for other causes.

\medterm{Hemorrhoidectomy} An operation to remove enlarged hemorrhoids that continue to bleed, prolapse, clot, or cause severe symptoms despite other treatment. Recovery may include significant pain, bleeding, difficulty passing urine, or temporary bowel discomfort.

\medterm{Hemorrhoids} Swollen, engorged vascular cushions and surrounding connective tissue in the lower rectum and anal canal, commonly triggered by chronic straining, constipation, pregnancy, or heavy lifting. They are classified by their position relative to the pain-sensitive dentate line: internal hemorrhoids (above the dentate line) are typically painless and manifest with bright red rectal bleeding or mucosal prolapse during defecation; external hemorrhoids (below the dentate line) are covered by sensitive skin and can develop acute, exquisitely painful blood clots (thrombosed external hemorrhoids). Management ranges from dietary fiber, stool softeners, and sitz baths to office procedures (rubber band ligation) and surgical hemorrhoidectomy for advanced disease.
\synonyms{Piles; Hemorrhoidal Disease; Haemorrhoids}

\medterm{Hemosiderin} An iron-containing pigment left in cells after hemoglobin from red blood cells is broken down. It can build up after repeated bleeding, excessive red-cell destruction, or too much iron in the body.

\medterm{Hemosiderosis} Excess deposition of hemosiderin, an iron-storage pigment, in tissues. It may follow repeated bleeding or transfusions and can affect organs such as the lungs or liver.

\medterm{Hemosiderotic Synovitis} Synovial inflammation with hemosiderin deposition caused by recurrent bleeding into a joint, commonly associated with hemophilia, anticoagulation, trauma, or a vascular synovial lesion.

\medterm{Hemospermia} The presence of blood in semen, usually caused by inflammation or infection of the genitourinary tract and often self-limited, but persistent or recurrent cases require evaluation.

\medterm{Hemostasis} The body’s process for stopping bleeding after a blood vessel is injured. Blood vessels narrow, platelets form a plug, and clotting proteins strengthen it; later, the body removes the clot as healing occurs.

\medterm{Hemostasis And Coagulation} The linked processes that stop bleeding. Blood vessels narrow, platelets make a plug, and clotting proteins form a stronger clot. Problems with these processes can cause excessive bleeding or unwanted clots.

\medterm{Hemostasis Valve} A valve or sealing device used with an intravascular catheter or sheath to limit blood loss while allowing passage of wires or instruments.

\medterm{Hemostat} A surgical clamping instrument with jaws that compress a blood vessel or tissue to control bleeding during a procedure. Many hemostats have a ratcheted locking mechanism that maintains pressure without continuous hand force.

\synonyms
Hemostatic Forceps; Artery Clamp

\medterm{Hemostatic Agent} A medicine or material used to stop or reduce bleeding. It may form a physical seal, help the blood clot, slow clot breakdown, or replace a missing clotting protein.

\medterm{Hemothorax} Blood collecting in the space between the lung and chest wall, usually after injury or surgery. It can cause chest pain, breathlessness, low blood pressure, and reduced breath sounds and often requires urgent drainage and treatment of the bleeding source.

\medterm{Hemovac Drain} A closed-suction surgical drain with a compressible reservoir used to remove blood or other fluid from a wound or operative site.

\medterm{Hendra Virus} A rare zoonotic virus carried by fruit bats that can infect horses and occasionally humans through close contact with infected horses. Human disease may cause severe respiratory or neurologic illness and requires public-health management.

\medterm{Henipavirus} A genus of paramyxoviruses that includes Nipah and Hendra viruses, which can cause severe encephalitis and respiratory disease.

\medterm{Henna} A plant dye used on skin, hair, and textiles. Natural henna is usually orange-brown; so-called black henna may contain para-phenylenediamine, which can cause severe allergic skin reactions and permanent scarring.

\medterm{Henoch-SchöNlein Purpura} An older name for IgA vasculitis, a disease in which small blood vessels become inflamed. It can cause purple spots on the legs or buttocks, joint pain, abdominal pain, and kidney problems.

\medterm{Henoch-Schonlein Purpura} An inflammation of small blood vessels, now called IgA vasculitis, that mainly affects children. It can cause raised purple spots on the legs or buttocks, joint pain, abdominal pain, and kidney inflammation. Most cases improve, but kidney monitoring is important.

\synonyms
IgA Vasculitis; Anaphylactoid Purpura; Schönlein-Henoch Syndrome

\medterm{HEPA Filter} A high-efficiency air filter that removes most very small airborne particles. It can reduce dust, allergens, and some infectious particles when properly fitted and maintained, but does not remove every gas or chemical.

\synonyms
High-Efficiency Particulate Air Filter

\medterm{Hepar} Latin for liver; it is used in medical terminology and in names of substances or preparations relating to the liver.

\medterm{Heparan Sulfate} A sulfated glycosaminoglycan in cell surfaces and extracellular matrix that binds growth factors, regulates signaling, and contributes to filtration-barrier and connective-tissue structure.

\medterm{Heparin} An anticoagulant medicine that reduces blood clotting by enhancing antithrombin activity. It may be given by injection to prevent or treat blood clots.

\medterm{Heparin Binding} The attachment of a protein or other molecule to heparin, a highly charged substance related to blood clotting. Binding can alter the molecule’s activity, movement, or laboratory measurement.

\medterm{Heparin-Induced Thrombocytopenia} An immune reaction to heparin that lowers the platelet count while making dangerous blood clots more likely. It usually begins several days after heparin exposure. Suspected cases require stopping heparin and urgent medical treatment with a different anticoagulant.

\medterm{Hepatectomy} Surgery to remove part or, rarely, nearly all of the liver, usually to treat a tumor or severe localized liver disease. The remaining liver must be large and healthy enough to work. Bleeding, bile leakage, infection, and liver failure are important risks.

\medterm{Hepatic} Relating to the liver. The liver processes nutrients and medicines, makes bile and important blood proteins, stores energy, and helps remove harmful substances; hepatic disease can affect one or several of these functions.

\medterm{Hepatic Artery} The main artery supplying oxygen-rich blood to the liver. It divides into branches that deliver blood to the right and left parts of the organ.

\medterm{Hepatic Artery Occlusion} Partial or complete blockage of the hepatic artery, usually caused by a blood clot, embolus, inflammation, injury, or a procedure. It may reduce blood flow to the liver and cause liver infarction or ischemic hepatitis, especially when portal-vein blood flow is also reduced.

\medterm{Hepatic Biopsy} Alternate name; see \textit{Liver Biopsy}.

\medterm{Hepatic Coma} Severe loss of consciousness caused by advanced liver failure and buildup of toxic substances in the blood. It is a medical emergency, often preceded by confusion, sleepiness, personality change, or hand-flapping movements.

\medterm{Hepatic Cyst} A fluid-filled sac in the liver. Most simple liver cysts are harmless, cause no symptoms, and are found by chance on imaging. Large, infected, bleeding, parasitic, or complex cysts may cause pain, pressure, fever, or other complications.

\synonyms
Simple Liver Cyst; Biliary Cyst

\medterm{Hepatic Dose Adjustment} Changing a medicine or its dose because liver disease may slow its breakdown or removal. The correct adjustment depends on the medicine, liver function, other illnesses, and monitoring of benefit and toxicity.

\medterm{Hepatic Duct} The duct formed by the union of the right and left liver ducts, carrying bile from the liver toward the common bile duct and intestine.

\medterm{Hepatic Encephalopathy} Confusion or changes in alertness, behavior, movement, or consciousness caused by severe liver disease and buildup of substances that affect the brain. Symptoms can progress to coma and may be precipitated by factors such as infection, bleeding, constipation, or medicines.

\medterm{Hepatic Fissure} A natural groove or cleft on the liver's surface. It separates parts of the liver and may contain or provide a route for blood vessels, bile ducts, or supporting ligaments.

\medterm{Hepatic Flexure} The bend where the right-sided ascending colon turns into the transverse colon near the liver. It is a normal anatomical landmark that may be seen during imaging, surgery, or colonoscopy.

\medterm{Hepatic Granulomas} Small collections of immune cells that form in the liver in response to infections, medicines, sarcoidosis, or other inflammatory conditions. They often cause no symptoms, but the underlying condition may cause abnormal liver tests, bile-duct blockage, fibrosis, or portal hypertension.

\medterm{Hepatic Hemangioma} A common, usually harmless growth made from small blood vessels in the liver. Most cause no symptoms and are found by chance on imaging; larger ones may cause discomfort or rarely bleeding.

\medterm{Hepatic Hydrothorax} A buildup of fluid around a lung caused by severe liver disease and high pressure in the abdomen. It is usually on the right side and can cause breathlessness; treatment focuses on the liver disease and fluid control.

\medterm{Hepatic Ischemia} Liver injury caused by inadequate hepatic blood flow or oxygen delivery, often during
shock, severe hypoxemia, cardiac failure, or hepatic vascular occlusion. It can produce
a rapid, marked rise in aminotransferases and may progress to acute liver failure.

\medterm{Hepatic Lamina} A sheet of cells in a developing embryo that contributes to the liver and its bile ducts.

\medterm{Hepatic Necrosis} Death of liver cells caused by severe ischemia, toxins, viral infection, autoimmune injury, or other insults; extensive necrosis can produce acute liver failure.

\medterm{Hepatic Portal System} A network of veins that carries blood from the digestive organs, spleen, and pancreas to the liver before the blood returns to the heart. The liver processes nutrients, medicines, and harmful substances in this blood.

\medterm{Hepatic Steatosis} Excess accumulation of fat in liver cells, commonly called fatty liver; causes include metabolic disease, alcohol, medicines, and other conditions, and inflammation or fibrosis may develop in some people.

\medterm{Hepatic Surgery} Surgical treatment of liver disease, including partial liver resection, transplantation, and
procedures for tumors or biliary disorders. The operation depends on the diagnosis, tumor extent, liver function, and
overall health.

\medterm{Hepatic Synthetic Dysfunction} Reduced ability of the liver to make important proteins, including albumin and blood-clotting factors. It can cause swelling, easy bleeding, or abnormal clotting tests, although other problems such as vitamin K deficiency or blood-thinning medicines can also alter these tests.

\medterm{Hepatic Trauma} Liver injury caused by blunt or penetrating trauma, potentially producing internal bleeding and requiring urgent imaging, monitoring, or surgery based on severity.

\medterm{Hepatic Vein} One of the veins that carries blood away from the liver and into the large vein leading to the heart. The main hepatic veins drain different regions of the liver.

\medterm{Hepatic Vein Obstruction} Obstruction of hepatic venous outflow, usually from thrombosis or compression, causing abdominal pain, hepatomegaly, ascites, and potentially liver failure.

\medterm{Hepatitis} Hepatitis is inflammation of the liver caused by viruses, alcohol, medications, toxins, metabolic disease, or autoimmunity and may be acute or chronic.

\medterm{Hepatitis A} A short-term liver infection caused by hepatitis A virus. It spreads when virus from infected stool is swallowed, often through close contact or contaminated food or water. It may cause fever, tiredness, nausea, abdominal discomfort, or jaundice. Most people recover completely; chronic infection is not expected. Vaccination helps prevent it.

\medterm{Hepatitis A Virus} The virus that causes hepatitis A. It spreads when virus from infected stool is swallowed, commonly through contaminated food, water, or close contact. It causes a short-term liver infection that may produce fever, nausea, tiredness, dark urine, pale stools, or jaundice; vaccination helps prevent it.

\medterm{Hepatitis B} A liver infection caused by the hepatitis B virus (HBV). It spreads when infected blood, semen, or certain other body fluids enter another person’s body, including through sex, shared needles or equipment, or from an infected parent to a baby around the time of birth. Infection may be acute (short-term) or chronic (long-term). Many people have no symptoms. When symptoms occur, they may include tiredness, fever, loss of appetite, nausea, abdominal pain, joint pain, dark urine, pale stools, or yellow skin and eyes (jaundice). Chronic infection, especially when acquired in infancy or childhood, can cause liver damage, cirrhosis, liver failure, or liver cancer. Vaccination prevents hepatitis B.

\medterm{Hepatitis B Virus} The virus that causes hepatitis B. HBV spreads when infected blood, semen, or certain other body fluids enter another person’s body, including through sex, shared needles or equipment, or around childbirth. It can cause acute (short-term) or chronic (long-term) liver infection. Chronic infection can cause liver damage, cirrhosis, liver failure, or liver cancer. Vaccination prevents hepatitis B.

\medterm{Hepatitis C} A liver infection caused by the hepatitis C virus (HCV). It spreads mainly when infected blood enters another person’s body, especially through shared injection equipment; it can also spread through some other blood exposures and, less often, through sex or from an infected parent to a baby. Acute infection often causes no symptoms, and some people clear the virus without treatment. In many others it becomes chronic (long-term) and can cause liver damage, cirrhosis, liver failure, or liver cancer. Direct-acting antiviral medicines cure more than 95\% of infections in most people; no vaccine is available.

\medterm{Hepatitis C Virus} The virus that causes hepatitis C. HCV spreads mainly when infected blood enters another person’s body, especially through shared needles or other injection equipment. It can also spread through some medical exposures, from an infected parent to a baby, or less often through sex. Many people have no symptoms. The infection may be acute (short-term) or become chronic (long-term), which can cause liver damage, cirrhosis, liver failure, or liver cancer. Modern antiviral medicines cure more than 95\% of infections in most people, but no vaccine is available.

\medterm{Hepatitis D} A liver infection caused by the hepatitis D virus (HDV). HDV can infect only people who also have hepatitis B because it needs hepatitis B virus (HBV) to reproduce. Infection with both viruses at the same time is called coinfection; getting HDV after already having HBV is called superinfection. Superinfection is more likely to cause long-term liver disease, cirrhosis, liver failure, or death.

\medterm{Hepatitis D} Another name for hepatitis D, a liver infection caused by HDV. HDV needs HBV to reproduce. It may occur as coinfection with HBV or as superinfection in someone who already has HBV; superinfection is more likely to cause long-term liver disease, cirrhosis, liver failure, or death.

\medterm{Hepatitis D Virus} A virus that can multiply only when hepatitis B virus is also present. It spreads through infected blood or body fluids and can make hepatitis B more severe, increasing the risk of cirrhosis and liver cancer. Hepatitis B vaccination also prevents hepatitis D.

\medterm{Hepatitis E} A liver infection caused by the hepatitis E virus (HEV). It usually spreads when people swallow water or food contaminated with stool, and it can also spread through undercooked meat from infected animals. It usually causes a short-term illness, often with tiredness, nausea, vomiting, abdominal pain, dark urine, pale stools, or jaundice. Severe illness is more likely during pregnancy and in some people with weakened immune systems. Rarely, infection becomes chronic in people with weakened immunity. No vaccine is available in the United States.

\medterm{Hepatitis E Virus} The virus that causes hepatitis E. HEV usually spreads through water or food contaminated with stool and can also spread through undercooked meat from infected animals. It commonly causes short-term liver infection, but it may cause severe illness during pregnancy and can occasionally become chronic in people with weakened immune systems. No vaccine is available in the United States.

\medterm{Hepatitis Serology} Blood testing for antibodies, antigens, and other markers of hepatitis viruses. The results may show a current or past infection, immunity from vaccination, or whether infection has become chronic. Different viruses require different tests. For example, hepatitis B testing commonly includes hepatitis B surface antigen and antibodies to hepatitis B surface and core antigens; hepatitis C testing usually starts with an antibody test and uses an HCV RNA test to show current infection. Results depend on the virus, timing, vaccination history, and other laboratory findings, so a single result may not give the full answer.

\medterm{Hepatitis Virus} Any of the viruses that can infect and inflame the liver. The five main hepatitis viruses are hepatitis A, B, C, D, and E. They differ in how they spread and in whether they cause short-term or long-term infection. Hepatitis A usually causes a short-term illness; hepatitis B and C can become chronic and damage the liver; hepatitis D occurs only with hepatitis B; and hepatitis E usually causes short-term illness but can be severe during pregnancy or in people with weakened immunity. Vaccines prevent hepatitis A and B, and hepatitis B vaccination also prevents hepatitis D.

\medterm{Hepatitis Virus Panel} A group of blood tests, usually done on one blood sample, that looks for signs of hepatitis A, B, and C infection. The tests may detect viral antigens or antibodies and may show a current infection, a past infection, or immunity after vaccination. The usual panel does not automatically test for hepatitis D or E, which require separate tests when appropriate. A positive antibody result does not always mean that the virus is still present; additional tests may be needed to confirm current infection.

\medterm{Hepatitis Viruses} The five main viruses that infect and inflame the liver: hepatitis A, B, C, D, and E. Hepatitis A and E usually cause short-term infection, although hepatitis E can be severe during pregnancy and may become chronic in some people with weakened immunity. Hepatitis B and C can become chronic and cause cirrhosis or liver cancer. Hepatitis D occurs only in people who also have hepatitis B. Vaccines prevent hepatitis A and B, hepatitis B vaccination also prevents hepatitis D, and effective medicines can cure hepatitis C; no vaccine is available for hepatitis C or E in the United States.

\medterm{Hepatobiliary} Relating to the liver, gallbladder, and bile ducts. The hepatobiliary system makes bile in the liver, stores and concentrates it in the gallbladder, and carries it through the bile ducts to the small intestine, where it helps digest fats. It also includes the liver’s roles in processing nutrients, medicines, and waste products.

\medterm{Hepatobiliary Scintigraphy} A nuclear-medicine scan, often called a HIDA scan, that uses a small injected radioactive tracer to follow bile movement through the liver, bile ducts, gallbladder, and small intestine. A special camera records the tracer over time. The scan can show blocked bile flow, whether the gallbladder fills and empties normally, bile leaks, and findings that support acute gallbladder inflammation or other biliary disorders.

\synonyms
Biliary Scintigraphy; Cholescintigraphy; HIDA Scan; Hepatobiliary Iminodiacetic Acid Scan

\medterm{Hepatobiliary System} The liver, gallbladder, and bile ducts. The liver makes and secretes bile, the gallbladder stores and concentrates it, and the bile ducts carry it to the small intestine, where it helps digest fats and carry some waste products out of the liver. The liver also processes nutrients and many medicines and toxins.

\medterm{Hepatoblastoma} A rare cancer that begins in the liver and is the most common type of liver cancer in children. It usually affects children younger than 3 years. A child may have no symptoms until the tumor grows larger; possible findings include a lump or swelling in the abdomen, abdominal pain, poor appetite, nausea, vomiting, or unexplained weight loss. Evaluation may include imaging and blood tests, especially the tumor marker alpha-fetoprotein (AFP); a tissue sample may be used when needed to confirm the diagnosis.

\medterm{Hepatocellular Adenoma} A rare, usually noncancerous tumor made of mature liver cells. It occurs most often in women of reproductive age and is associated with estrogen exposure, although it can also occur with anabolic steroid use, certain inherited metabolic diseases, or without a known cause. It is often found by chance on imaging and may cause no symptoms. Larger tumors can bleed or rupture, and a small number can change into liver cancer, especially certain subtypes.

\medterm{Hepatocellular Carcinoma} A cancer that begins in liver cells called hepatocytes. It is the most common primary liver cancer in adults. It often develops in a liver affected by long-term disease, especially cirrhosis, chronic hepatitis B or C, fatty liver inflammation, or long-term heavy alcohol use, but it can occur without a known risk factor. Early disease may cause no symptoms. Later findings may include a lump or pain in the upper-right abdomen, abdominal swelling, jaundice, tiredness, poor appetite, or unexplained weight loss. Diagnosis may use imaging, liver blood tests, and the tumor marker alpha-fetoprotein (AFP); a tissue sample is sometimes needed.

\medterm{Hepatocellular Jaundice} Yellowing of the skin or eyes caused by damaged liver cells processing or releasing bilirubin poorly. Causes include hepatitis, alcohol, medicines, toxins, fatty liver disease, and other liver disorders.

\synonyms
Hepatocellular Jaundice; Hepatocellular Icterus

\medterm{Hepatocellular Jaundice} Jaundice caused by impaired liver cells. Injured hepatocytes cannot take up, change into a water-soluble form, or remove bilirubin normally, so bilirubin builds up in the blood and makes the skin and whites of the eyes yellow. Urine may become dark and stools may become pale; itching, tiredness, nausea, or abdominal discomfort may also occur. Causes include viral or autoimmune hepatitis, alcohol-related liver disease, medicines or toxins, cirrhosis, hemochromatosis, and Wilson disease.

\medterm{Hepatocerebral Degeneration} Usually referring to acquired hepatocerebral degeneration, a rare long-term brain disorder linked to advanced liver disease or a portosystemic shunt, an abnormal connection that lets blood bypass the liver. Harmful substances that are normally cleared by the liver, including excess manganese and other neurotoxins, can affect brain areas involved in movement and thinking. Symptoms may include tremor, slowed movement or parkinsonism, poor coordination, unsteady walking, speech changes, abnormal movements, mood or personality changes, and problems with thinking or memory. It is different from the short-lived episodes of hepatic encephalopathy, although the conditions can overlap.

\medterm{Hepatocyte} The main working cell of the liver. Hepatocytes process carbohydrates, fats, proteins, medicines, and other substances from the blood. They make bile, albumin, clotting factors, and other blood proteins; store energy as glycogen; change ammonia into urea; and help remove or break down harmful substances. Injury to hepatocytes can reduce these functions and raise liver enzymes in the blood.

\medterm{Hepatojugular Reflux} A bedside physical sign, also called abdominojugular reflux. Firm, steady pressure on the upper abdomen briefly increases the amount of blood returning to the heart. In a normal response, the jugular veins in the neck rise briefly and then fall while the pressure continues. A sustained rise in jugular venous pressure suggests elevated right-sided heart filling pressure or limited right-ventricle function, often in heart failure. It is a clinical finding, not a diagnosis by itself.

\medterm{Hepatolenticular Degeneration} Another name for Wilson disease, a rare inherited disorder caused by changes in the \textit{ATP7B} gene. The altered gene prevents the liver from removing excess copper into bile, so copper builds up to harmful levels in the liver, brain, eyes, kidneys, and other tissues. It may cause hepatitis, cirrhosis, or liver failure; tremor, stiffness, poor coordination, speech or mood changes; and green, gold, or brown Kayser--Fleischer rings around the cornea. The disorder is inherited in an autosomal recessive pattern.

\medterm{Hepatologist} A physician who specializes in diseases of the liver and its connected bile ducts and gallbladder. Hepatologists commonly diagnose and manage hepatitis, fatty liver disease, cirrhosis, jaundice, liver failure, liver tumors, and complications of chronic liver disease. Some also care for related diseases of the pancreas and people being assessed for liver transplantation. Hepatology is usually a subspecialty of gastroenterology.

\medterm{Hepatology} The branch of medicine concerned mainly with the liver and also with the gallbladder, bile ducts, and related pancreas disorders. It covers how these organs work, how their diseases are diagnosed and managed, and complications such as hepatitis, fatty liver disease, cirrhosis, jaundice, bile-flow disorders, liver cancer, and liver failure. Hepatology is usually a subspecialty of gastroenterology.

\medterm{Hepatoma} An older, nonspecific term for a liver tumor. In some contexts it is used to mean hepatocellular carcinoma, a cancer arising from liver cells, but the word alone does not show whether a tumor is benign or malignant. \textit{Hepatocellular carcinoma} is the preferred term when a malignant liver-cell tumor is meant.

\medterm{Hepatomegaly} An enlarged liver. It may be noticed when the liver is felt below the right ribs during an examination or seen on ultrasound, CT, or MRI. Hepatomegaly is a physical finding, not a disease itself. Causes include acute hepatitis, metabolic dysfunction-associated steatotic liver disease, alcohol-related liver disease, congestion from heart or blood-flow problems, bleeding into the liver, bile-duct blockage, storage or infiltrative diseases, and primary or metastatic cancer. It may cause no symptoms or may be associated with right-upper-abdominal discomfort.

\medterm{Hepatopulmonary Syndrome} A complication of liver disease or portal hypertension in which blood vessels inside the lungs become abnormally widened. Blood then passes through the lungs without receiving enough oxygen, causing low oxygen levels in the blood. Symptoms may include shortness of breath, bluish skin, or clubbing of the fingers. Breathing difficulty and low oxygen may become worse when sitting or standing and improve when lying down; this pattern is called platypnea and orthodeoxia. The condition is identified by liver disease or portal hypertension, impaired oxygenation, and evidence of widened vessels inside the lungs.

\medterm{Hepatorenal Syndrome} Serious kidney dysfunction that occurs mainly in advanced cirrhosis, usually with portal hypertension and ascites. Widening of blood vessels in the abdominal circulation lowers the amount of blood effectively reaching the kidneys. The kidneys respond by narrowing their blood vessels and retaining salt and water, so filtration falls even though the kidneys may not initially have a primary structural disease. Hepatorenal syndrome may cause acute kidney injury (HRS-AKI), which develops rapidly, or more persistent kidney dysfunction. Diagnosis requires excluding other important causes of kidney injury, such as shock, dehydration, medicines, infection, or intrinsic kidney disease.

\medterm{Hepatosplenomegaly} Enlargement of both the liver and spleen at the same time. It is a physical finding rather than a disease itself and may be found during an examination or on imaging. Causes include infections, chronic liver disease with portal hypertension, blood-cell destruction, blood cancers, inherited storage or metabolic disorders, and other infiltrative or inflammatory diseases. It may cause abdominal fullness or discomfort, or no symptoms; other findings depend on the underlying cause. Blood tests and imaging may help identify the cause.

\medterm{Hepatotoxicity} Toxic injury to the liver caused by a medicine, herbal product, dietary supplement, poison, chemical, or other exposure. The injury may mainly damage liver cells (hepatocellular injury), slow or block bile flow (cholestatic injury), or show both patterns. It may cause no symptoms but can lead to tiredness, nausea, poor appetite, abdominal pain, itching, dark urine, or jaundice, with abnormal liver blood tests. Most cases improve after the cause is removed, but severe injury can cause acute liver failure.

\medterm{Hepatovirus} A genus in the Picornaviridae family of small RNA viruses. Its best-known human member is hepatitis A virus (HAV), which causes hepatitis A and usually spreads by the fecal--oral route, when virus from stool reaches the mouth through contaminated food, water, hands, or close contact. Other hepatoviruses infect different animal hosts and are not established causes of human disease.

\medterm{Hepcidin} A peptide hormone made mainly by liver cells and the main regulator of iron balance. When iron stores are high or inflammation is present, hepcidin binds to ferroportin, the protein that exports iron from intestinal cells, macrophages, and liver cells. Ferroportin is then removed, reducing iron absorption from food and the release of recycled or stored iron into the blood. When iron is low or the bone marrow needs more iron to make red blood cells, hepcidin levels usually fall.

\medterm{Heptachlor} A persistent organochlorine insecticide formerly used mainly on crops and for termite control. It breaks down slowly into heptachlor epoxide, which can remain in soil and the food chain for a long time. Exposure can occur through contaminated soil, food, water, air, or contact with old treated materials. Studies, especially in animals, associate exposure with effects on the liver, nervous system, reproductive system, and development; evidence of health effects in humans is limited. Most uses have been canceled or severely restricted in the United States.

\medterm{HER2} Short for human epidermal growth factor receptor 2, a protein on the surface of cells that helps control normal cell growth. It is made from the \textit{ERBB2} gene. Some cancer cells have too much HER2 protein or extra copies of the \textit{ERBB2} gene; these cancers may grow more quickly and spread more easily. Testing HER2 protein levels or gene copies can help classify a cancer and identify whether HER2-targeted treatment may be useful.

\medterm{HER2-Positive} A cancer whose cells have increased expression or amplification of human epidermal growth factor receptor 2 (HER2). Some HER2-targeted medicines may be used according to the cancer type and clinical context.

\medterm{HER2-Positive Breast Cancer} Breast cancer cells with unusually high levels of the HER2 growth receptor or extra copies of its gene. This can make the cancer grow faster, but it also allows treatment with medicines that specifically block HER2.

\medterm{Herbal Remedies And Supplements For Weight Loss} Plant-based weight-loss products have variable evidence, uncertain composition, and potential stimulant, liver, cardiovascular, or drug-interaction risks; they should not replace evidence-based obesity treatment.

\medterm{Herbal Supplement} A concentrated plant-derived product marketed for health or symptom relief; composition and evidence vary, and contamination, drug interactions, liver injury, and pregnancy risks must be considered.

\medterm{Herbaspirillum} A genus of gram-negative bacteria found in soil and plants; human infection is rare but can occur opportunistically.

\medterm{Herbicide} A chemical used to kill or control unwanted plants; exposure can irritate skin, eyes, lungs, or the gastrointestinal tract, with toxicity determined by the specific compound and dose.

\medterm{Herd Immunity} Protection that occurs when enough people in a community are immune to an infection that it spreads less easily. The proportion needed varies widely by infection, and herd immunity does not protect everyone equally.

\medterm{Hereditary} Transmitted through genetic material from a parent or arising from an inherited variant; a hereditary trait may be dominant, recessive, X-linked, mitochondrial, or multifactorial.

\medterm{Hereditary Amyloidosis} A group of inherited disorders in which mutations cause misfolded proteins to form amyloid deposits in organs such as the nerves, heart, kidneys, or gastrointestinal tract.

\medterm{Hereditary Angioedema} An inherited, usually bradykinin-mediated disorder causing repeated episodes of swelling in the skin, hands, feet, stomach or bowel wall, or throat. Types 1 and 2 commonly result from deficiency or dysfunction of C1 inhibitor, often due to variants in the \textit{SERPING1} gene; other forms may have normal C1-inhibitor levels. Attacks may cause severe abdominal pain or swelling that narrows the upper airway and often occur without itchy hives.

\synonyms
HAE; Hereditary Angioneurotic Edema; HANE; C1-INH Deficiency Angioedema

\medterm{Hereditary Cancer} Cancer occurring because a person inherited a disease-causing genetic variant that raises cancer risk. Genetic counseling and appropriate risk-based screening may be recommended.

\medterm{Hereditary Coproporphyria} An inherited disorder of a chemical pathway used to make heme, the oxygen-carrying part of hemoglobin. Attacks may cause severe abdominal pain, vomiting, nerve or mental symptoms, and sometimes blistering skin after sunlight exposure.

\medterm{Hereditary Diffuse Gastric Cancer} An inherited condition, often caused by a harmful CDH1 gene change, that greatly increases the risk of a hard-to-detect stomach cancer and also raises the risk of lobular breast cancer. Genetic counseling, specialist screening, and discussion of risk-reducing stomach surgery may be recommended.

\synonyms
HDGC; Familial Gastric Cancer Syndrome

\medterm{Hereditary Elliptocytosis} An inherited red-cell membrane disorder in which erythrocytes are elliptical on the blood smear; most people are asymptomatic, but hemolytic anemia can occur in severe forms.

\medterm{Hereditary Fructose Intolerance} An inherited disorder in which the body cannot properly break down fructose, a sugar in fruit, honey, and table sugar. Eating it can cause vomiting, sweating, low blood glucose, abdominal pain, or liver injury.

\medterm{Hereditary Hemorrhagic Telangiectasia} An inherited disorder in which fragile, abnormal blood vessels form in the nose, skin, lungs, brain, or liver. Recurrent nosebleeds and anemia are common; bleeding or abnormal vessels in internal organs may be serious.

\synonyms
HHT; Osler-Weber-Rendu Syndrome

\medterm{Hereditary Hemorrhagic Telangiectasia} An inherited blood-vessel disorder causing fragile vessels that can bleed, especially in the nose and digestive tract. Abnormal vessels may also form in the lungs, brain, or liver.

\synonyms
HHT; Osler-Weber-Rendu Syndrome

\medterm{Hereditary Hyperbilirubinemia} An inherited disorder of bilirubin production, conjugation, or excretion that causes persistent or intermittent jaundice, such as Gilbert, Crigler--Najjar, Dubin--Johnson, or Rotor syndrome.

\medterm{Hereditary Leiomyomatosis And Renal Cell Cancer} An inherited condition that can cause painful skin muscle tumors, early or multiple uterine fibroids, and an aggressive type of kidney cancer. It is usually related to an FH gene change; genetic counseling and regular specialist skin, uterine, and kidney assessment are important.

\synonyms
HLRCC; Reed Syndrome

\medterm{Hereditary Multiple Exostoses} An inherited disorder in which many benign, cartilage-capped bone growths called osteochondromas form near the ends of growing bones. They usually develop during childhood or adolescence and stop forming when skeletal growth ends. They may cause pain, limited joint movement, unequal limb lengths, short stature, bone deformity, or pressure on nearby nerves and blood vessels. The condition is usually caused by a change in the \textit{EXT1} or \textit{EXT2} gene and follows an autosomal dominant inheritance pattern. Most growths remain benign, but a small number can change into a cancerous cartilage tumor. Also called hereditary multiple osteochondromas.

\synonyms
Hereditary Multiple Osteochondromas; Diaphyseal Aclasis

\medterm{Hereditary Ovalocytosis} An inherited red-cell membrane disorder in which many erythrocytes are oval or elliptical; it is often mild but can cause hemolytic anemia in some individuals.

\medterm{Hereditary Paraganglioma-Pheochromocytoma Syndrome} An inherited condition that increases the risk of tumors arising from nerve-related cells near blood vessels or in the adrenal glands. Some tumors release hormones that cause high blood pressure, headache, sweating, or palpitations; genetic counseling and lifelong imaging and hormone monitoring may be needed.

\synonyms
Hereditary Paraganglioma; SDH-Associated Pheochromocytoma

\medterm{Hereditary Persistence of Fetal Hemoglobin} An inherited condition in which production of fetal hemoglobin continues into adulthood. It is usually harmless and may lessen the severity of some beta-globin disorders such as sickle-cell disease.

\medterm{Hereditary Spastic Paraplegia} A group of inherited neurologic disorders causing progressive stiffness and weakness of the legs because of corticospinal-tract degeneration. Some forms also affect speech, vision, sensation, coordination, or other organs.

\medterm{Hereditary Spherocytic Anemia} An inherited red-cell membrane disorder, usually hereditary spherocytosis, in which fragile spherical erythrocytes are destroyed prematurely, causing hemolytic anemia, jaundice, and splenomegaly.

\medterm{Hereditary Spherocytosis} An inherited blood disorder in which red blood cells are round rather than disc-shaped and are broken down too quickly, mainly in the spleen. It may cause anemia, jaundice, gallstones, and an enlarged spleen.

\medterm{Hereditary Urea Cycle Abnormality} An inherited problem removing ammonia, a waste product, from the blood. Ammonia can build up and cause vomiting, sleepiness, confusion, seizures, coma, and brain injury, especially during illness or fasting.

\medterm{Heredity} The passing of genetic characteristics from parents to children through genes and chromosomes. Heredity influences appearance, body function, disease risk, medicine response, and susceptibility to some inherited conditions.

\medterm{Heredo-Ataxia} An inherited disorder causing gradually worsening problems with balance, walking, coordination, or speech. Different gene changes cause different forms, so diagnosis usually includes neurologic assessment and genetic testing.

\medterm{Heredopolyneuritis} An inherited disorder affecting many peripheral nerves. It may cause gradually worsening weakness, muscle wasting, numbness, reduced reflexes, or difficulty walking, often beginning in the feet and lower legs.

\medterm{Hering-Breuer Reflex} An automatic lung response that slows or stops breathing in when the lungs are stretched too much. It helps protect the lungs from overinflation, especially during large breaths.

\medterm{Hermanski-Pudlak Syndrome} Alternate spelling; see \textit{Hermansky-Pudlak Syndrome}.

\medterm{Hermansky-Pudlak Syndrome} A group of inherited disorders that can cause reduced skin, hair, and eye pigment; easy bleeding because platelets do not work normally; and, in some types, lung scarring, bowel inflammation, or low white-cell counts. Genetic and specialized blood testing help confirm the diagnosis.

\medterm{Hermaphroditism} An outdated and potentially stigmatizing term historically used for differences in sex development. Modern clinical language generally uses intersex or a specific difference of sex development, based on the person's anatomy, chromosomes, hormones, and lived identity.

\medterm{Hernia} A bulge formed when an organ or tissue pushes through a weak area in the muscle or connective tissue that normally holds it in place. Hernias commonly occur in the abdominal wall and may become trapped or lose their blood supply, causing severe pain.

\medterm{Hernia Mesh} A sheet of synthetic or biologic material used to strengthen the body wall during some hernia repairs. It may reduce recurrence, but complications can include infection, persistent pain, movement of the mesh, or another hernia.

\medterm{Hernia Repair} Surgery that returns the displaced tissue to its proper place and closes or strengthens the weak area in the body wall. Repair may use stitches, mesh, or a minimally invasive approach; the choice depends on the hernia and the person’s health.

\synonyms
Herniorrhaphy; Hernioplasty

\medterm{Herniated Disc} A spinal disk that bulges or tears so that its inner material irritates or presses on a nearby nerve. It can cause back or neck pain, pain spreading into an arm or leg, numbness, tingling, or weakness.

\medterm{Herniated Disk} Displacement of disk material through a tear in the outer ring of an intervertebral disc, most commonly in the lumbar spine (often at L4-L5 and L5-S1) or the cervical spine, which may press on a nearby nerve and cause pain, numbness, tingling, or weakness.

\medterm{Herniation} Protrusion of an organ, tissue, or structure through an abnormal opening or weak area, as in a herniated disk or abdominal-wall hernia.

\medterm{Herniorrhaphy} An operation to repair a hernia, in which tissue or an organ has pushed through a weak area in the body wall. The surgeon returns the tissue and closes or reinforces the weakness, sometimes with mesh; recurrence, pain, infection, or injury to nearby structures are possible.

\medterm{Heroin} Heroin is a highly addictive opioid converted to morphine in the body; overdose causes respiratory depression, coma, and death, while repeated use causes dependence and withdrawal.

\medterm{Heroin Dependence} Legacy, drug-specific label; see \textit{Opioid Use Disorder}.

\medterm{Heroin Overdose} A life-threatening reaction to too much heroin, causing extreme sleepiness, pinpoint pupils, slow or stopped breathing, blue lips, coma, or death. Call emergency services and give naloxone if available; repeat doses may be needed.

\medterm{Herpangina} A common childhood viral illness causing fever, sore throat, and small blisters or ulcers at the back of the mouth. It spreads through respiratory or stool-contaminated secretions and usually improves on its own.

\medterm{Herpes} Infection caused by a herpesvirus. The term commonly refers to herpes simplex infection, which causes oral or
genital lesions, but other herpesviruses cause distinct diseases.

\medterm{Herpes Encephalitis} Inflammation of the brain caused by herpes simplex virus infection.

\medterm{Herpes Gestationis} An older name for pemphigoid gestationis, an uncommon autoimmune blistering rash of pregnancy. It usually causes intensely itchy hives followed by blisters around the abdomen and may recur in later pregnancies.

\medterm{Herpes Labialis} Alternate name; see \textit{Cold Sore}.

\medterm{Herpes Simplex} Infection with herpes simplex virus type 1 (HSV-1) or type 2 (HSV-2). HSV-1 most often causes oral herpes or cold sores, while HSV-2 most often causes genital herpes, although either type can infect either site. Many infections cause no symptoms; when symptoms occur, painful blisters or ulcers may appear and recur. The virus remains inactive in nerve cells and can reactivate later. It spreads mainly through skin-to-skin or oral contact, including when no sore is visible, and rarely passes from a parent to a baby during delivery. Rare complications include eye infection, brain infection, or severe infection in a newborn or a person with weakened immunity.

\medterm{Herpes Simplex Encephalitis} A life-threatening brain infection caused most often by herpes simplex virus type 1. It can cause fever, headache, confusion, seizures, behavior changes, or loss of consciousness and needs urgent antiviral treatment.

\medterm{Herpes Simplex Infection} A viral infection caused by herpes simplex virus type 1 (HSV-1) or type 2 (HSV-2). HSV-1 usually causes oral herpes, while HSV-2 usually causes genital herpes, but either type can affect the mouth or genitals. Many infections are symptom-free; painful blisters or ulcers may develop and return because the virus remains inactive in nerve cells and can reactivate. HSV spreads mainly through oral or skin-to-skin contact, including contact with skin that looks normal. Rare complications include infection of the eye or brain and severe infection in newborns or people with weakened immunity.

\synonyms
HSV Infection; Cold Sores and Genital Herpes

\medterm{Herpes Simplex Keratitis} Infection or inflammation of the cornea caused by herpes simplex virus, usually producing eye pain, redness, tearing, light sensitivity, or blurred vision. Recurrent disease can scar the cornea and impair sight.

\medterm{Herpes Simplex Skin Infection} A herpes simplex infection of the skin or moist body linings that causes grouped blisters or shallow sores, often around the mouth or genitals. Widespread infection in someone with eczema can be severe.

\medterm{Herpes Simplex Type 1} HSV-1, a DNA virus commonly causing oral or perioral blisters; it can also cause genital infection and remains latent with possible recurrence.

\medterm{Herpes Simplex Virus} A double-stranded DNA virus belonging to the Herpesviridae family that causes oral and
genital herpes. HSV has two serotypes: HSV-1, which typically causes oral herpes, and
HSV-2, which typically causes genital herpes. The virus establishes latent infections in
sensory ganglia and can reactivate periodically.

\medterm{Herpes Simplex Virus Infection} Infection with herpes simplex virus type 1 or type 2, commonly causing recurrent oral,
genital, skin, eye, or central nervous system disease.

\medterm{Herpes Vegetans} A rare long-lasting herpes simplex infection that causes thick, wart-like or ulcerated growths, usually on genital or anal skin. It occurs mainly in people with weakened immunity and may be difficult to treat.

\medterm{Herpes Viral Culture Of Lesion} A laboratory culture of fluid or tissue from a lesion to detect herpes simplex virus or varicella-zoster virus; molecular testing is generally more sensitive and is often preferred.

\medterm{Herpes Zoster} A painful, usually one-sided blistering rash caused by reactivation of latent varicella-zoster virus in a sensory nerve
ganglion. It typically follows a band-like distribution and may cause persistent nerve
pain after the rash resolves.

\synonyms
Shingles; Zoster; Acute Posterior Ganglionitis

\medterm{Herpes Zoster} A painful, vesicular rash caused by reactivation of latent varicella-zoster virus (VZV,
human herpesvirus 3) in the dorsal root ganglia.

\synonyms
Shingles; Acute Zoster Ganglionitis

\medterm{Herpes Zoster Ophthalmicus} Shingles affecting the forehead, eyelid, or nose near the eye. It can inflame the eye and threaten vision, so a painful blistering rash in this area needs same-day medical assessment and treatment.

\medterm{Herpesviridae} A family of DNA viruses that can remain inactive in the body after infection and reactivate later. Human members include herpes simplex viruses, varicella-zoster virus, cytomegalovirus, Epstein--Barr virus, and human herpesviruses 6, 7, and 8. They can cause skin or mouth lesions, chickenpox, shingles, glandular fever, congenital infection, or disease in people with weakened immunity.

\medterm{Herpesviridae Infection} An infection caused by a herpesvirus; symptoms range from localized lesions to serious disease in newborns or people with weakened immunity.

\medterm{Herpetic Gingivostomatitis} The classic and most common manifestation of primary herpes simplex virus type 1 (HSV-1) infection in infants and young children aged 6 months to 5 years. It is characterized by high fever, irritability, painful submandibular and cervical lymphadenopathy, and widespread friable, inflamed gums with clusters of painful vesicles that rupture into shallow, coalescing ulcers across the buccal mucosa, tongue, and lips, often severely impairing oral intake.

\synonyms
Primary Herpetic Gingivostomatitis; Oral HSV-1 Stomatitis

\medterm{Herpetic Glossitis} Painful inflammation and sores of the tongue caused by herpes simplex virus. It may occur with blisters or ulcers elsewhere in the mouth and can be more severe or persistent in people with weakened immunity.

\medterm{Herpetic Keratitis} An infection of the cornea caused by herpes simplex virus (predominantly HSV-1), recognized globally as the leading infectious cause of corneal blindness. It classically presents with unilateral eye pain, photophobia, foreign-body sensation, tearing, and pathognomonic branching dendritic ulcers with terminal bulbs visualized under fluorescein dye examination; recurrent episodes predispose to deep stromal keratitis, corneal scarring, and permanent vision loss.

\synonyms
Dendritic Keratitis; Herpes Simplex Corneal Infection

\medterm{Herpetic Stomatitis} Inflammation and ulceration of the mouth caused by herpes simplex virus, most commonly primary HSV-1 infection, with painful vesicles or ulcers and sometimes fever or gingivitis.

\medterm{Herpetic Whitlow} A painful herpes simplex infection of a finger, causing redness, swelling, and grouped blisters. It can spread by contact and should not be cut open because this may spread infection.

\medterm{Herpetiform} Resembling herpes lesions, typically describing grouped vesicles or an arrangement of lesions; the term does not by itself establish herpesvirus infection.

\medterm{Hertz} A unit used to measure frequency, equal to one complete cycle or repetition per second. It is used for sound, electrical signals, brain waves, heart signals, and other repeating events.

\medterm{Herxheimer Reaction} A temporary worsening of fever, chills, headache, muscle aches, or other symptoms soon after antibiotics begin for some infections, especially syphilis. It usually settles within a day and is caused by the body’s response to dying bacteria.

\medterm{Hesitation Wounds} Multiple superficial, tentative parallel incised cuts surrounding a deeper primary fatal
cut wound, classically observed in suicides involving sharp instruments (wrists, neck).

\medterm{Hesselbach Triangle} A naturally weaker area of the lower front abdominal wall. Direct groin hernias can push through this area, which lies between a central abdominal muscle, blood vessels, and the groin ligament.

\medterm{Heterochromatin} DNA that is tightly packed inside a cell’s nucleus and is usually less active in making proteins. It helps organize chromosomes and control which genes are available for use.

\medterm{Heterochromia} A difference in iris color between the eyes or within one iris, which may be congenital or acquired from disease, trauma, medication, or inflammation.

\medterm{Heterocyclic Amine} A nitrogen-containing chemical compound formed mainly when meat or fish is cooked at high temperatures; several heterocyclic amines are mutagens and potential dietary carcinogens.

\medterm{Heterogeneity} Variation in composition, structure, or behavior within a tissue, sample, lesion, or population; in tumors, heterogeneity can affect biopsy interpretation and treatment response.

\medterm{Heterogeneous Echotexture} An ultrasound appearance in which different parts of a tissue look unlike one another. It is a finding, not a diagnosis, and may result from inflammation, scarring, fat, bleeding, a tumor, or normal variation.

\medterm{Heterologous} Coming from a different species, source, or type. In medicine, the term may describe a graft, tissue, antibody, hormone, or genetic material obtained from a different organism or cell type.

\medterm{Heterophil Antibody} A heterophil antibody reacts with antigens from unrelated animal species; testing for these antibodies can help diagnose some infections, including infectious mononucleosis.

\medterm{Heterophile Antibody} A heterophile antibody is an antibody that reacts with antigens from unrelated species and can interfere with some laboratory immunoassays.

\medterm{Heterophyes Heterophyes} A small intestinal fluke acquired by eating raw or undercooked fish; it can cause abdominal symptoms and occasionally tissue complications.

\medterm{Heteroplasia} Development of tissue in an abnormal location or of a type not normally found there. The effects depend on the tissue involved and may range from no symptoms to obstruction, pain, or impaired function.

\medterm{Heteroplasmy} The presence of more than one type of mitochondrial DNA in a cell or person, usually a mixture of normal and changed copies. The proportion of changed copies can influence whether disease develops and how severe it is.

\medterm{Heterosexual} Describing a person who experiences sexual or romantic attraction mainly toward people of a different sex or gender. Sexual orientation is distinct from gender identity and is not a disease or disorder.

\medterm{Heterotaxy Syndrome} A birth condition in which organs in the chest and abdomen are arranged differently from usual. Heart defects, abnormal spleen function, and problems with the intestines or other organs may occur.

\medterm{Heterotopic} Located in an abnormal or unusual place; a heterotopic pregnancy is simultaneous implantation outside the uterus and, usually, inside the uterus.

\medterm{Heterotopic Breast Tissue Of The Vulva} Ectopic mammary-like tissue in the vulva that may undergo benign, hormonal, inflammatory, or malignant changes similar to breast tissue; diagnosis is histologic.

\medterm{Heterotopic Pregnancy} A pregnancy in which one embryo implants inside the uterus while another implants outside it, usually in a fallopian tube. It is uncommon but more likely after fertility treatment; an intrauterine pregnancy does not rule out an ectopic pregnancy.

\medterm{Heterotransplant} A transplant of cells, tissue, or an organ between members of different species. Because the immune difference is substantial, rejection is a major concern and long-term use requires specialized treatment.

\medterm{Heterozygote} A heterozygote has two different alleles at a genetic locus, one inherited from each parent; phenotype depends on dominance, the variant's effect, and the surrounding genetic and environmental context.

\medterm{Heterozygous} Having two different alleles of a particular gene at a genetic locus. The observable effect depends on the gene and the relationship between the alleles.

\medterm{Hexachlorobenzene} A persistent organochlorine compound that can cause porphyria, liver injury, neurologic effects, and other toxicity after significant exposure.

\medterm{Hexachlorophane} An antibacterial chemical formerly used in some skin cleansers. Repeated or extensive exposure can damage the nervous system, especially in infants, so its medical use is highly restricted.

\medterm{Hexachlorophene} A topical antiseptic phenolic compound whose excessive or prolonged use can cause neurotoxicity and skin injury; many uses are restricted.

\medterm{Hexadactyly} The presence of six fingers or six toes on a hand or foot, also called polydactyly.

\medterm{Hexadecadienoic Acid} A long-chain fatty acid containing 16 carbon atoms and two double bonds, occurring in lipid metabolism and studied as a biochemical or signaling molecule.

\medterm{Hexadecatetraenoic Acid} A 16-carbon polyunsaturated fatty acid containing four double bonds, present in lipid pathways and investigated for effects on membranes and signaling.

\medterm{Hexamethonium} An older medicine that blocks nerve signals at autonomic ganglia and can lower blood pressure. It is rarely used because it also causes blurred vision, dry mouth, constipation, urinary retention, and other widespread effects.

\medterm{Hexamine} A chemical used in industry and, in some settings, as methenamine, a urinary antiseptic. It releases formaldehyde in acidic urine and is unsuitable for some kidney or liver conditions.

\medterm{Hexane} A solvent found in some industrial products and glue. Long-term inhalation can damage peripheral nerves, causing numbness, weakness, and difficulty walking; swallowing or heavy exposure can also injure the lungs and brain.

\medterm{Hexanoic Acid} A six-carbon saturated fatty acid, also called caproic acid, found in some fats and odors and used in chemical synthesis; concentrated exposure irritates skin and mucosa.

\medterm{Hexokinase} An enzyme that begins the use of glucose by adding a phosphate group to it inside cells. This traps glucose in the cell and allows it to enter energy-producing or storage pathways. Most tissues use hexokinase; the liver also uses a related enzyme called glucokinase.

\medterm{Hexose} A six-carbon monosaccharide, such as glucose, fructose, or galactose, that serves as an energy substrate or metabolic intermediate.

\medterm{Hexylresorcinol} A phenolic antiseptic and local anesthetic used in some throat lozenges and topical products; excessive use or allergy can cause adverse effects.

\medterm{Hg} The chemical symbol for mercury; in measurements such as mmHg, it refers to millimeters of mercury, a unit of pressure.

\medterm{Hiatal} Relating to a hiatus, an opening or passage in an anatomical structure; a hiatal hernia occurs when abdominal tissue passes through the diaphragm's esophageal hiatus.

\medterm{Hiatal Hernia} Protrusion of part of the stomach through the esophageal opening in the diaphragm into the
chest. It may be sliding or paraesophageal and can be associated with reflux or other symptoms.

\medterm{Hiatus} A natural opening or passage in a body structure. For example, the diaphragm has an opening through which the esophagus passes; tissue can sometimes bulge through such an opening.

\medterm{Hiatus Hernia} Another term for hiatal hernia: protrusion of part of the stomach through the esophageal hiatus
of the diaphragm into the chest. It may be sliding or paraesophageal.

\medterm{Hib} Haemophilus influenzae type b, a bacterium that can cause serious infections such as meningitis, pneumonia, bloodstream infection, and swelling of the throat. Routine vaccination greatly reduces these infections in children.

\medterm{Hibernoma} A rare benign tumor of brown fat, usually a slow-growing painless soft-tissue mass that may feel warm because of its high metabolic and vascular activity.

\medterm{Hiccough} Alternate spelling; see \textit{Hiccups}.

\medterm{Hiccough} Alternate spelling; see \textit{Hiccups}.

\medterm{Hiccup} Singular form; see \textit{Hiccups}.

\medterm{Hiccups} Involuntary, spasmodic contractions of the diaphragm followed by sudden closure of the glottis,
producing a characteristic sound. Persistent or intractable hiccups may reflect
irritation of the diaphragm or phrenic and vagal pathways, metabolic disease, medication
exposure, or central nervous system disease.

\medterm{Hidradenitis} A chronic inflammatory skin disorder causing repeated painful lumps, abscesses, draining tunnels, and scars, usually in the armpits, groin, buttocks, or under the breasts.

\medterm{Hidradenitis Suppurativa} A long-lasting skin disease causing repeated painful lumps, pockets of pus, draining tunnels, and scars, usually in the armpits, groin, buttocks, or under the breasts. It is not caused by poor hygiene and is not contagious.

\medterm{Hidradenocarcinoma} A rare malignant sweat-gland tumor that may arise de novo or from a hidradenoma, presenting as a firm nodule and capable of local recurrence and metastasis.

\medterm{Hidradenoma} A usually harmless tumor arising from sweat-gland cells, typically appearing as a single slow-growing skin lump. It may be flesh-colored, red, cyst-like, or ulcerated, so examination or removal may be needed to distinguish it from other tumors.

\medterm{Hidradenoma Papilliferum} A benign papillary tumor of apocrine or mammary-like glands, usually presenting as a solitary flesh-colored or red nodule in the vulvar or anogenital region.

\medterm{Hidrocystoma} A benign cyst of a sweat gland near the eyelid or skin, appearing as a small translucent or bluish papule that may enlarge with heat or sweating.

\medterm{Hidrosis} Sweating or perspiration; abnormal hidrosis may mean excessive sweating, reduced sweating, or sweating in an unusual pattern.

\medterm{High Altitude} Elevation high enough that reduced atmospheric pressure lowers available oxygen, commonly beginning around 2,500 meters and increasing the risk of altitude illness.

\medterm{High Altitude Illness} Health problems caused by traveling to high elevation before the body adapts to lower oxygen levels. Symptoms range from headache and nausea to dangerous brain or lung swelling requiring immediate descent.

\medterm{High Arch} An unusually high curve along the inside of the foot. It may run in families or result from a nerve or muscle disorder and can cause pain, calluses, ankle instability, or difficulty finding comfortable shoes.

\synonyms
Pes Cavus; High-Arched Foot; Cavus Foot

\medterm{High Blood Cholesterol Levels} Too much cholesterol in the blood, especially LDL cholesterol, which can build up in artery walls and increase the risk of heart attack and stroke. It usually causes no symptoms and is found with a blood test.

\medterm{High Blood Potassium} Elevated potassium concentration in the bloodstream; see \textit{Hyperkalemia}.

\medterm{High Blood Pressure} Persistently elevated arterial blood pressure. Diagnosis generally requires properly obtained readings on more than one occasion or other appropriate confirmation. It often causes no symptoms but increases the risk of cardiovascular, cerebrovascular, kidney, and eye disease.

\medterm{High Blood Pressure And Eye Disease} Long-standing high blood pressure can damage retinal blood vessels and the optic nerve, causing hypertensive retinopathy or vision loss.

\medterm{High Blood Pressure Medications} Medicines used to lower arterial blood pressure, including diuretics, ACE inhibitors, ARBs, calcium-channel blockers, beta blockers, and other agents selected according to comorbidities and risk.

\medterm{High Blood Protein} An abnormally increased concentration of proteins in the blood, which may result from dehydration, inflammation, infection, or a plasma-cell disorder.

\medterm{High Blood Sugar} Hyperglycemia, an abnormally elevated concentration of glucose in the blood; persistent hyperglycemia is commonly associated with diabetes but can have other causes.

\medterm{High Cholesterol} An abnormally high level of cholesterol in the blood, especially low-density lipoprotein cholesterol. It often causes no symptoms but can contribute to plaque buildup in the arteries and cardiovascular disease.

\medterm{High Hemoglobin Count} A hemoglobin level above the expected range, which may reflect dehydration, chronic low oxygen, smoking, medication effects, or a blood disorder.

\medterm{High Output Failure} High-output heart failure, in which cardiac output is elevated or normal but inadequate for the body's unusually high metabolic or circulatory demand. Causes include severe anemia, thyrotoxicosis, beriberi, and large arteriovenous fistulas; treatment targets the cause and may include heart-failure therapy.

\medterm{High Potassium Level} Hyperkalemia is an elevated blood potassium concentration that can cause dangerous cardiac conduction abnormalities; confirmation, medication review, kidney assessment, ECG, and urgent treatment depend on level and symptoms.

\medterm{High Protein Diet} An eating pattern providing more protein than a person's usual or estimated requirement; suitability depends on nutritional needs, kidney and liver function, and the overall diet.

\medterm{High Red Blood Cell Count} An increased number or concentration of red blood cells in the blood, which may thicken the blood and result from dehydration, low oxygen, smoking, medications, or a bone marrow disorder.

\synonyms
Erythrocytosis; Polycythemia

\medterm{High Uric Acid Level} An increased concentration of uric acid in the blood, which may contribute to gout or kidney stones and can result from reduced kidney clearance, diet, medicines, or metabolic disease.

\medterm{High White Blood Cell Count} An increased number of white blood cells, often associated with infection, inflammation, stress, medicines, or a blood or bone marrow disorder.

\medterm{High White Blood Cell Count In Pregnancy} A mild rise in white blood cells can be normal during pregnancy, labor, or stress. Marked or persistent elevation with fever, pain, urinary or respiratory symptoms, or abnormal differential counts may indicate infection, inflammation, medication effect, or another disorder.

\medterm{High-Altitude Pulmonary Edema} Life-threatening fluid buildup in the lungs after rapid travel to high altitude. It causes breathlessness, cough, weakness, and sometimes pink or frothy sputum; immediate descent and emergency treatment are needed.

\medterm{High-Density Lipoprotein} A particle that carries cholesterol through the blood. Often called “HDL cholesterol,” its level is one part of heart-risk assessment and does not by itself determine whether someone will have heart disease.

\medterm{High-Dependency Unit} A hospital area for people who need closer observation and more treatment than a general ward provides but do not need the full support of an intensive-care unit. Staff monitor vital signs frequently and provide specialized care.

\medterm{High-Fiber Foods} Foods rich in soluble or insoluble dietary fiber, including legumes, whole grains, vegetables, fruit, nuts, and seeds, support bowel regularity and may improve cholesterol and glucose control when increased with adequate fluid.

\medterm{High-Flow Nasal Cannula} A device that delivers warmed, moist oxygen through small tubes in the nose at a high flow rate. It can improve breathing and oxygen levels in people with serious lung illness and is also used after removal of a breathing tube.

\medterm{High-Fructose Corn Syrup} A liquid sweetener containing varying proportions of glucose and fructose, used in many processed foods and drinks. Like other added sugars, high intake increases total energy exposure and is associated with dental caries and cardiometabolic risk; it is not inherently unique in effect compared with equivalent added sugars.

\medterm{High-Grade Endometrial Stromal Sarcoma} An aggressive cancer arising from supporting tissue in the uterus. It can grow into the uterine wall, blood vessels, or distant organs and requires specialist treatment based on its extent and cell features.

\medterm{High-Grade Serous Ovarian Carcinoma} An aggressive cancer involving the ovary, fallopian tube, or nearby lining of the abdomen. It often spreads before causing clear symptoms and may cause bloating, abdominal pain, or changes in bowel or bladder habits.

\medterm{High-Grade Surface Osteosarcoma} An aggressive bone cancer that starts on the outer surface of a bone. It may cause pain or swelling and can spread to the lungs or other sites, so treatment usually requires specialist cancer care.

\medterm{High-Intensity Focused Ultrasound} A treatment that concentrates sound-wave energy to heat and destroy a selected area of tissue while limiting damage nearby. Uses include some tumors and other conditions, depending on the device and treatment site.

\medterm{High-Level Disinfection} A cleaning process that kills nearly all microorganisms but may not destroy large numbers of bacterial spores. It is used for some reusable medical equipment that touches moist body surfaces and cannot be heat-sterilized.

\medterm{Highly Pathogenic Avian Influenza} Highly pathogenic influenza A infection in birds that can cause severe disease and mortality in poultry; selected strains can rarely infect humans after animal exposure.

\medterm{High-Pitched Cry} An unusually shrill infant cry that can occur with pain, neurologic dysfunction, drug withdrawal, or serious illness and warrants assessment when persistent or accompanied by abnormal behavior or feeding.

\medterm{High-Resolution Computed Tomography} A CT scan that uses very thin images and special settings to show fine detail, especially in the lungs. It can help detect scarring, inflammation, widened airways, and small-airway disease, but it uses X-rays.

\synonyms
HRCT; High-Resolution Chest CT

\medterm{High-Sensitivity Cardiac Troponin} A blood test that detects very small amounts of proteins released when heart muscle is injured. A raised result can occur with a heart attack or many other illnesses, so clinicians interpret repeated results with symptoms, ECG findings, and other tests.

\synonyms
hs-cTn; High-Sensitivity Troponin Assay

\medterm{Hilar} Relating to a hilum, the gateway where vessels, nerves, ducts, or bronchi enter or leave an organ, especially the lung or kidney.

\medterm{Hilar Cholangiocarcinoma} A cancer at the place where the liver’s main bile ducts join. It can block bile flow and cause yellow skin or eyes, itching, dark urine, pale stools, weight loss, or repeated bile-duct infections.

\medterm{Hill's Sign} Hill's sign is an older physical finding described in some cases of aortic regurgitation; it is not sufficient by itself to diagnose the condition.

\medterm{Hill-Sachs Lesion} A dent or compression fracture in the back of the upper-arm bone caused when the shoulder dislocates forward, sometimes contributing to repeated dislocations.

\medterm{Hilum} A depression or opening where blood vessels, nerves, ducts, or other structures enter or leave an organ.

\medterm{Hilus} The hilus is an indentation or entry point on an organ where vessels, nerves, ducts, or airways enter or leave.

\medterm{Hindbrain} The hindbrain is the lower part of the developing brain, giving rise to the medulla, pons, and cerebellum and helping control vital functions and coordination.

\medterm{Hindgut} The hindgut is the embryonic part of the digestive tract that forms much of the distal colon, rectum, and upper anal canal.

\medterm{Hip} The joint and surrounding region where the thigh bone meets the pelvis, allowing movement and bearing body weight.

\medterm{Hip Arthroscopy} A minimally invasive procedure that uses a camera and instruments inserted through small incisions to diagnose and treat selected disorders inside or around the hip joint.

\medterm{Hip Bursitis} Hip bursitis is inflammation of a bursa near the hip, commonly causing pain over the outer hip that worsens with lying on that side, stairs, or prolonged walking. Activity changes, physical therapy, medicines, or injection may help after other causes are considered.

\synonyms
Trochanteric Bursitis; Greater Trochanteric Pain Syndrome

\medterm{Hip Circumference} The measurement around the widest part of the hips or buttocks, used in anthropometry and sometimes combined with waist circumference to assess body-fat distribution.

\medterm{Hip Dislocation} The ball at the top of the thighbone comes out of the hip socket. It is usually caused by major injury and causes severe pain and inability to move the leg; nearby nerves and the blood supply to the bone may also be injured.

\medterm{Hip Disorders} The hip joints are ball-and-socket joints connecting the femur to the pelvis, providing
stability and a wide range of motion for walking, running, and weight-bearing
activities. Common hip conditions include osteoarthritis, fractures, bursitis, and
labral tears.

\medterm{Hip Dysplasia} A problem present from birth or developing early in life in which the hip socket and upper thighbone do not fit together securely. The hip may be unstable or dislocate. Early examination and ultrasound or X-ray help guide treatment and protect hip development.

\synonyms
Developmental Dysplasia of the Hip (DDH); Congenital Hip Dislocation

\medterm{Hip Fracture} A break in the upper thighbone, commonly near the femoral neck or just below it. It often follows a fall in someone with fragile bones and causes hip or groin pain, inability to bear weight, and sometimes a shortened or outward-turned leg.

\medterm{Hip Fracture In Older Adults} A fracture of the upper thighbone in an older person, commonly after a fall or minor force when bone strength is reduced. It can cause pain, loss of mobility, delirium, blood clots, and loss of independence.

\synonyms
Femoral Neck Fracture; Geriatric Hip Fracture

\medterm{Hip Fracture Surgery} An operation to stabilize or replace a fractured proximal femur, using internal fixation, hemiarthroplasty, or total hip arthroplasty according to fracture pattern and patient factors.

\medterm{Hip Impingement} Abnormal contact between parts of the hip joint during movement, often due to cam or pincer morphology. It may cause groin pain, stiffness, and reduced motion and can contribute to labral or cartilage damage.

\medterm{Hip Injury} Damage involving the hip joint, surrounding muscles, tendons, ligaments, or bones; causes range from strain and bursitis to fracture and dislocation, with inability to bear weight requiring prompt assessment.

\medterm{Hip Joint} The ball-and-socket joint where the top of the thighbone fits into the pelvis. It supports body weight and allows the leg to move in many directions; strong ligaments and a cartilage rim help keep it stable.

\medterm{Hip Joint Injection} Injection of local anesthetic, corticosteroid, or another prescribed agent into or around the hip joint for diagnosis or temporary relief of pain and inflammation.

\medterm{Hip Joint Replacement} Arthroplasty replaces damaged parts of the hip with prosthetic components to relieve pain and improve mobility, commonly for advanced osteoarthritis, fracture, or other destructive joint disease.

\medterm{Hip Labral Tear} A tear in the ring of cartilage surrounding the hip socket. It may cause groin pain, clicking, catching, stiffness, or reduced movement and can result from injury, repeated stress, or abnormal hip shape.

\medterm{Hip Pain} Pain in or around the hip caused by joint disease, tendon or muscle injury, bursitis, fracture, nerve problems, or pain referred from the back.

\medterm{Hip Prosthesis} An artificial hip joint used to replace a damaged hip. It usually has a stem placed in the thighbone, a ball, and a socket; parts may be metal, ceramic, or durable plastic and may be fixed with or without bone cement.

\medterm{Hip Replacement} A surgical procedure in which the diseased hip joint is replaced with a prosthetic
implant (arthroplasty). Total hip replacement (THR) replaces both the femoral head and
the acetabulum.

\medterm{Hip-Fracture Rehabilitation} Recovery care after a broken hip that helps restore walking, strength, balance, and independence. It may include early movement, pain control, exercises, nutrition, prevention of blood clots, delirium prevention, and treatment to reduce future fracture risk.

\medterm{Hippocampus} A pair of structures deep in the brain that help form and retrieve new memories and support learning and navigation. Damage to them can cause difficulty forming new memories.

\medterm{Hippocrates} An ancient Greek physician traditionally regarded as a major founder of Western clinical medicine; the Hippocratic corpus emphasizes observation, prognosis, and professional conduct.

\medterm{Hippocratic Oath} A traditional promise that doctors make to practice medicine ethically, help patients, protect their privacy, and avoid harm. Modern versions differ by country and institution.

\medterm{Hippotherapy} A treatment in which a trained therapist uses the movement of a horse to support physical, occupational, or speech therapy. It may help some people improve balance, posture, muscle control, coordination, or communication, but it is tailored to individual goals and safety needs.

\medterm{Hippus} Hippus is rhythmic, involuntary widening and narrowing of the pupil, which may occur normally or with neurologic or eye disease.

\medterm{Hirayama Disease} A condition affecting the lower neck spinal cord, usually in teenage boys and young men. It causes slowly worsening, often one-sided weakness and wasting of the hand and forearm muscles. Neck bending can press the spinal cord and reduce its blood supply. Symptoms commonly stop progressing after several years. A neck MRI taken during bending helps confirm the diagnosis; early neck protection may limit damage.

\synonyms
Monomelic Amyotrophy; Juvenile Cervical Flexion Myelopathy

\medterm{Hirschsprung Disease} A condition present from birth in which a section of the large intestine lacks the nerve cells needed to relax and move stool. Newborns may fail to pass stool, develop a swollen abdomen, or have severe constipation.

\synonyms
Congenital Aganglionic Megacolon; Colonic Aganglionosis

\medterm{Hirschsprung's Disease} A birth condition in which part of the large intestine lacks the nerve cells needed to relax and move stool. Newborns may fail to pass stool, develop a swollen abdomen, vomiting, or severe constipation.

\medterm{Hirsutism} Excess coarse hair growth in a woman or person with ovaries in areas where men commonly grow hair, such as the face, chest, or abdomen. It may result from increased androgen hormones, polycystic ovary syndrome, medicines, or other conditions.

\synonyms
Excessive Terminal Hair Growth; Hypertrichosis in Women

\medterm{Hirudin} A substance first found in medicinal leech saliva that prevents blood from clotting by blocking an important clotting protein. Related medicines have been used or studied for selected clotting conditions.

\medterm{Hirudins} A family of thrombin-inhibiting anticoagulant peptides, including natural and recombinant agents, used or studied to prevent or treat clotting in selected clinical situations.

\medterm{Hirudo Medicinalis} The medicinal leech, used in selected medical procedures to relieve venous congestion after reconstructive surgery.

\medterm{His} Usually referring to the bundle of His, the atrioventricular conduction bundle that carries electrical impulses from the AV node into the right and left bundle branches.

\medterm{His Bundle Electrography} An invasive electrophysiologic recording of electrical activity in the atrioventricular His bundle, used to assess conduction-system disease and localize atrioventricular block.

\medterm{His Disease} Trench fever, a louse-borne infection caused by \textit{Bartonella quintana}, classically producing recurrent or five-day fever, headache, and severe shin or leg pain.

\medterm{Histamine} A chemical messenger released during allergic and inflammatory reactions. It can cause itching, redness, swelling, mucus production, narrowing of the airways, and increased stomach acid.

\medterm{Histamine Degradation} The body’s breakdown and removal of histamine after it has acted. This limits effects such as itching, swelling, airway narrowing, and increased stomach acid.

\medterm{Histamine Receptor} A protein on or in a cell that responds to histamine. Different receptor types help control allergic symptoms, stomach-acid production, nerve signaling, and immune responses.

\medterm{Histamine Release} Discharge of histamine from mast cells or basophils during allergic, immune, or nonimmune activation, causing effects such as itching, vasodilation, swelling, mucus production, and bronchoconstriction.

\medterm{Histidine} Histidine is an amino acid used in protein synthesis and a precursor of histamine; it is essential during growth and can be metabolized to urocanic acid and other products.

\medterm{Histidine Decarboxylase} An enzyme that converts the amino acid histidine into histamine. It is found in human cells and some microorganisms; histamine influences allergy, stomach acid, blood-vessel responses, and immunity.

\medterm{Histidinemia} An inherited metabolic disorder caused by reduced histidase activity, leading to elevated histidine in blood and urine; many affected people have few or no clinically important symptoms.

\medterm{Histiocyte} An immune cell found in tissues that can engulf germs, damaged cells, or foreign material and help alert other immune cells.

\medterm{Histiocytic Sarcoma} A rare aggressive cancer of cells that normally help process immune signals and foreign material. It may arise in lymph nodes, the intestine, skin, or other organs and requires specialist pathology and cancer care.

\medterm{Histiocytoma} A tumor or tumor-like proliferation of histiocytes; the term may refer to a benign fibrous skin lesion or, less commonly, a malignant histiocytic neoplasm depending on classification.

\medterm{Histiocytosis} A group of uncommon disorders in which certain immune cells build up abnormally in the skin, bones, organs, or blood. Effects range from isolated skin or bone problems to serious disease affecting several organs.

\medterm{Histochemical} Histochemical means using chemical stains or reactions in tissue sections to identify particular substances or cell components.

\medterm{Histocompatibility} The degree to which donor and recipient tissues are immunologically alike. Matching tissue markers, especially HLA markers, can lower rejection risk, but antibody testing, organ type, and the recipient's health also influence transplant success.

\medterm{Histocompatibility Antigen Test} A test identifying human leukocyte antigen types, which influence tissue compatibility for transplantation and are relevant to some immune-mediated and genetic associations.

\medterm{Histocompatibility Testing} Laboratory typing of HLA and other immune markers to assess donor--recipient compatibility before transplantation and to investigate selected immune or transfusion-related questions.

\medterm{Histogenesis} The developmental process by which unspecialized cells differentiate and organize into mature tissues with characteristic structures and functions.

\medterm{Histological Techniques} Laboratory methods for preparing, staining, sectioning, and examining tissue under a microscope to reveal cellular architecture, deposits, organisms, or disease-related changes.

\medterm{Histology} The study of cells and tissues under a microscope. A laboratory prepares a thin tissue section, applies stains that show its structures, and examines it for normal features or changes caused by disease.

\medterm{Histology Cassette} A labeled perforated container that holds a tissue specimen during processing, allowing fixation, dehydration, paraffin infiltration, and later embedding for microscopic sectioning.

\medterm{Histone} A protein around which DNA winds so it can fit inside the cell nucleus. Chemical changes to histones help control which genes are turned on or off.

\medterm{Histone Acetylation} Addition of acetyl groups to histone proteins that usually loosens DNA packaging and can increase access of transcriptional machinery to genes.

\medterm{Histone Acetyltransferase} An enzyme that transfers acetyl groups to histone proteins, generally promoting a more open chromatin state and regulating gene transcription.

\medterm{Histone Deacetylases} Enzymes that remove acetyl groups from histones and other proteins, influencing chromatin structure, gene expression, cell differentiation, and disease biology.

\medterm{Histone Deacetylation} Removal of acetyl groups from histones, often tightening chromatin and reducing transcription; histone deacetylases are targets of some anticancer drugs.

\medterm{Histone Modification} A chemical change to a histone protein that alters how tightly DNA is packed and can change which genes a cell uses. These changes help regulate cell growth, development, and response to disease.

\medterm{Histopathologic Grade} A microscopic estimate of how abnormal and aggressive tumor cells appear compared with their tissue of origin; grade contributes to prognosis and treatment planning alongside stage.

\medterm{Histopathology} The diagnosis of disease by examining tissue under a microscope. A pathologist studies a biopsy or surgical specimen for abnormal cells, inflammation, infection, scarring, or cancer.

\medterm{Histoplasma} A genus of fungi that includes Histoplasma capsulatum, which can cause histoplasmosis after inhalation of environmental spores.

\medterm{Histoplasma Capsulatum} A fungus found in soil contaminated with bird or bat droppings that can cause histoplasmosis when
its spores are inhaled.

\medterm{Histoplasma Complement Fixation} A blood test that detects complement-fixing antibodies to Histoplasma and can support evaluation of histoplasmosis, although sensitivity, timing, cross-reactions, and immune status affect interpretation.

\medterm{Histoplasmosis} A fungal infection acquired by breathing in spores from soil contaminated with bird or bat droppings. It often affects the lungs and may be mild, but can spread to other organs, especially in people with weakened immunity.

\medterm{Histrionic Personality Disorder} A long-term pattern of intense emotional expression, strong need for attention, and behavior that may seem dramatic, suggestible, or uncomfortable when not the focus. Diagnosis requires a persistent pattern that causes significant problems and is made by a qualified mental-health professional.

\medterm{Hitchhiker's Thumb} A descriptive term for an unusually flexible thumb that can bend backward at the distal joint. It is usually a harmless inherited variation unless associated with pain, instability, or a broader connective-tissue disorder.

\medterm{HIV} Abbreviation for human immunodeficiency virus; see \textit{Human immunodeficiency virus}.

\medterm{Hiv Antibodies} Proteins made by the immune system after HIV infection and detected by blood or oral-fluid tests. A positive screening result needs confirmatory testing, and a negative result may be unreliable soon after exposure.

\medterm{HIV Infection} Infection with human immunodeficiency virus, which progressively damages immune function if untreated.

\medterm{Hiv Seropositivity} The presence of detectable HIV antibodies or other blood-test evidence of infection. A reactive screening test is not final until confirmatory testing is completed.

\medterm{HIV Test} A blood or oral-fluid test that looks for HIV infection. Laboratory tests may detect viral material or antibodies; a reactive screening result needs confirmatory testing. Tests can be negative early in infection, so repeat testing may be needed after a recent exposure.

\medterm{HIV Transmission And Prevention} The spread of HIV through infected blood, semen, vaginal fluids, rectal fluids, or breast milk, prevented by testing, treatment that suppresses viral load, condoms, sterile equipment, and pre- or post-exposure prophylaxis.

\medterm{HIV/AIDS} Common combined term referring to HIV infection and its most advanced stage, AIDS; see \textit{HIV infection} and \textit{AIDS}.

\medterm{HIV/AIDS Immunopathogenesis} HIV enters and gradually destroys or disables important immune cells. Without treatment, immune defenses become weaker and serious infections or cancers may develop. Medicines stop most viral multiplication and allow immune recovery.

\medterm{HIV/AIDS In Pregnant Women And Infants} HIV care in pregnancy and infancy uses maternal antiretroviral therapy, viral-load monitoring, delivery and feeding planning, infant prophylaxis, and early infant diagnostic testing to reduce vertical transmission.

\medterm{HIV/AIDS Nutrition} Nutritional care for people living with HIV or AIDS. Infection, poor appetite, diarrhea, medicine side effects, or other illnesses may cause weight loss or nutrient shortages; care may include food support, supplements, treatment of infections, and review of medicines.

\medterm{HIV-Associated Rheumatic Disease} Joint, muscle, bone, or connective-tissue problems associated with HIV infection or its immune effects. They may include inflammatory arthritis, reactive arthritis, muscle disease, osteoporosis, or symptoms related to infection or treatment.
Antiretroviral therapy suppresses viral replication and reduces transmission and complications.

\medterm{Hives} The common term for urticaria, describing raised, intensely itchy welts (wheals) on the skin. They happen when immune cells in the skin release histamine, causing fluid to leak from tiny blood vessels. Individual welts typically appear quickly, change shape, and fade away within 24 hours without leaving a scar. Hives lasting under six weeks are acute, often triggered by viral infections, foods, insect stings, or medications. Hives lasting longer than six weeks are chronic and often occur without a known external trigger.

\synonyms
Urticaria; Urticarial Weals

\medterm{Hives And Angioedema} Hives are raised itchy welts; angioedema is deeper swelling of the skin or mucous membranes, often caused by allergy, medicines, or other immune reactions.

\medterm{HLA} A group of cell-surface proteins helping the immune system recognize the body's own cells, reject foreign tissue, and influence transplant compatibility.

\synonyms
Human Leukocyte Antigen System; Major Histocompatibility Complex (MHC)

\medterm{Hla Antigen} A marker on the surface of cells that helps immune cells recognize what belongs to the body. HLA markers are important for matching donors and recipients in transplantation and are linked with risk of some immune diseases.

\medterm{HLA System} The HLA system is a group of cell-surface proteins that present antigens to immune cells and influence transplant matching and autoimmune risk.

\medterm{HLA-B27} An inherited immune-system marker associated with increased risk of some inflammatory conditions affecting the spine, joints, eyes, or bowel. A positive test does not prove that a person has one of these diseases.

\medterm{HLA-B27 Antigen} A cell-surface immune marker associated with increased risk of ankylosing spondylitis and related spondyloarthritis, reactive arthritis, and uveitis; a positive result is not diagnostic by itself.

\medterm{Hmgb Proteins} Proteins that help organize genetic material inside cells. When released from damaged cells, some can also signal tissue injury and promote inflammation.

\medterm{HMG-CoA Reductase} An enzyme the liver uses to make cholesterol. Statin medicines partly lower blood cholesterol by blocking this enzyme.

\medterm{Hmgn Proteins} Proteins that help control how genetic material is arranged and used inside cells. They can influence cell growth, specialization, and repair of genetic material.

\medterm{Hoarding} Alternate name; see \textit{Hoarding Disorder}.

\medterm{Hoarding Disorder} A mental-health condition involving persistent difficulty discarding possessions, even when they have little value. Saving items causes distress and can fill living areas so that normal activities or safety are impaired.

\medterm{Hoarseness} A change in the voice that makes it rough, weak, breathy, or strained. Common causes include a viral infection, voice overuse, smoking, reflux, or a growth on the vocal folds. Hoarseness lasting more than a few weeks needs assessment.

\medterm{Hodgkin Lymphoma} A cancer of a type of immune cell, usually beginning in lymph nodes. It often causes painless swelling of the nodes, fever, night sweats, itching, or weight loss and is diagnosed by examining a tissue sample.

\medterm{Hodgkin Lymphoma In Children} A malignant lymphoma arising from lymphoid tissue in children, characterized by Hodgkin and Reed-Sternberg cells and treated according to subtype, stage, and risk.

\medterm{Hodgkin's And Non-Hodgkin's Lymphoma Comparison} Hodgkin lymphoma is defined by characteristic Reed-Sternberg cells and commonly spreads in an orderly pattern, while non-Hodgkin lymphomas comprise many B-cell, T-cell, and other lymphoid cancers. Biopsy, immunophenotyping, and staging determine treatment and prognosis.

\medterm{Hodgkin's Disease} Alternate name; see \textit{Hodgkin Lymphoma}.

\medterm{Hodgkin's Disease} Historical term; see \textit{Hodgkin lymphoma}.

\medterm{Hodgkin's Lymphoma} A B-cell lymphoma characterized by the presence of Reed-Sternberg cells (large,
binucleate or multinucleate cells with owl-eye nucleoli) in a background of reactive
inflammatory cells. See Hodgkin Lymphoma.

\medterm{Hoffmann Sign} A reflex test where flicking the middle fingernail makes the thumb and index finger quickly twitch inward, pointing to irritation or pressure on the spinal cord in the neck.

\medterm{Hoffmann Syndrome} An uncommon form of adult muscle weakness caused by severe, untreated underactive thyroid disease. It may cause stiff or painful muscles, cramps, enlarged-looking muscles, slow reflex relaxation, and high muscle-enzyme levels; treating the thyroid problem usually improves it.

\synonyms
Hypothyroid Myopathy with Pseudohypertrophy

\medterm{Hofmann's Syndrome} A rare form of adult hypothyroid muscle disease causing muscle enlargement, stiffness, cramps, weakness, and slowed reflexes. Treating the low thyroid hormone level usually improves the muscle problems.

\medterm{Holistic Medicine} An approach that considers physical, psychological, social, behavioral, and sometimes spiritual factors in health and illness. It may support patient-centered care, but individual therapies should be assessed for evidence, safety, and compatibility with effective conventional treatment.

\medterm{Holmium} A rare-earth metal used in some medical lasers and radioactive tracers. Holmium laser energy can cut or vaporize tissue, especially in urologic procedures; radioactive holmium requires strict handling and safety controls.

\medterm{Holmium Laser Enucleation Of The Prostate} A procedure performed through the urine passage that uses a laser to separate and remove excess prostate tissue blocking urine flow. The tissue is removed from the bladder and examined; temporary bleeding, burning, leakage, infection, or sexual and urinary changes can occur.

\synonyms
HoLEP; Holmium Prostate Enucleation

\medterm{Holoprosencephaly} A birth difference in which the developing brain does not divide fully into two cerebral hemispheres. Severity varies; it may affect facial development, movement, feeding, seizures, and intellectual development.

\medterm{Holorachischisis} A severe birth defect in which a long section of the spine and spinal cord fails to close during development. It usually causes major weakness, loss of sensation, and other serious nervous-system problems.

\medterm{Holosystolic Murmur Differential} A heart murmur heard throughout the period when the heart's lower chambers contract. Common causes include leakage of the mitral or tricuspid valve and a hole between the lower chambers; its sound, location, and other findings help identify the cause.

\medterm{Holter Electrocardiography} Continuous ambulatory ECG recording, usually for 24 hours or longer, used to detect intermittent arrhythmias, correlate symptoms with rhythm, and assess heart-rate patterns.

\medterm{Holter Monitor} A small portable recorder that continuously tracks the heart’s electrical rhythm, usually for a day or longer. It can detect intermittent abnormal rhythms and compare them with symptoms such as palpitations, dizziness, or fainting.

\medterm{Holter Monitor} A portable ambulatory electrocardiographic recorder that continuously monitors heart rhythm, usually for 24 hours or longer, to detect intermittent arrhythmias and correlate them with symptoms.

\medterm{Holter Monitoring} Portable, continuous electrocardiographic recording over an extended period, commonly 24 hours or longer, to detect intermittent rhythm abnormalities and correlate them with symptoms.

\medterm{Holt-Oram Syndrome} An inherited condition causing abnormalities of the bones in one or both arms or hands and heart defects. The heart problem may be a hole between chambers or an abnormal rhythm; the severity varies widely.

\medterm{Homans Sign} Calf discomfort when the foot is gently bent upward toward the shin, historically checked as a clue for a blood clot in the leg.

\medterm{Homeobox Gene} A gene that helps control how an embryo’s body is laid out and how organs and tissues develop. Changes in some homeobox genes can cause birth differences or disease.

\medterm{Homeodomain Proteins} Transcription factors containing a homeodomain that bind regulatory DNA sequences and direct developmental patterning and cell differentiation.

\medterm{Homeopathy} An alternative-health system that uses highly diluted preparations selected according to “like cures like”; reliable evidence has not shown homeopathic products to treat serious disease beyond placebo effects.

\medterm{Homeostasis} The body’s ability to keep internal conditions within a safe range despite changes outside it. It regulates temperature, blood sugar, blood pressure, fluid, salts, and acidity through coordinated feedback systems.

\medterm{Homeotherm} An organism that maintains a relatively stable internal body temperature despite changes in environmental temperature, largely through physiological regulation.

\medterm{Homicide} The killing of one person by another, classified legally according to intent and circumstances; medical assessment may address injury patterns, cause of death, toxicology, and safeguarding.

\medterm{Homo Sapiens} The scientific name for modern humans, the living human species with shared biological features, genetic variation, and worldwide populations.

\medterm{Homocysteine} An amino acid formed during methionine metabolism. Persistently elevated blood levels may occur with genetic, nutritional, renal, or other disorders.

\medterm{Homocystinuria} A rare inherited disorder in which the body cannot properly process homocysteine, an amino acid. It may cause eye problems, abnormal bones, learning or developmental problems, and an increased risk of blood clots.

\medterm{Homograft} A transplant of tissue or an organ between two people of the same species but with different genetic makeup. The recipient’s immune system may reject it, so matching and immune-suppressing treatment are important.

\medterm{Homologous} Similar because of a shared origin, structure, or genetic relationship. In medicine, homologous tissue or cells come from the same species, while homologous chromosomes carry corresponding genes.

\medterm{Homologous Recombination} A DNA-repair and genetic-exchange process that uses a matching or nearly matching sequence as a template to repair breaks or reshuffle genetic material.

\medterm{Homologous Transplantation} Transplantation of cells, tissue, or an organ between genetically different people of the same species, requiring attention to compatibility, rejection, infection, and immunosuppression.

\medterm{Homonymous Hemianopia} A visual field deficit characterized by the loss of the same half of the visual field (either both right halves or both left halves) in both eyes. It results from a retrochiasmal lesion affecting the optic tract, lateral geniculate nucleus, optic radiation, or primary visual cortex on the side opposite the visual defect.

\medterm{Homosexual} Describing a person who experiences enduring romantic or sexual attraction to people of the same sex; sexual orientation is not a mental disorder.

\medterm{Homosexuality} Sexual or romantic attraction to people of the same sex or gender; it is not a mental disorder or pathology.

\medterm{Homovanillic Acid} A breakdown product of dopamine found in cerebrospinal fluid, blood, and urine; levels may reflect dopamine metabolism but are interpreted cautiously because many factors affect them.

\medterm{Homozygote} An individual or cell carrying two identical alleles at a particular gene locus, whether the alleles are normal, pathogenic, or another variant.

\medterm{Homozygous} Having two copies of the same version of a gene, one inherited from each parent. The effect depends on the gene and whether that version is harmless, protective, or disease-causing.

\medterm{Honey-Bee Stings} Injuries caused by a honey-bee injecting venom into the skin. A sting usually causes local pain, redness, and swelling; a severe allergic reaction can cause breathing difficulty, fainting, or shock and requires emergency treatment.

\medterm{Honolulu Fever} An infection caused by Streptobacillus bacteria, usually after contact with rats or contaminated food. It may cause fever, vomiting, headache, muscle or joint pain, and a rash and requires antibiotics.

\medterm{Hookworm} An intestinal infection caused by parasitic worms whose larvae live in contaminated soil. The larvae can enter through bare skin, and the adult worms may cause abdominal symptoms, blood loss, anemia, or poor growth.

\medterm{Hookworm Infection} Intestinal infection caused by hookworms acquired when larvae in contaminated soil
penetrate the skin or are ingested. Chronic infection can cause abdominal symptoms,
iron-deficiency anemia, protein loss, and impaired growth in children.

\medterm{Hookworm Infestation} Intestinal infection by hookworm nematodes, which attach to the intestinal wall and may cause iron-
deficiency anemia, abdominal symptoms, or poor growth.

\medterm{Hordeolum} Alternate medical name; see \textit{Stye}.

\medterm{Horizontal Gene Transfer} The movement of genetic material between microorganisms rather than from parent to offspring. It can spread traits such as antibiotic resistance through direct contact, viruses, or uptake of DNA from the surroundings.

\medterm{Hormonal Effects In Newborns} Transient newborn findings caused by maternal hormones or normal neonatal endocrine transitions, such as breast enlargement or vaginal bleeding; persistent or severe findings require evaluation.

\medterm{Hormone} A chemical messenger made by glands or other tissues and carried in the blood to target cells. It binds specific receptors to regulate functions such as growth, metabolism, reproduction, stress responses, and fluid balance.

\medterm{Hormone Levels} Measurements of hormones in blood, urine, or other specimens used to evaluate endocrine function; interpretation depends on the specific hormone, timing, age, sex, medications, and clinical context.

\medterm{Hormone Receptor} A protein that binds a specific hormone and converts the signal into cellular responses, either at the cell surface or inside the cell nucleus.

\medterm{Hormone Replacement Therapy} Treatment that supplies a hormone the body does not make enough of, such as thyroid, adrenal, sex, or pituitary hormones. Menopausal hormone therapy can relieve hot flushes and vaginal symptoms for selected people, but benefits and risks depend on the hormone, dose, duration, and health history.

\medterm{Hormone Therapy} Treatment that adds, lowers, blocks, or changes the action of hormones. In some breast and prostate cancers, it blocks hormones that the cancer needs to grow.

\medterm{Hormone Therapy For Breast Cancer} Treatment that blocks or lowers estrogen or progesterone effects in hormone-sensitive breast cancer. It reduces stimulation of cancer cells and may be used after surgery or for advanced disease.

\medterm{Hormone Therapy For Prostate Cancer} Treatment that lowers androgen production or blocks androgen action to slow growth of hormone-sensitive prostate cancer.

\medterm{Hormone-Receptor Positive} A cancer whose cells express receptors for hormones such as estrogen or progesterone. Hormone-blocking treatment may be useful for selected cancers with these receptors.

\medterm{Hormones} Chemical messengers released by glands or tissues that travel through blood or nearby spaces. They regulate growth, metabolism, reproduction, stress, fluid balance, mood, sleep, and many other body processes.

\medterm{Hormone-Sensitive Lipase} An enzyme in fat cells that breaks stored triglycerides into fatty acids and glycerol. Its activity rises when the body needs fuel, including during fasting or exercise, and falls after insulin increases.

\medterm{Horn} A horn is a pointed, hard growth of keratinized tissue, usually on the skin or as a normal anatomical projection such as a uterine or spinal horn.

\medterm{Horner Syndrome} A group of eye and face changes caused by damage to certain nerves. One side may have a smaller pupil, a drooping eyelid, and less sweating; the cause can range from migraine to a serious neck, chest, or brain problem.

\synonyms
Horner's Oculosympathetic Palsy; Bernard-Horner Syndrome

\medterm{Horner's Syndrome} A group of findings caused by disruption of sympathetic nerves, including a drooping eyelid, a small pupil, and reduced sweating on one side of the face.

\medterm{Hornet Stings} Hornet stings inject venom and usually cause local pain and swelling; multiple stings or allergy can cause a serious systemic reaction.

\medterm{Horripilation} Piloerection, or the temporary raising of hair and development of gooseflesh due to cold, fear, emotion, or autonomic stimulation.

\medterm{Horseshoe Kidney} A birth difference in which the lower parts of the two kidneys join across the middle of the body. Many people have no symptoms, but kidney stones, urine-flow problems, infections, or high blood pressure can occur.

\medterm{Hospice} Care for a person nearing the end of life when the focus is comfort rather than cure. Hospice supports symptoms, emotional and practical needs, and family caregivers at home, in a hospice facility, or in a hospital; eligibility rules vary by location.

\medterm{Hospice Care} Team-based care for a person nearing the end of life that prioritizes comfort, symptom control, communication, and family support rather than disease-directed cure.

\medterm{Hospice Rehabilitation} Rehabilitation provided during hospice care to maximize comfort, dignity, and participation rather than to restore function at all costs. It may include positioning, pain relief, safe transfers, equipment, and caregiver education.

\medterm{Hospital} A healthcare facility providing inpatient, emergency, surgical, diagnostic, intensive, rehabilitation, and specialist services for people needing medical care.

\medterm{Hospital Medicine} The medical specialty focused on comprehensive care of hospitalized patients, including diagnosis, treatment, coordination across teams, transitions of care, and prevention of inpatient complications.

\medterm{Hospital-Acquired Infection} An infection developing during or soon after hospital care that was not present on admission. It may involve wounds, urine, lungs, blood, or devices and can be caused by resistant organisms.
Examples: CLABSI, CAUTI, surgical site infection.

\medterm{Hospital-Acquired Infections} Infections that develop during or after hospitalisation. Major types include central
line-associated bloodstream infections, catheter-associated urinary tract infections,
ventilator-associated pneumonia, and surgical site infections. Risk factors include
invasive devices, prolonged hospitalisation, and antimicrobial use.

\medterm{Hospital-Acquired Pneumonia} Lung infection that begins at least 48 hours after hospital admission and was not present when the person arrived. It may cause fever, cough, breathing difficulty, or low oxygen; treatment depends on severity, cultures, and resistance risk.

\medterm{Hospital-Associated Disability} New or worsened difficulty with walking, self-care, or other daily activities that develops during or after hospitalization. It may result from illness, bed rest, weakness, delirium, poor nutrition, or treatments and can persist after discharge.

\medterm{Hospitalism} Physical or emotional problems associated with prolonged hospitalization, such as weakness, confusion, reduced mobility, sleep disruption, or loss of independence. Early movement, orientation, social contact, and individualized care can reduce these effects.

\medterm{Hospitalization} Admission to a hospital for monitoring, diagnosis, treatment, surgery, rehabilitation, or intensive support that cannot safely be delivered in an outpatient setting.

\medterm{Hospitals As Health Educators} Hospitals educate patients, families, staff, and communities about prevention, treatment, medicines, recovery, and when to seek care.

\medterm{Host} A person or other living organism that can be infected by a virus or other pathogen
under natural conditions. The host may be a human, animal, or plant. In medical
contexts, the host is the patient in whom an infection or disease develops.

\medterm{Host Cell} A cell in which an infectious agent, parasite, or symbiont lives, replicates, or uses cellular machinery; the host-cell response influences disease and clearance.

\medterm{Host Response} The body's reaction to an infection, foreign material, transplanted tissue, or other stimulus. It may include inflammation, immune defense, healing, scar formation, rejection, or tolerance.

\medterm{Hostility} An attitude or behavior marked by antagonism, anger, distrust, or intent to oppose another person; persistent severe hostility may occur with psychiatric, neurologic, substance-related, or situational problems.

\medterm{Hot Flash} Alternate singular form; see \textit{Hot Flashes}.

\medterm{Hot Flashes} Sudden episodes of intense heat, sweating, flushing, and sometimes a racing heart or chills. They are common during menopause but can also result from hormone changes, medicines, cancer treatment, or other illness.

\medterm{Hot Flushes} Alternate name; see \textit{Hot Flashes}.

\medterm{Hot Tub Folliculitis} A temporary bacterial infection of hair follicles, commonly caused by Pseudomonas aeruginosa after exposure to contaminated hot tubs or pools, producing itchy follicular papules or pustules.

\medterm{Hounsfield Unit} A number on a CT scan that describes how strongly a tissue absorbs X-rays. Air has a very low value, water is near zero, and bone has a high value; the numbers help distinguish tissues and lesions.

\medterm{Hounsfield Units} Numbers used on CT scans to describe how strongly tissues absorb X-rays. Air has a very low value, water is near zero, and bone has a high value; the numbers help distinguish tissues and abnormal areas.

\medterm{Hour} A unit of time equal to 60 minutes, commonly used in medicine to describe symptom duration, dosing intervals, monitoring periods, and treatment timing.

\medterm{Hour-Glass Contraction} An abnormal tight band that divides the uterus into upper and lower sections during difficult labor. It can prevent the baby from being delivered normally and requires urgent assessment by an obstetric team.

\medterm{House Officer} A physician receiving postgraduate clinical training in a hospital, such as an intern, resident, or fellow; usage varies by country and institution.

\medterm{House Staff} The physicians and other trainees who provide patient care within a hospital under the supervision of attending clinicians, commonly including interns, residents, and fellows.

\medterm{Household Glue Poisoning} Harm from swallowing or inhaling household glue. Small accidental amounts may cause irritation or stomach upset, while solvents can cause dizziness, unconsciousness, breathing problems, or heart-rhythm disturbances; seek poison-control advice.

\medterm{Housemaid's Knee} A common name for prepatellar bursitis, inflammation and fluid accumulation in the bursa over the kneecap, often caused by prolonged kneeling, trauma, infection, or inflammatory disease. Painful swelling, redness, warmth, or fever requires assessment for septic bursitis.

\medterm{Howell-Jolly Bodies} Small leftover pieces of red-cell DNA seen on a blood smear. They are usually removed by a working spleen, so their presence may indicate absent or reduced spleen function; they can also occur with severe vitamin-B12 or folate deficiency.

\medterm{HPA} A hormone system linking parts of the brain with the adrenal glands. It helps the body respond to stress by controlling the release of cortisol, which affects blood pressure, blood sugar, immune activity, sleep, and energy.

\synonyms
Hypothalamic-Pituitary-Adrenal Axis

\medterm{HPLC} A laboratory method that separates substances in a liquid sample so they can be identified and measured. It is used to test medicines, body chemicals, poisons, and other samples.

\medterm{HPV} Human papillomavirus is a group of viruses that infect the skin and moist body linings. Some types cause warts, while long-lasting infection with high-risk types can cause precancerous changes and cancers of the cervix, anus, genitals, or throat.

\medterm{HPV Test} A test that detects high-risk human papillomavirus genetic material or related markers in cervical samples to screen for cervical cancer risk or evaluate abnormal cervical cytology.

\medterm{HPylori} A shorthand name for \textit{Helicobacter pylori}, a bacterium that can live in the stomach and cause inflammation, ulcers, and increased stomach-cancer risk. Confirmed infection is treated with a combination of medicines and follow-up testing.

\medterm{HRCT} Alternate name; see \textit{High-Resolution Computed Tomography}.

\medterm{HRT} Hormone replacement therapy uses estrogen, with a progestogen when needed, to treat selected menopausal
symptoms caused by declining ovarian hormone production. It can relieve vasomotor and
genitourinary symptoms and help prevent bone loss, but benefits and risks depend on the
person’s age, medical history, regimen, and duration of use.

\medterm{Hürthle Cell Tumour} A growth arising from hormone-making cells in the thyroid. It may be benign or cancerous; ultrasound and examination of thyroid tissue help determine its nature and treatment.

\medterm{Hs-cTn} Alternate name; see \textit{High-Sensitivity Cardiac Troponin}.

\medterm{Ht} Common abbreviation for hematocrit, the percentage of blood volume occupied by red blood cells; in other contexts it may mean height.

\medterm{Huber Needle} A special non-coring needle used to access an implanted port for giving medicines, fluids, blood products, or taking blood samples. Correct placement and sterile technique reduce infection, bleeding, and port damage.

\medterm{Human Calcitonin} A peptide hormone produced by thyroid parafollicular cells that can lower blood calcium by inhibiting osteoclast activity; synthetic calcitonin is used selectively for bone and calcium disorders.

\medterm{Human Chorionic Gonadotropin} A hormone produced mainly by the placenta and measured in blood or urine to detect and monitor pregnancy. Levels normally rise early in pregnancy, but the result must be interpreted with timing, symptoms, and other findings.

\synonyms
HCG; Beta-HCG

\medterm{Human Chromosome Count} Most human body cells contain 46 chromosomes arranged in 23 pairs, including one pair of sex chromosomes; mature egg and sperm cells normally contain 23.

\medterm{Human Embryo} The developing human organism from fertilization through approximately the eighth week, during which the major body plan and organ systems are established.

\medterm{Human Gastrin} A peptide hormone produced mainly by gastric G cells that stimulates gastric acid secretion and supports growth of the gastric mucosa after food enters the stomach.

\medterm{Human Genome} The complete DNA sequence and genetic content of a human, organized into nuclear chromosomes and mitochondrial DNA; variants influence traits, disease risk, and treatment response.

\medterm{Human Herpesvirus 6} A herpesvirus that commonly infects children and can cause roseola; it may reactivate and cause serious disease in people with weakened immunity.

\medterm{Human Herpesvirus 8} A herpesvirus associated with Kaposi sarcoma, primary effusion lymphoma, and multicentric Castleman disease, especially in people with weakened immunity.

\medterm{Human Immunodeficiency Virus} A virus that attacks immune cells, especially CD4 cells, and can gradually weaken the body’s defenses. Without effective treatment it may lead to AIDS and serious infections or cancers. HIV-1 is widespread globally and HIV-2 is less common; daily treatment can suppress the virus and protect health.

\medterm{Human Leukocyte Antigen} A group of cell-surface proteins encoded by genes in the major histocompatibility complex. HLA molecules present antigens to T cells and are important in immune responses and tissue matching for transplantation.

\synonyms
HLA Complex; Human MHC

\medterm{Human Metapneumovirus} A respiratory virus that commonly causes cough, fever, runny nose, and wheezing. Infants, older adults, and people with weakened immunity may develop bronchiolitis, pneumonia, or severe breathing difficulty.

\medterm{Human Microbiome} The community of bacteria, viruses, fungi, archaea, and other microorganisms living on and inside the human body. These organisms and their genes can influence digestion, immune activity, metabolism, and disease, although finding a microorganism does not by itself prove that it causes illness.

\medterm{Human Nociceptin} A neuropeptide that binds the nociceptin/orphanin FQ receptor and modulates pain, stress, reward, autonomic function, and other brain and spinal-cord processes.

\medterm{Human Papillomavirus} A large group of viruses spread mainly through intimate skin-to-skin or sexual contact. Most infections cause no symptoms and clear without treatment. Some types cause warts, while persistent infection with high-risk types can cause precancerous changes and cancers of the cervix, anus, genitals, or throat. Vaccination and recommended cervical screening reduce risk.

\medterm{Human Papillomavirus} A large family of viruses spread mainly through intimate skin-to-skin or sexual contact. Some types cause warts and others can cause cancers; vaccination and recommended screening reduce the risk of serious disease.

\medterm{Human Papillomavirus Infection} Alternate name; see \textit{Human Papillomavirus}.

\medterm{Human Placental Lactogen} A hormone made by the placenta during pregnancy. It helps prepare the breasts for feeding and changes how the pregnant person uses fat and sugar, helping provide nutrients to the developing fetus.

\medterm{Humectant} A substance that attracts and holds water, often used in skin creams, lotions, foods, and medicines. In skincare, it draws water into the outer skin layer and can reduce dryness when used properly.

\medterm{Humeral Epicondyle} A bony prominence at the distal humerus that provides attachment for forearm muscles and ligaments; the medial and lateral epicondyles are located on opposite sides of the elbow.

\medterm{Humeral Fracture} A break in the upper-arm bone, involving the proximal, shaft, or distal humerus; pain, swelling, deformity, and possible nerve or blood-vessel injury guide urgent assessment and treatment.

\medterm{Humeroradial Joint} The part of the elbow where the upper-arm bone meets the radius, one of the forearm bones. It helps the elbow bend and straighten and supports rotation of the forearm.

\medterm{Humeroulnar Joint} The main hinge of the elbow, where the upper-arm bone meets the ulna. It allows the arm to bend and straighten and provides much of the elbow’s stability.

\medterm{Humerus} The long bone of the upper arm, extending from the shoulder to the elbow. It joins the shoulder blade above and the radius and ulna below; fractures can injure nearby nerves or blood vessels.

\medterm{Humidifier} A device that adds moisture to air; medical humidifiers require appropriate cleaning and water handling because contaminated equipment can disperse infectious or inflammatory material.

\medterm{Humidifier Fever} A flu-like illness caused by inhalation of microorganisms or biologic contaminants from a contaminated humidifier or ventilation system, often involving fever, chills, cough, and malaise.

\medterm{Humidity} The amount of water vapor in air, influencing heat loss, respiratory comfort, skin hydration, and survival or transmission of some microorganisms.

\medterm{Humor} In anatomy, a fluid of the body, especially the aqueous or vitreous humor of the eye; in ordinary language, it means a sense of amusement or comedy.

\medterm{Humoral Immunity} The part of the immune system that protects mainly through antibodies made by B cells. Antibodies can block germs and toxins and help other immune cells remove them.

\medterm{Hunchback} A nontechnical term usually referring to pronounced kyphosis, an excessive forward curvature of the upper spine. Causes include osteoporosis, vertebral compression fractures, Scheuermann disease, degenerative change, and congenital disorders; evaluation depends on age, progression, pain, and neurologic findings.

\medterm{Hunger} The physical drive to eat, influenced by stomach fullness, blood glucose, hormones such as ghrelin, emotions, habits, illness, and the availability or appeal of food.

\medterm{Hungry Bone Syndrome} A severe drop in blood calcium after treatment of overactive parathyroid glands, often after surgery. Healing bones rapidly take up calcium and other minerals, causing tingling, cramps, spasms, seizures, or abnormal heart rhythms.

\medterm{Hunter Syndrome} An inherited condition, usually affecting boys, in which iduronate-2-sulfatase deficiency causes complex sugars to accumulate in tissues. It can cause coarse facial features, enlarged organs, stiff joints, hearing loss, breathing problems, and developmental difficulties.

\medterm{Hunter's Syndrome} Another name for Hunter syndrome, an inherited disorder in which complex sugars accumulate because the body cannot break them down properly. It mainly affects boys and can damage several organs over time.

\medterm{Huntington Disease} A progressive inherited brain disorder caused by an abnormally long repeated section of DNA in the \textit{HTT} gene. The altered gene produces an abnormal huntingtin protein that damages and eventually kills nerve cells involved in movement, thinking, mood, and behavior. Symptoms may include involuntary, dance-like movements (chorea), muscle stiffness, poor coordination, changes in speech or swallowing, depression, irritability, and worsening memory or judgment. It usually begins in adulthood, but a less common juvenile form can begin during childhood or adolescence. The condition is inherited in an autosomal dominant pattern, so each child of a parent with one altered copy has a 50\% chance of inheriting it.

\synonyms
Huntington's Chorea; HD

\medterm{Hurler Syndrome} Mucopolysaccharidosis type I (MPS I), the severe form, an autosomal recessive lysosomal
storage disorder caused by alpha-L-iduronidase deficiency (IDUA gene mutations).
Accumulation of dermatan sulfate and heparan sulfate leads to progressive multisystem
involvement.

\medterm{Hurthle Cell} A thyroid follicular cell with abundant granular mitochondria-rich cytoplasm, seen in Hashimoto thyroiditis and Hürthle-cell neoplasms.

\medterm{Hurthle Cell Cancer} A rare thyroid cancer arising from Hurthle cells, with behavior ranging from indolent to aggressive.
\synonyms
Oxyphilic Thyroid Carcinoma; Oncocytoid Carcinoma

\medterm{Hutchinson's Teeth} Hutchinson's teeth are permanently notched, widely spaced upper front teeth classically associated with congenital syphilis.

\medterm{Hyaline} A glassy, translucent material or appearance, commonly describing proteinaceous deposits or homogeneous extracellular material seen microscopically in tissues.

\medterm{Hyaline Cartilage} The most common type of cartilage, containing fine type II collagen fibers in a firm,
glassy matrix. It supports airways, covers articular surfaces, and forms much of the
fetal skeleton and growth plates.

\medterm{Hyaline Casts} Clear, tube-shaped particles sometimes seen in urine. A small number may be normal, especially after exercise or dehydration; larger numbers can occur when the kidneys are under stress and must be interpreted with other urine and blood tests.

\medterm{Hyaline Membrane} A protein-rich layer that forms along the walls of damaged air sacs in the lungs. It blocks oxygen transfer and is classically seen in newborn respiratory distress caused by inadequate surfactant.

\medterm{Hyaline Membrane Disease} A former name for breathing failure in newborns, especially premature infants, caused mainly by too little lung surfactant. The air sacs collapse easily, making oxygen transfer difficult.

\medterm{Hyalitis} Inflammation of the vitreous body of the eye, usually associated with intraocular infection, autoimmune inflammation, trauma, or another ocular disorder. It may cause floaters, blurred vision, or visual loss and requires ophthalmic assessment.

\medterm{Hyaloid Canal} A narrow channel through the vitreous body of the eye that marks the former path of the hyaloid artery from the optic disc toward the lens during fetal development.

\medterm{Hyaluronic Acid} A naturally occurring carbohydrate polymer that binds water in connective tissue, skin, joints, and the eye.

\synonyms
Hyaluronan; Sodium Hyaluronate

\medterm{Hyaluronidase} An enzyme that breaks down hyaluronic acid in connective tissue. Medicines containing it can help injected fluids or medicines spread through tissue and may also dissolve certain hyaluronic-acid cosmetic fillers.

\medterm{Hyaluronoglucosaminidase} An enzyme that degrades hyaluronan and related glycosaminoglycans; hyaluronidase preparations can increase tissue permeability and assist dispersion of injected fluids or medicines.

\medterm{Hybridoma} A cell line formed by fusing an antibody-producing B lymphocyte with an immortal myeloma cell, allowing continuous production of a monoclonal antibody.

\medterm{Hycanthone} An older antiparasitic drug formerly used against schistosomiasis; its use was abandoned or restricted because of serious toxicity, including liver injury and carcinogenic concern.

\medterm{Hydatid Cyst} A cyst caused by the larval stage of Echinococcus granulosus, a tapeworm transmitted
from dogs to humans through ingestion of contaminated food. The liver is the most
commonly affected organ, followed by the lungs. Cysts grow slowly and may become
symptomatic from mass effect or rupture, causing anaphylaxis.

\medterm{Hydatid Disease} A zoonotic parasitic infection caused by the larval stage of tapeworms of the genus \textit{Echinococcus} (principally \textit{Echinococcus granulosus}). Humans acquire the infection by accidentally ingesting microscopic eggs shed in the feces of definitive canine hosts via contaminated food, water, soil, or direct animal contact. Larvae hatch in the small intestine, penetrate the bowel wall, and enter the circulation, most frequently lodging in the liver (approx.\ 70\%) and lungs (approx.\ 20\%), where they form slow-growing, fluid-filled cysts containing daughter cysts and hydatid sand. Most cysts remain asymptomatic for years until local mass effect causes abdominal discomfort, hepatomegaly, cough, or hemoptysis. Spontaneous or traumatic cyst rupture can cause severe anaphylaxis and secondary dissemination. Diagnosis is established by imaging (ultrasound, CT, or MRI showing cyst wall calcification, daughter cysts, or detached internal membranes) and confirmatory echinococcal serology. Management includes antiparasitic therapy (albendazole), image-guided aspiration and sterilization (PAIR: Puncture, Aspiration, Injection, Re-aspiration), or surgical cyst excision.
\synonyms{Echinococcosis; Cystic echinococcosis; Hydatidosis}

\medterm{Hydatidiform Mole} A form of gestational trophoblastic disease in which abnormal trophoblastic tissue grows in
the uterus. A complete mole has no viable fetus; a partial mole may contain abnormal fetal tissue. Both require follow-up
because persistent disease can occur.

\medterm{Hydralazine} Hydralazine is a direct arteriolar vasodilator used for selected hypertension, including severe hypertension in pregnancy; reflex tachycardia, fluid retention, headache, and drug-induced lupus are recognized adverse effects.

\medterm{Hydramnios} Excess amniotic fluid during pregnancy, also called polyhydramnios, caused by fetal swallowing impairment, maternal diabetes, multiple gestation, fetal anomalies, infection, or sometimes no identifiable cause.

\medterm{Hydranencephaly} A rare brain development disorder in which most of the upper brain is replaced by fluid-filled spaces. It can cause severe developmental and movement problems, seizures, feeding difficulty, and a shortened life expectancy.

\medterm{Hydrarthrosis} Abnormal accumulation of clear synovial fluid within a joint, producing swelling and restricted movement; causes include osteoarthritis, inflammatory arthritis, trauma, and infection.

\medterm{Hydration} The body’s water balance and the process of replacing water and electrolytes lost through urine, stool, sweat, breathing, vomiting, or bleeding.

\medterm{Hydrazine} A highly reactive chemical used in fuels and industry that can cause irritant, neurologic, hepatic, and systemic toxicity after significant exposure.

\medterm{Hydrazine Sulfate} A hydrazine-containing chemical once promoted experimentally for cancer-related symptoms; evidence of benefit is inadequate and toxicity can include liver injury, neuropathy, and seizures.

\medterm{Hydrazines} A family of nitrogen-containing chemical compounds used in industry, agriculture, and pharmaceuticals; some are irritants, toxic, mutagenic, or hepatotoxic.

\medterm{Hydrazones} Chemical compounds formed by reaction of hydrazine derivatives with aldehydes or ketones, important in medicinal chemistry, analytical chemistry, and some drug structures.

\medterm{Hydroa Vacciniforme} A rare photodermatosis causing recurrent vesicles or papules on sun-exposed skin that heal with varioliform scars; some cases are associated with Epstein-Barr virus.

\medterm{Hydrocarbon} An organic compound composed only of carbon and hydrogen; ingestion or inhalation of petroleum hydrocarbons can cause chemical pneumonitis and systemic toxicity.

\medterm{Hydrocarbon Pneumonia} Lung inflammation caused by inhaling or aspirating petroleum products or other hydrocarbons, which can produce cough, breathing difficulty, and chemical pneumonitis.

\medterm{Hydrocele} A collection of serous fluid within the tunica vaginalis surrounding the testis.
Congenital hydrocele results from patent processus vaginalis; acquired hydrocele may
follow infection, trauma, or tumor. Transilluminates on examination. Ultrasound confirms
fluid collection.

\medterm{Hydrocele Repair} Surgery to close or remove the fluid-containing sac around a testicle when a hydrocele is persistent, large, painful, associated with a hernia, or obscures examination.

\medterm{Hydrocephalus} Hydrocephalus is an abnormal buildup of cerebrospinal fluid (CSF) in the brain’s fluid-filled spaces, called ventricles. It can enlarge the ventricles and sometimes increase pressure on brain tissue. It may result from blocked CSF flow, reduced absorption, bleeding, infection, a tumor, injury, or a developmental condition. Symptoms vary with age and may include headache, vomiting, sleepiness, vision problems, poor balance, walking difficulty, memory changes, or an increasing head size in infants.

\medterm{Hydrocephalus Acquired} Acquired hydrocephalus is abnormal accumulation of cerebrospinal fluid after birth due to obstruction, impaired absorption, or rarely excess production; raised intracranial pressure may cause headache, vomiting, gait or cognitive change, and visual symptoms.

\medterm{Hydrocephalus Ex-Vacuo} Enlargement of the brain's ventricles secondary to loss or shrinkage of surrounding brain tissue, rather than primarily from obstructed cerebrospinal-fluid flow.

\medterm{Hydrocephaly} Hydrocephalus: abnormal accumulation or circulation of cerebrospinal fluid that enlarges the ventricles and may increase pressure on the brain.

\medterm{Hydrochloric Acid} The strong acid in stomach juice that helps digest proteins and kill many swallowed microorganisms. The stomach’s mucus lining and other protective mechanisms normally prevent this acid from damaging its tissues.

\medterm{Hydrochloric Acid Poisoning} A corrosive injury after swallowing, inhaling, or contacting hydrochloric acid. It can burn the mouth, throat, esophagus, stomach, skin, or lungs; do not induce vomiting and seek emergency care immediately.

\medterm{Hydrochlorothiazide} Hydrochlorothiazide is a thiazide diuretic that inhibits sodium-chloride reabsorption in the distal convoluted tubule and treats hypertension and mild oedema; hypokalaemia, hyponatraemia, hyperuricaemia, and photosensitivity may occur.

\medterm{Hydrocodone} An opioid analgesic used for moderate to severe pain and, in some formulations, cough; it can cause respiratory depression, sedation, constipation, tolerance, dependence, and overdose.

\medterm{Hydrocodone And Acetaminophen Overdose} Taking too much of this opioid-and-pain-reliever combination can slow breathing and cause coma, while acetaminophen can severely damage the liver. Emergency treatment is needed even before symptoms appear.

\medterm{Hydrocodone Bitartrate} The bitartrate salt of hydrocodone, an opioid analgesic that can cause sedation, constipation, respiratory depression, tolerance, dependence, and overdose.

\medterm{Hydrocodone/Oxycodone Overdose} Taking too much hydrocodone or oxycodone can cause extreme sleepiness, pinpoint pupils, slow or stopped breathing, coma, and death. Call emergency services, give naloxone if available, and provide rescue breathing if trained.

\medterm{Hydrocolpos} Distension of the vagina with fluid due to obstruction of the vaginal outlet, often from
an imperforate hymen or transverse septum, presenting as a pelvic mass in neonates or
adolescents.

\medterm{Hydrocortisone} A corticosteroid identical or similar to the body’s cortisol. It reduces inflammation and immune activity and treats adrenal hormone deficiency; prolonged or high-dose use can cause infection, high blood sugar, skin thinning, or adrenal suppression.

\medterm{Hydrocyanic Acid Poisoning} A life-threatening poisoning that prevents cells from using oxygen. It can cause sudden headache, confusion, seizures, breathing failure, cardiovascular collapse, and death; emergency antidote treatment is required immediately.

\medterm{Hydroflumethiazide} An older thiazide diuretic that increases urinary salt and water loss to lower blood pressure or reduce swelling. It can disturb sodium, potassium, and glucose levels and may cause dehydration or increased uric acid.

\medterm{Hydrofluoric Acid Exposure} Contact with hydrofluoric acid, which can cause deep tissue injury and systemic fluoride
toxicity, including dangerous low calcium and cardiac arrhythmias.

\medterm{Hydrofluoric Acid Poisoning} A medical emergency after contact with or inhalation of hydrofluoric acid. It can cause deep tissue burns and dangerous low calcium, heart rhythm changes, or death, even when the skin injury looks minor.

\medterm{Hydrogel} A water-rich three-dimensional material that can hold and release substances. Hydrogels are used in wound dressings, contact lenses, drug delivery, tissue engineering, and research, with safety depending on composition and use.

\medterm{Hydrogen Bond} A weak attraction between a hydrogen atom attached to an electronegative atom and another electronegative atom. These bonds help maintain the shapes of proteins and DNA and influence water’s properties.

\medterm{Hydrogen Breath Test} A test for some causes of bloating, diarrhea, or abdominal pain. After drinking a measured sugar solution, the person breathes into collection tubes at intervals. The gases can suggest poor sugar absorption or excess bacteria in the small intestine.

\medterm{Hydrogen Cyanide} A rapidly acting cellular toxin that inhibits cytochrome oxidase and prevents oxidative phosphorylation; significant exposure can cause severe lactic acidosis, seizures, cardiovascular collapse, and death.

\medterm{Hydrogen Peroxide} An oxidizing chemical used in some antiseptic and industrial products; ingestion or concentrated exposure can injure tissue and release gas, causing embolic or respiratory complications.

\medterm{Hydrogen Peroxide Poisoning} Harm after swallowing or inhaling hydrogen peroxide. It may irritate or burn tissues and release gas that blocks blood vessels; concentrated products can cause severe injury, so urgent poison-control or emergency advice is required.

\medterm{Hydrogen Sulfide} A toxic, flammable gas with a rotten-egg odor at low concentrations that inhibits cellular respiration and can rapidly cause respiratory failure, neurologic effects, or death.

\medterm{Hydrogenation} A chemical process that adds hydrogen to a compound; in foods it can convert liquid oils into more solid fats.

\medterm{Hydrogen-Ion Concentration} The concentration of hydrogen ions in a solution, expressed through pH and central to acid--base physiology, enzyme activity, and interpretation of blood-gas results.

\medterm{Hydromorphone} Hydromorphone is a potent opioid analgesic for severe pain; sedation, constipation, respiratory depression, tolerance, and dependence are important risks, especially with alcohol or other sedatives.

\medterm{Hydromorphone Overdose} Taking too much hydromorphone, an opioid, can cause extreme sleepiness, pinpoint pupils, slow or stopped breathing, coma, and death. Call emergency services and give naloxone if available; urgent monitoring is essential.

\medterm{Hydromyelia} Abnormal widening of the central canal of the spinal cord, which may be congenital or associated with obstruction of cerebrospinal-fluid flow and neurologic symptoms.

\medterm{Hydronephrosis} Swelling of a kidney because urine cannot drain normally. Causes include a stone, narrowing, reflux, a tumor, or prostate enlargement. Prolonged or severe blockage can damage the kidney and may cause pain, infection, or reduced kidney function.

\medterm{Hydronephrosis In Newborns} Swelling of a newborn’s kidney caused by widened urine-collecting spaces. It may result from temporary changes, blockage, or urine flowing backward and can affect one or both kidneys.

\medterm{Hydronephrosis Of One Kidney} Dilation of one kidney’s renal pelvis and calyces caused by impaired urine drainage, often from a stone, stricture, reflux, or obstruction of the ureter or ureteropelvic junction.

\medterm{Hydrophobia} An involuntary, agonizing spasm of the throat and vocal cord muscles triggered by attempting to drink or by the sight, sound, or mention of water. It is the hallmark neurological sign of furious (encephalitic) rabies, caused by viral damage to brainstem swallowing centers. The severe choking and pain make swallowing impossible, leading to a terrified avoidance of liquids, excessive drooling (hypersalivation), and often aerophobia (spasms triggered by a draft of air). Once hydrophobia appears, rabies is nearly 100\% fatal, highlighting the critical need for immediate post-exposure rabies vaccination following animal bites.
\synonyms{Rabic Hydrophobia; Rabid Pharyngeal Spasms; Fear of Water in Rabies}

\medterm{Hydropneumothorax} Air and fluid collected together in the space around a lung. It can compress the lung and cause chest pain or severe breathing difficulty, often requiring urgent hospital treatment and drainage.

\medterm{Hydrops} Abnormal accumulation of fluid in tissues or body cavities. The term may refer to generalized fetal edema, a localized collection, or organ-specific fluid distension; the cause depends on the clinical setting.

\medterm{Hydrops Fetalis} Severe fluid buildup in an unborn baby, affecting at least two areas such as under the skin, around the lungs, around the heart, or in the abdomen. Causes include blood-group incompatibility, anemia, infection, genetic conditions, heart problems, and complications of twin pregnancy.

\medterm{Hydrops Of The Gallbladder} Distension of the gallbladder by clear mucus or watery fluid, usually caused by prolonged obstruction of the cystic duct by a gallstone. It may cause right-upper-quadrant pain and can progress to inflammation or infection.

\medterm{Hydrosalpinx} Dilatation of a fallopian tube caused by accumulation of clear fluid, usually after inflammation, infection, endometriosis, or previous pelvic surgery. It may affect fertility.

\medterm{Hydrotherapy} Exercise or treatment performed in water. Water supports body weight and provides gentle resistance, which may reduce pain and help movement, strength, and balance during rehabilitation. The program should match the person’s condition and safety needs.

\medterm{Hydrothorax} An abnormal collection of clear fluid in the pleural space around a lung. It may result from heart, liver, kidney, or other disease and can cause breathlessness, cough, or chest pressure.

\medterm{Hydroureter} Abnormal enlargement of a ureter caused by impaired urine drainage.

\medterm{Hydroxamic Acid} An organic functional group containing hydroxylamine attached to a carbonyl, found in chelators and some medicines, including histone-deacetylase inhibitors.

\medterm{Hydroxocobalamin} A cobalt-containing form of vitamin B12 used as an antidote for cyanide poisoning. It
binds cyanide to form cyanocobalamin, which is excreted in urine; it can cause red
discoloration of skin and body fluids and interfere with some laboratory tests.

\medterm{Hydroxyapatite} A calcium-and-phosphate mineral that gives much of their hardness to bones and teeth. Bone cells continually remove and replace it as bone grows, repairs itself, and responds to changing physical forces.

\medterm{Hydroxychloroquine} Hydroxychloroquine is an antimalarial and immunomodulatory drug used for malaria prevention or treatment and selected autoimmune diseases; cumulative retinal toxicity requires ophthalmic monitoring.

\medterm{Hydroxycholesterols} Oxidized forms of cholesterol produced in the body or during food processing. They participate in cell signaling and cholesterol regulation and may contribute to inflammation or blood-vessel disease when present in excess.

\medterm{Hydroxylation} An enzymatic or chemical reaction that adds a hydroxyl group to a molecule, often altering the activity, solubility, metabolism, or excretion of a drug or biomolecule.

\medterm{Hydroxyproline} An amino acid formed mainly by post-translational modification of collagen; its measurement has been used as an indicator of collagen turnover but is not a routine diagnostic test.

\medterm{Hydroxyprolinemia} A rare inherited metabolic finding caused by impaired hydroxyproline breakdown, leading to elevated hydroxyproline in blood or urine and usually little clinical illness.

\medterm{Hydroxysteroid} A steroid hormone or steroid-related compound containing a hydroxyl group; hydroxysteroids include intermediates and products of adrenal and gonadal steroid metabolism.

\medterm{Hydroxyurea} A medicine used for some blood cancers and sickle-cell disease. In sickle-cell disease it can reduce painful episodes and other complications by helping red blood cells remain less likely to sickle; blood counts and other safety tests are required.

\medterm{Hydroxyzine} Hydroxyzine is a sedating H1-antihistamine used for itching, anxiety, and preoperative sedation; drowsiness, anticholinergic effects, and QT prolongation can occur.

\medterm{Hydroxyzine Overdose} Taking too much hydroxyzine can cause severe drowsiness, confusion, agitation, seizures, low blood pressure, or dangerous heart-rhythm changes. Seek urgent poison-control or emergency advice, especially after a large or intentional dose.

\medterm{Hydrozoa} A group of mostly marine animals that includes hydroids and some jellyfish. Contact with certain species can cause painful stings or skin irritation; the medical risk depends on the species and exposure.

\medterm{Hygiene} Practices that reduce exposure to pathogens and support health, including hand hygiene, oral care, skin care, food safety, and appropriate environmental cleanliness.

\medterm{Hygiene Hypothesis} The theory that limited early-life microbial exposure may increase susceptibility to allergies and immune disorders.

\medterm{Hygroma} A fluid-filled swelling, often caused by leakage or blockage of lymphatic vessels. A cystic hygroma usually appears in the neck of a fetus or child, while other hygromas may occur near joints or after injury.

\medterm{Hymen} A thin fold of mucous membrane that partially covers the external opening of the vagina.
The hymen is a remnant of fetal development, with no known physiological function in
postnatal life.

\medterm{Hymenectomy} Hymenectomy is a procedure to remove or open part of the hymen, usually to treat an obstructed hymenal opening or another symptomatic abnormality.

\medterm{Hymenolepiasis} Intestinal infection with dwarf tapeworms, especially Hymenolepis nana, acquired by ingesting eggs; it may cause abdominal symptoms, diarrhea, or be asymptomatic.

\medterm{Hymenolepis Diminuta} A tapeworm that can infect people who accidentally swallow infected insects, usually in contaminated food. Infection may cause abdominal discomfort, diarrhea, or no symptoms and is treated with an antiparasitic medicine.

\medterm{Hymenolepis Nana} A small tapeworm that spreads when microscopic eggs are swallowed, often through contaminated hands, food, or water. It can cause abdominal pain, diarrhea, poor appetite, or no symptoms and is treatable with antiparasitic medicine.

\medterm{Hymenoptera} An insect order that includes bees, wasps, hornets, and ants; stings or bites can cause pain, allergy, or life-threatening anaphylaxis.

\medterm{Hyoid Bone} U-shaped bone in the anterior neck at C3, unique for not directly articulating with any
other bone. Suspended by muscles and the stylohyoid ligament. Provides attachment for
suprahyoid (digastric, stylohyoid, mylohyoid, geniohyoid) and infrahyoid (sternohyoid,
sternothyroid, thyrohyoid, omohyoid) muscles.

\medterm{Hyoscine} A medicine that blocks muscarinic acetylcholine receptors. Depending on the preparation, it can prevent motion sickness, reduce spasms, or lessen secretions; dry mouth, blurred vision, constipation, confusion, and urinary retention may occur.

\medterm{Hyoscyamine} An antimuscarinic medicine that reduces smooth-muscle spasm and secretions; it may be used for selected gastrointestinal or urinary symptoms and can cause dry mouth, blurred vision, urinary retention, or confusion.

\medterm{Hyperactivity} Excessive or developmentally inappropriate motor activity, restlessness, or difficulty remaining still; persistent impairment may occur in ADHD, mania, substance effects, or other conditions.

\medterm{Hyperacusis} Reduced tolerance to ordinary environmental sounds, which may be perceived as excessively loud, painful, or distressing. It can occur with hearing disorders, migraine, head injury, neurologic conditions, or without an identifiable cause.

\medterm{Hyperadrenocorticism} Excess production or exposure to adrenal glucocorticoids, especially cortisol; in humans it causes Cushing syndrome, which may result from adrenal, pituitary, ectopic, or medication-related causes.

\medterm{Hyperaemia} Increased blood flow to a tissue, making it look red or feel warm. It may be a normal response to activity or a response to inflammation or renewed blood flow after temporary blockage.

\medterm{Hyperaesthesia} Unusually strong sensitivity to touch, pain, temperature, or other sensations. It may occur with nerve injury, inflammation, migraine, medicines, or disorders affecting the brain or spinal cord.

\medterm{Hyperaldosteronism} Excess production of the adrenal hormone aldosterone. It can cause high blood pressure, low potassium, muscle weakness, cramps, thirst, and frequent urination; an adrenal growth or overactivity of both glands may be responsible.

\medterm{Hyperaldosteronism Complications} Beyond hypertension and hypokalemia, primary hyperaldosteronism causes target organ
damage including left ventricular hypertrophy, fibrosis, stroke, coronary artery
disease, atrial fibrillation, and chronic kidney disease. These complications occur at
higher rates than in essential hypertension of similar severity.

\medterm{Hyperalgesia} Feeling more pain than expected from something that normally hurts. It can happen with inflammation, nerve injury, changes in nervous-system pain processing, or opioid use; it differs from allodynia, which is pain from a normally painless touch.

\medterm{Hyperammonemia} An abnormally high level of ammonia in the blood. It may result from severe liver disease, an inherited metabolic disorder, medicines, or other illness and can cause confusion, sleepiness, vomiting, seizures, coma, or brain swelling.

\medterm{Hyperamylasemia} An elevated blood amylase concentration caused by pancreatic or salivary-gland disease, reduced clearance, macroamylasemia, or other conditions; it is not diagnostic of pancreatitis by itself.

\medterm{Hyperargininemia} An inherited urea-cycle disorder caused by arginase deficiency, leading to high arginine and progressive spasticity, developmental impairment, seizures, or liver-related findings.

\medterm{Hyperbaric} Having a pressure greater than the surrounding atmospheric pressure; hyperbaric oxygen therapy delivers oxygen at increased pressure for selected medical indications.

\medterm{Hyperbaric Chamber} A pressurized chamber in which a person breathes oxygen or another gas at pressure above sea level, used for conditions such as decompression sickness, air embolism, and selected wounds.

\medterm{Hyperbaric Oxygen Therapy} Breathing oxygen inside a chamber pressurized above normal air pressure. The treatment increases oxygen dissolved in the blood and is used for selected problems such as carbon-monoxide poisoning, decompression sickness, and certain severe wounds or infections.

\medterm{Hyperbilirubinaemia} Hyperbilirubinaemia is an abnormally high level of bilirubin in the blood, which can cause jaundice and may result from increased breakdown of red cells, liver disease, or blocked bile flow.

\medterm{Hyperbilirubinemia} An abnormally high concentration of bilirubin in blood that may cause jaundice; causes include hemolysis, impaired liver processing, hepatocellular injury, and bile-duct obstruction.

\medterm{Hypercalcaemia} Alternate British spelling; see \textit{Hypercalcemia}.

\medterm{Hypercalcemia} An abnormally high level of calcium in the blood (total calcium greater than $10.5\,\text{mg/dL}$ or $2.6\,\text{mmol/L}$). More than 90\% of cases are caused by primary hyperparathyroidism (usually a benign parathyroid adenoma) or malignancy. Hallmark symptoms are classically summarized as ``stones, bones, abdominal groans, and psychiatric moans'': kidney stones and frequent urination; bone pain and osteoporosis; constipation, nausea, and pancreatitis; and fatigue, muscle weakness, confusion, and depression. ECG characteristically shows a shortened QT interval. Severe hypercalcemia is an emergency managed with intensive intravenous saline hydration, calcitonin, and bisphosphonates to inhibit bone resorption.
\synonyms{High Blood Calcium; Hypercalcaemia; Elevated Serum Calcium}

\medterm{Hypercalcemia Of Malignancy} A high blood-calcium level caused by cancer. It may occur because the cancer releases substances that raise calcium or damages bone. Severe cases can cause thirst, vomiting, weakness, confusion, kidney injury, or abnormal heart rhythms and require urgent treatment.

\medterm{Hypercalciuria} Excessive urinary calcium excretion, which can promote calcium kidney stones and, in some settings, bone mineral loss; evaluation considers diet, medications, kidney function, and endocrine causes.

\medterm{Hypercapnia} An abnormally high level of carbon dioxide in the blood, usually because breathing is not removing enough of it. It may cause headache, sleepiness, confusion, flushed skin, tremor, or coma and can occur with severe lung or muscle disease.

\medterm{Hypercementosis} Excessive deposition of cementum around a tooth root, producing a bulbous root contour that can complicate extraction but is often asymptomatic.

\medterm{Hypercholesterolaemia} An abnormally high level of cholesterol in the blood, especially low-density lipoprotein cholesterol. Over time, it can contribute to fatty buildup in arteries and increase the risk of heart attack or stroke.

\medterm{Hypercholesterolemia} A high level of cholesterol, especially LDL cholesterol, in the blood. It usually causes no symptoms but can gradually narrow arteries and raise the risk of heart attack and stroke; causes may be inherited or related to other diseases, medicines, or diet.

\medterm{Hypercoagulable State} A condition in which blood clots more easily than usual, increasing the risk of deep-vein thrombosis, pulmonary embolism, stroke, or other clots. Causes may be inherited or acquired, including cancer, pregnancy, some medicines, inflammation, or prolonged immobility.

\medterm{Hypercyanotic Tetralogy} A sudden episode of worsening blue or gray color, rapid breathing, distress, or faintness in a child with Tetralogy of Fallot. It results from reduced blood flow to the lungs and is an emergency requiring immediate medical treatment.

\medterm{Hyperelastic Skin} Skin that stretches farther than usual and recoils abnormally; generalized hyperextensibility may occur in connective-tissue disorders such as Ehlers-Danlos syndromes.

\medterm{Hyperemesis} Excessive or persistent vomiting that can cause dehydration, abnormal body salts, weight loss, and poor nutrition. It may result from pregnancy, infection, medicines, digestive disease, brain illness, or other causes.

\medterm{Hyperemesis Gravidarum} Severe nausea and vomiting during pregnancy that causes dehydration, weight loss, or abnormal body salts and prevents adequate food or fluid intake. It often begins early in pregnancy and may require fluids, nutrition, medicines, and hospital care.

\medterm{Hyperemia} Increased blood flow to a tissue, causing redness and warmth; it may be an active physiologic response to exercise or inflammation or a reactive response after temporary ischemia.

\medterm{Hyperemia Grade} A graded clinical or endoscopic estimate of tissue redness and vascular congestion, used to describe the severity of inflammation or mucosal irritation.

\medterm{Hyperesthesia} Increased sensitivity to sensory stimulation—touch (tactile), pain (hyperalgesia), or
other stimuli—so that normally non-painful or mildly painful stimuli produce exaggerated
discomfort.

\medterm{Hyperextension} Movement that extends a joint beyond its usual neutral or anatomical range, potentially stretching or tearing ligaments, capsules, muscles, or nerves.

\medterm{Hyperfibrinogenemia} An abnormally high level of fibrinogen, a blood protein needed for clotting. It often rises with inflammation, infection, pregnancy, obesity, or estrogen exposure and may be associated with increased clotting risk.

\medterm{Hyperflexion} Excessive bending of a joint beyond its normal range, which can injure ligaments, discs, muscles, or other supporting structures.

\medterm{Hypergammaglobulinemia} An excess of antibody proteins in the blood. It may result from long-term infection, inflammation, liver disease, or an abnormal immune-cell growth; testing determines whether many antibody types or one type is increased.

\medterm{Hyperglycaemia} An abnormally high level of glucose in the blood. It may cause thirst, frequent urination, tiredness, or blurred vision and, if severe or prolonged, can damage organs or cause an emergency.

\medterm{Hyperglycemia} A higher-than-normal level of glucose in the blood. Severe short-term elevation can cause excessive urination, dehydration, and dangerous changes in body chemistry. Long-term elevation can damage the eyes, kidneys, nerves, heart, and blood vessels.

\medterm{Hyperglycinemia} An abnormally high glycine concentration in blood, most notably in nonketotic hyperglycinemia, a disorder of glycine cleavage that can cause neonatal hypotonia, seizures, and developmental impairment.

\medterm{Hypergonadotropic Hypogonadism} Reduced testicular or ovarian hormone production caused by gonad failure, with increased pituitary signals and possible infertility, absent periods, or delayed sexual development.

\medterm{Hypergraphia} Excessive, compulsive, or unusually prolonged writing or drawing, sometimes associated with temporal-lobe epilepsy, mania, or other neurologic or psychiatric conditions.

\medterm{Hyperhidrosis} Sweating that is excessive for the temperature or activity. It may mainly affect the palms, soles, underarms, or face without another illness, or may be caused by medicines, hormone disorders, infection, or cancer.

\medterm{Hyperhomocysteinemia} An elevated blood homocysteine concentration associated with genetic enzyme defects, folate or vitamin B12 deficiency, kidney disease, medicines, and increased vascular or pregnancy-related risk in selected settings.

\medterm{Hypericum Perforatum} St. John’s wort, a plant used in some complementary products for depressive symptoms; it can cause drug interactions by inducing hepatic enzymes and transporters and should not be combined casually with medicines.

\medterm{Hyper-IgE Syndrome} A rare inherited immune disorder marked by very high levels of immunoglobulin E, eczema, and repeated skin or lung infections. Some forms also cause unusual bones, teeth, or blood vessels.

\medterm{Hyper-IgM Syndrome} See \textit{Hyper-IgM Syndromes}.

\medterm{Hyper-IgM Syndromes} Inherited immune disorders in which the body makes normal or high amounts of IgM antibodies but too little IgG, IgA, or IgE. People may develop repeated bacterial, viral, or fungal infections because antibody defenses do not mature normally.

\medterm{Hyperimmunization} Repeated or intensive exposure to an antigen, such as through vaccination or laboratory immunization, intended to produce a high antibody level. The term does not by itself describe an illness or harmful “over-vaccinated” state.

\medterm{Hyperimmunoglobulin D Syndrome} A hereditary autoinflammatory disorder, usually caused by mevalonate kinase deficiency, characterized by recurrent fever, enlarged lymph nodes, abdominal symptoms, rash, and joint or muscle pain; it is also called HIDS or MKD.

\medterm{Hyperimmunoglobulin E Syndrome} A rare primary immunodeficiency characterized by very high IgE, eczema, recurrent skin and lung infections, and sometimes skeletal or dental abnormalities.

\medterm{Hyperinsulinemia} An abnormally high insulin level in the blood, often caused by insulin resistance, excess insulin medicine, or rarely an insulin-producing tumor. It may occur with normal or low glucose and requires interpretation with glucose results.

\medterm{Hyperinsulinism} Excessive insulin secretion or action causing recurrent or persistent hypoglycemia; causes include congenital disorders, tumors, medications, and critical illness.

\medterm{Hyperkalaemia} Alternate British spelling; see \textit{Hyperkalemia}.

\medterm{Hyperkalemia} A dangerous elevation of potassium in the blood (greater than $5.0\,\text{mEq/L}$ or $\text{mmol/L}$), representing a medical emergency due to its disruptive effect on cardiac electrical conduction. Most commonly caused by acute or chronic kidney failure, potassium-sparing medications (such as ACE inhibitors, ARBs, or spironolactone), cellular destruction (burns, rhabdomyolysis, tumor lysis), or acidosis. While mild cases may cause muscle weakness, fatigue, or numbness, severe elevation causes progressive ECG abnormalities: tall peaked T waves, prolonged PR interval, widened QRS complex, sine-wave pattern, and fatal cardiac arrest. Emergency treatment requires intravenous calcium gluconate to stabilize the heart muscle, insulin with glucose to shift potassium into cells, and diuretics, binding agents, or urgent hemodialysis to remove potassium from the body.
\synonyms{High Blood Potassium; Hyperkalaemia; Elevated Serum Potassium}

\medterm{Hyperkalemic Periodic Paralysis} An inherited muscle disorder causing repeated attacks of weakness or temporary paralysis. Attacks may follow rest after exercise, stress, or potassium-rich food and usually improve between episodes.

\medterm{Hyperkeratosis} Thickening of the outer keratin layer of skin or mucosa, producing callus, scale, or rough plaques in response to friction, inflammation, infection, inherited disorders, or other causes.

\medterm{Hyperkinesia} Excessive or involuntary movement, such as chorea, dystonia, tics, or tremor, resulting from abnormal activity in motor-control circuits or from drugs and other diseases.

\medterm{Hyperlipasaemia} An abnormally high blood lipase concentration, often associated with pancreatic inflammation but also seen with kidney failure, gastrointestinal disease, medications, and other conditions.

\medterm{Hyperlipidaemia} An abnormally high level of fats in the blood, including cholesterol or triglycerides. It may have no symptoms but can increase the risk of artery disease, heart attack, stroke, or pancreatitis.

\medterm{Hyperlipidemia} Too much cholesterol or triglyceride fat in the blood. It usually causes no symptoms but can build up in artery walls, increasing the risk of heart attack and stroke; inherited disorders and other illnesses, medicines, diet, or alcohol can contribute.

\medterm{Hyperlipoproteinemia} An excess of one or more lipoproteins in blood, producing high cholesterol or triglycerides and increasing risk of atherosclerosis, pancreatitis, or inherited lipid disease.

\medterm{Hyperlordosis} Exaggerated inward curvature of the lumbar spine, often producing a swayback appearance and sometimes associated with pain, muscular imbalance, obesity, pregnancy, or structural spinal disease.

\medterm{Hyperlordotic Posture} An exaggerated inward curvature of the lumbar spine, often producing increased anterior pelvic tilt and a prominent lower-back arch.

\medterm{Hyperlysinemia} An inherited metabolic abnormality with elevated lysine in blood caused by altered lysine degradation; many people have few symptoms, but some have developmental or neurologic findings.

\medterm{Hypermagnesemia} An abnormally high concentration of magnesium in the blood, often related to kidney failure or excessive magnesium intake; severe cases can cause weakness, low blood pressure, depressed reflexes, or arrhythmias.

\medterm{Hypermetropia} A common refractive error of the eye in which parallel rays of light enter the relaxed eye and focus behind the retina rather than directly upon it. This occurs because the eyeball is anatomically too short along its anteroposterior axis (axial hypermetropia) or because the refractive power of the cornea or crystalline lens is insufficient (curvature hypermetropia). While distant objects may remain relatively clear in mild cases, viewing close objects requires continuous ciliary muscle accommodation, leading to accommodative asthenopia (eyestrain, frontal headache, and ocular fatigue during close work). In young children, significant uncorrected hypermetropia can cause accommodative esotropia (inward squint) and amblyopia. Optical correction is achieved using convex (plus/converging) spherical lenses in spectacles or contact lenses, or via corneal refractive surgery.
\synonyms{Hyperopia; Farsightedness; Long-sightedness}

\medterm{Hypermobile Ehlers-Danlos Syndrome} A heritable connective-tissue disorder characterized mainly by generalized joint hypermobility, instability,
chronic musculoskeletal pain, and other variable features involving skin, autonomic function, or digestion.
Diagnosis is clinical and requires assessment for other connective-tissue disorders and related conditions.

\medterm{Hypermobile Joints} Joints that move farther than usual. This may cause no problems, or may lead to pain, repeated dislocations, sprains, fatigue, or joint instability.

\medterm{Hypermobility} Greater-than-usual movement of one or more joints. It may be a harmless trait, but when associated with pain, instability, dislocations, fatigue, or other features it may be part of a hypermobility spectrum disorder.

\medterm{Hypermobility Spectrum Disorder} A condition involving joint hypermobility together with symptoms such as pain, joint instability, soft-tissue injury, or fatigue, without meeting the criteria for another defined connective-tissue disorder.

\medterm{Hypermobility Syndrome} Excessive joint movement associated with pain, instability, fatigue, or other symptoms; evaluation distinguishes generalized joint hypermobility from a defined hypermobility spectrum disorder or hypermobile Ehlers-Danlos syndrome.

\medterm{Hypermotility} Abnormally rapid or excessive movement of an organ, especially the digestive tract. It may cause cramping, urgency, diarrhea, or poor absorption and can result from infection, medicines, inflammation, or other disorders.

\medterm{Hypernatraemia} An abnormally high level of sodium in the blood, usually because the body has too little water. It may cause thirst, weakness, confusion, seizures, or coma and requires careful correction.

\medterm{Hypernatremia} A high blood-sodium level, usually because the body has lost too much water or cannot obtain enough. It can cause intense thirst, weakness, confusion, seizures, or coma and must be corrected carefully.

\medterm{Hypernatremia Result} A high blood sodium result, usually reflecting too little free water relative to body sodium; severity and rate of onset determine neurologic risk and the urgency and speed of correction.

\medterm{Hyperopia} Alternate name; see \textit{Hypermetropia}.

\medterm{Hyperornithinemia} An elevated blood ornithine concentration caused by defects in ornithine metabolism or transport, including disorders that can produce retinal degeneration or neurologic disease.

\medterm{Hyperosmolar} Having more dissolved particles than another fluid. When blood becomes too concentrated, water moves out of cells and can cause dehydration, confusion, or other nervous-system problems.

\medterm{Hyperosmolar Hyperglycemic State} A life-threatening emergency in which very high blood glucose causes severe dehydration and concentrated body fluids, usually with little or no ketone buildup. It may cause extreme thirst, weakness, confusion, seizures, or coma and needs urgent hospital treatment.

\medterm{Hyperosmolar Non-Ketotic Coma} A life-threatening complication of very high blood glucose, usually in type 2 diabetes, causing severe dehydration and sometimes confusion, seizures, or coma. It requires emergency hospital treatment.

\medterm{Hyperostosis} Abnormal excessive formation or thickening of bone, which may be localized or widespread and may occur in skeletal dysplasias, inflammation, endocrine disease, or medication effects.

\medterm{Hyperoxaluria} An excess of oxalate in the urine, which can combine with calcium to form kidney stones or calcium deposits. Severe inherited forms can progressively damage the kidneys.

\medterm{Hyperoxia} Exposure to an abnormally high oxygen concentration or tissue oxygen level, which can cause oxygen toxicity, absorption atelectasis, retinal injury in premature infants, or lung injury depending on dose and duration.

\medterm{Hyperparakeratosis} Thickening of the outer skin layer in which cells retain their nuclei instead of fully maturing. It may produce rough scale and can result from irritation, inflammation, infection, or an inherited skin disorder.

\medterm{Hyperparathyroidism} Excessive secretion of parathyroid hormone leading to hypercalcemia. Classified as primary (parathyroid adenoma or hyperplasia), secondary (chronic kidney disease, vitamin D deficiency), or tertiary (autonomous PTH secretion after prolonged secondary hyperparathyroidism).

\medterm{Hyperpathia} An unusually strong and sometimes delayed pain response to a painful stimulus. It may occur after nerve injury or with disorders that change how the nervous system processes pain.

\medterm{Hyperphagia} An unusually strong or excessive appetite. It can occur with diabetes, some genetic or brain disorders, certain medicines, or mental health conditions.

\medterm{Hyperphosphatemia} A high phosphate level in the blood, most often caused by kidney disease. It can disturb calcium balance, weaken bones, and contribute to calcium deposits in blood vessels or soft tissues.

\medterm{Hyperpigmentation} Darkening of an area of skin because the body makes or stores more pigment there. Causes include sun exposure, inflammation, injury, medicines, hormone disorders, and some skin or internal diseases.

\medterm{Hyperpituitarism} Excess production of one or more hormones from the pituitary gland, usually because of a benign pituitary tumor. The effects depend on the hormone involved and may affect growth, menstrual cycles, fertility, metabolism, or stress responses.

\medterm{Hyperplasia} An increase in the number of cells in a tissue or organ, which may make it larger. It can be a normal response or result from disease, hormones, or other stimuli.

\medterm{Hyperplastic Polyp} A usually noncancerous overgrowth of lining cells, often in the colon or rectum, whose cancer risk depends on size, number, and location.

\medterm{Hyperpnea} Breathing that is deeper, and sometimes faster, than normal, increasing air movement and occurring with exercise, illness, metabolic changes, or anxiety.

\medterm{Hyperpnoea} Breathing that is deeper or faster than usual to meet increased demand for oxygen or remove extra carbon dioxide. It may occur during exercise, fever, anxiety, or lung, heart, or metabolic illness.

\medterm{Hyperpolarization} An electrical change that makes a nerve or muscle cell less likely to send a signal. It occurs when the inside of the cell becomes more electrically negative after certain ion channels open.

\medterm{Hyperprolactinaemia} An abnormally elevated concentration of prolactin in the blood. Causes include medicines, pregnancy, hypothyroidism, pituitary tumors, and other disorders.

\medterm{Hyperprolactinemia} An abnormally high level of the hormone prolactin in the blood. Causes include pregnancy, breastfeeding, medicines, an underactive thyroid, kidney disease, or a pituitary tumor. It may disrupt periods, fertility, sexual function, or milk production.

\medterm{Hyperproteinemia} An abnormally high total blood-protein concentration, usually from dehydration, chronic inflammation, or a monoclonal immunoglobulin disorder.

\medterm{Hyperpyrexia} An extreme elevation of body temperature, generally above 41.5 °C, that can result from severe infection, heatstroke, malignant hyperthermia, drug toxicity, or other life-threatening conditions.

\medterm{Hyperreactio Luteinalis} Bilateral ovarian enlargement with multiple theca-lutein cysts caused by an exaggerated response to human chorionic gonadotropin, usually during pregnancy or with trophoblastic disease.

\medterm{Hyperreflexia} Exaggerated or overactive deep-tendon (stretch) reflexes. It is a neurologic sign, not a diagnosis, and often reflects impaired descending upper-motor-neuron control; it may occur with spasticity, clonus, or an extensor plantar response. Distribution and associated findings help distinguish central nervous-system disease from systemic or toxic causes.

\medterm{Hypersensitivity} An exaggerated or inappropriate immune response to an antigen, ranging from a mild reaction to
anaphylaxis.

\medterm{Hypersensitivity Pneumonitis} Lung inflammation caused by repeated breathing of substances such as mold, bird proteins, or other workplace dusts. It may cause cough, breathlessness, fever, or fatigue and can lead to permanent lung scarring if exposure continues.

\medterm{Hypersensitivity Reactions} Inappropriate or excessive immune responses causing tissue damage or disease.

\medterm{Hypersensitivity Type I} An immediate allergic reaction in which immunoglobulin E antibodies activate mast cells after exposure to an allergen. Symptoms can include itching, hives, wheezing, vomiting, swelling, or life-threatening anaphylaxis.

\medterm{Hypersensitivity Type II} An immune reaction in which IgG or IgM antibodies attach to targets on cells or tissues and recruit other immune mechanisms that injure them. Examples include some drug reactions and autoimmune blood-cell destruction.

\medterm{Hypersensitivity Type III} An immune reaction in which clusters of antibodies and foreign or self proteins settle in tissues or blood vessels, causing inflammation and injury. It can affect the skin, joints, kidneys, or blood vessels.

\medterm{Hypersensitivity Type IV} A delayed immune reaction driven by T cells rather than antibodies. It usually develops 48 to 72 hours after exposure and can cause contact dermatitis, some drug rashes, or inflammation around an infection.

\medterm{Hypersensitivity Vasculitis} Inflammation of small blood vessels, usually in the skin, often triggered by an infection, medicine, or immune reaction. It commonly causes raised purple spots or bruising on the legs.

\medterm{Hypersomnia} Being very sleepy during the day or sleeping much longer than usual even after enough sleep. It can be caused by sleep disorders, medicines, depression, substance use, or other illness; it is different from feeling tired without wanting to sleep.

\medterm{Hypersplenism} A condition in which an enlarged or overactive spleen removes blood cells too quickly. It may lower red cells, white cells, platelets, or all three and can cause anemia, infections, or easy bleeding.

\medterm{Hypertelorism} An unusually wide distance between the eyes because the bony eye sockets are farther apart than usual. It may be present from birth or occur as part of a genetic or developmental condition.

\medterm{Hypertension} Medical term; see \textit{High blood pressure}.

\medterm{Hypertension In Anesthesia} Abnormally high blood pressure during or around an operation. Pain, inadequate anesthesia, low oxygen, high carbon dioxide, shivering, bladder fullness, or stopping certain medicines may contribute and increase the risk of bleeding, heart injury, or stroke.

\medterm{Hypertensive} Having blood pressure that remains above the accepted range for the person's age, health, and measurement setting, increasing risk of heart, brain, kidney, or vessel disease.

\medterm{Hypertensive Cerebellar Hemorrhage} Bleeding into the cerebellum caused by rupture of a vessel damaged by high blood pressure.

\medterm{Hypertensive Crisis} A dangerously high blood pressure that needs prompt assessment. If it is causing new organ injury, such as stroke, heart failure, kidney injury, or vision loss, it is an emergency requiring immediate hospital treatment.

\medterm{Hypertensive Emergency} A severe rise in blood pressure that is causing new injury to the brain, heart, kidneys, eyes, or major blood vessels. Symptoms may include chest pain, breathlessness, weakness, confusion, seizures, or vision loss and require immediate hospital care.

\medterm{Hypertensive Encephalopathy} Brain dysfunction caused by a sudden, severe rise in blood pressure. Headache, confusion, seizures, vision changes, or reduced consciousness may occur and require immediate blood-pressure treatment.

\medterm{Hypertensive Heart Disease} Left ventricular hypertrophy and diastolic/systolic dysfunction resulting from chronic
hypertension. LVH is an independent cardiovascular risk factor and precursor to heart
failure. Diagnosis by ECG (voltage criteria) or echocardiography (wall thickness >1.1
cm). Treatment of hypertension can regress LVH.

\medterm{Hypertensive Intracerebral Hemorrhage} Bleeding into brain tissue after a blood vessel weakened by high blood pressure breaks, potentially causing sudden headache, weakness, confusion, unconsciousness, or death.

\medterm{Hypertensive Nephropathy} Kidney damage caused or worsened by long-term high blood pressure. It may lead to protein in the urine and gradually reduced kidney function; diabetes, other kidney diseases, and inherited risk can cause similar findings and must be considered.

\medterm{Hypertensive Nephrosclerosis} Scarring and narrowing of small kidney blood vessels caused by long-term high blood pressure. It can gradually reduce kidney function and may cause protein in the urine; controlling blood pressure helps slow progression.

\medterm{Hypertensive Retinopathy} Damage to the blood vessels and light-sensitive tissue at the back of the eye caused by high blood pressure. It may include narrowed vessels, bleeding, swelling, or fluid leakage. Severe changes can threaten vision and may indicate dangerously high blood pressure.

\medterm{Hypertensive Urgency} A markedly high blood pressure without clear evidence of new organ injury. It still requires prompt medical review, especially when symptoms are present; treatment is adjusted carefully rather than lowered suddenly without supervision.

\medterm{Hyperthermia} Hyperthermia is an abnormally high body temperature caused by excessive heat production or impaired heat loss, without a change in the hypothalamic temperature set point. Causes include heat exposure, exertion, medicines, endocrine disorders, infection-related illness, and impaired sweating. Severe hyperthermia can cause confusion, seizures, organ injury, or death and requires urgent treatment.

\medterm{Hyperthermia For Treating Cancer} A cancer treatment that deliberately heats a tumor or part of the body, sometimes with radiation or medicines, to damage cancer cells or make other treatment work better.

\medterm{Hyperthermia In Anesthesia} An abnormally high temperature during or after anesthesia. Causes include infection, drug reactions, thyroid storm, serotonin syndrome, and malignant hyperthermia, a rare emergency with muscle rigidity, rising carbon dioxide, and dangerous temperature elevation.

\medterm{Hyperthermic Intraperitoneal Chemotherapy} Heated cancer medicine circulated through the abdominal cavity during surgery after visible tumor is removed. It is used for selected cancers that have spread across the abdominal lining and can cause infection, bleeding, organ injury, or other serious complications.

\medterm{Hyperthyroidism} A condition in which the thyroid makes too much thyroid hormone. It may cause weight loss, sweating, heat intolerance, shaking, anxiety, diarrhea, or a fast or irregular heartbeat. Common causes include Graves disease and overactive thyroid nodules.

\synonyms
Thyrotoxicosis; Overactive Thyroid

\medterm{Hyperthyroidism} A state in which excess thyroid hormone effect causes symptoms such as heat intolerance, sweating, tremor, weight loss, or rapid heartbeat. It may result from increased thyroid production or other causes of thyrotoxicosis.

\synonyms
Thyrotoxicosis; Thyroid Hyperfunction

\medterm{Hyperthyroxinemia} An abnormally high concentration of circulating thyroxine, caused by hyperthyroidism, thyroiditis, medication, binding-protein changes, or assay interference; symptoms and TSH clarify the cause.

\medterm{Hypertonia} Abnormally increased muscle tightness or resistance when a limb is moved. It may cause stiffness, abnormal posture, spasms, or limited movement and can result from brain, spinal-cord, or other nervous-system disorders.

\medterm{Hypertonic Saline} Saline containing more sodium than normal blood, commonly given through a vein for severe symptoms from dangerously low blood sodium or selected urgent conditions. Close monitoring is required because correcting sodium too quickly can cause serious brain injury.

\medterm{Hypertonic Saline Inhalation} Inhaling a salty mist through a nebulizer to draw water into airway mucus and make it easier to clear. It may help some people with cystic fibrosis or bronchiectasis, but can briefly narrow the airways and should be used as instructed.

\medterm{Hypertonic Solution} A solution with greater effective osmotic concentration than the fluid with which it is compared, causing water to move out of cells across a semipermeable membrane.

\medterm{Hypertrichosis} Excessive hair growth in areas not usually controlled by male hormones. It can be present from birth or develop later because of inherited conditions, medicines, poor nutrition, hormone or metabolic disease, or rarely cancer. The cause should be assessed when growth is new or widespread.

\medterm{Hypertriglyceridemia} An abnormally high level of triglycerides, a type of fat, in the blood. Causes include diabetes, obesity, alcohol, some medicines, and inherited disorders. Very high levels increase the risk of acute pancreatitis and require medical management.

\medterm{Hypertrophic} Relating to hypertrophy, an enlargement of cells or tissue that increases the size of an organ or structure.

\medterm{Hypertrophic Cardiomyopathy} A disease in which the heart muscle, especially the left ventricle, becomes
abnormally thick and stiff. It may obstruct blood flow, impair filling, cause chest pain or breathlessness, and increase
the risk of arrhythmia and sudden death.

\synonyms
HCM; Hypertrophic Obstructive Cardiomyopathy (HOCM)

\medterm{Hypertrophic Cardiomyopathy Genetics} Many cases result from an inherited change in a gene for a heart-muscle protein, commonly MYH7 or MYBPC3. The condition can run in families, so genetic counseling and screening of close relatives may be recommended.

\medterm{Hypertrophic Cardiomyopathy Treatment} Management may include medicines to relieve symptoms, treatment of
arrhythmias or obstruction, an implantable defibrillator for selected patients, and surgery or ablation when indicated.
Treatment is guided by symptoms, anatomy, risk, and specialist assessment.

\medterm{Hypertrophic Gastritis} Thickening of the gastric mucosal folds caused by glandular overgrowth or inflammation, as in Ménétrier disease; protein loss, pain, nausea, and increased cancer risk may occur.

\medterm{Hypertrophic Myocardiopathy} Abnormal thickening of the heart muscle, often inherited. The thickened muscle may become stiff, obstruct blood leaving the heart, cause chest pain or breathlessness, or trigger dangerous abnormal rhythms.

\medterm{Hypertrophic Pyloric Stenosis} Thickening of the muscle at the outlet of the stomach, usually in young infants. It causes increasingly forceful vomiting after feeds, continued hunger, dehydration, and poor weight gain and is usually corrected with surgery.

\medterm{Hypertrophic Scar} A raised, thick scar that stays within the boundaries of the original wound. It forms when the body
makes too much scar tissue during healing, often after burns, surgery, or injury. Unlike a keloid,
it does not grow beyond the wound edges and may become flatter over time.

\medterm{Hypertrophy} Enlargement of a tissue or organ because its individual cells become larger. It may develop in response to increased workload, hormones, disease, or other stimulation.

\medterm{Hypertropia} Vertical misalignment of the eyes in which one eye deviates upward relative to the other, causing double vision, eye strain, or abnormal head posture.

\medterm{Hyperuricemia} An abnormally high concentration of uric acid in the blood, which may be asymptomatic or contribute to gout and uric-acid kidney stones.

\medterm{Hyperventilating} Breathing faster or deeper than the body needs, lowering carbon dioxide in the blood. It may cause light-headedness, tingling, chest tightness, or fainting and can result from anxiety or serious illness.

\medterm{Hyperventilation} Breathing more deeply or rapidly than the body needs, causing carbon dioxide levels in the blood to fall. It may cause dizziness, tingling, chest tightness, or fainting and can result from anxiety or serious heart, lung, or metabolic illness.

\medterm{Hyperventilation Syndrome} Repeated overbreathing, often during stress or panic, that lowers carbon dioxide in the blood and causes breathlessness, chest tightness, dizziness, tingling, palpitations, or hand and foot spasms. Other heart and lung causes must be considered.

\medterm{Hypervigilance} A state of excessive alertness to possible threat, often accompanied by scanning, startle, irritability, and sleep difficulty in anxiety, trauma-related disorders, stimulant use, or medical illness.

\medterm{Hyperviscosity Syndrome} A set of symptoms and signs caused by blood that is abnormally thick, usually from very high levels of an immunoglobulin paraprotein or of red or white blood cells. It can complicate Waldenström macroglobulinemia, multiple myeloma, polycythemia, leukemia, or severe infection. Features include bleeding from the nose or gums, blurred or lost vision from retinal changes, headache, dizziness, confusion, stroke-like symptoms, or heart failure; severe cases are an emergency managed with plasma exchange while the underlying disease is treated.

\medterm{Hypervitaminosis A} Illness caused by too much vitamin A, usually from supplements or medicines. It may cause headache, nausea, dry skin, hair loss, liver injury, bone pain, or birth defects during pregnancy.

\medterm{Hypervitaminosis D} Illness caused by too much vitamin D, leading to high blood calcium. It may cause nausea, vomiting, thirst, frequent urination, weakness, kidney stones, confusion, or kidney damage.

\medterm{Hypervolaemia} An excess of blood or body fluid. It may cause rapid weight gain, swelling, breathlessness, raised neck veins, or high blood pressure, especially with heart, kidney, or liver disease.

\medterm{Hypervolemia} Too much fluid in the bloodstream or body tissues. It may cause rapid weight gain, swelling, breathlessness, raised neck veins, or high blood pressure and can occur with heart, kidney, or liver disease.

\medterm{Hypha} A branching tubular filament that forms the basic structural unit of most molds and many fungi; networks of hyphae form a mycelium and may invade tissue in fungal disease.

\medterm{Hyphema} Blood collecting in the front part of the eye, usually after an injury. It may cause pain, blurred vision, light sensitivity, or increased eye pressure and requires prompt eye assessment to protect vision.

\medterm{Hypnagogic} Occurring during the transition from wakefulness to sleep. Hypnagogic sensations or hallucinations can be benign, but frequent or distressing episodes may occur with narcolepsy, sleep deprivation, medicines, or other sleep and neurologic disorders.

\medterm{Hypnagogic Hallucinations} Vivid sights, sounds, sensations, or feelings that occur while a person is falling asleep. They are often harmless but may be linked to sleep deprivation, narcolepsy, medicines, or other sleep disorders.

\medterm{Hypnic Headache} A rare primary headache that repeatedly awakens a person during sleep, usually in later life and often at a similar time each night.

\medterm{Hypnogram} A chart showing when a person is awake and in each sleep stage during a recording. It is commonly made from a sleep study and helps show sleep timing, awakenings, and stage patterns.

\synonyms
Sleep Architecture Diagram; Polysomnographic Hypnogram

\medterm{Hypnosis} A state of focused attention and deep relaxation in which a person may respond more readily to guided suggestions. It is not sleep or loss of control and should be used by a trained professional when offered as part of care.

\medterm{Hypnotherapy} Hypnotherapy uses focused attention and guided suggestion as an adjunct to care for selected problems such as pain, anxiety, or unwanted habits.

\medterm{Hypnotic} A medicine that helps a person fall asleep, stay asleep, or improve sleep. It is used mainly for short-term insomnia because some hypnotics can cause next-day drowsiness, falls, dependence, breathing problems, or dangerous effects with alcohol and other sedatives.

\medterm{Hypoactive Sexual Desire Disorder} Persistently low or absent sexual interest that causes personal distress or relationship difficulty. Desire can be affected by medicines, hormones, illness, pain, mood, stress, trauma, or relationship factors, so assessment looks for the underlying cause.

\medterm{Hypoadrenalism} Insufficient adrenal-cortex hormone production, especially cortisol and sometimes aldosterone; primary disease can cause fatigue, weight loss, low blood pressure, hyperpigmentation, and adrenal crisis.

\medterm{Hypoalbuminemia} An abnormally low blood albumin concentration caused by reduced production, inflammation, kidney or gastrointestinal loss, or severe malnutrition; it can contribute to edema and alter drug binding.

\medterm{Hypoalphalipoproteinemias} Disorders with low HDL-related apolipoprotein levels, which may impair reverse cholesterol transport and increase cardiovascular risk depending on the cause and overall lipid profile.

\medterm{Hypoblast} The inner cell layer of the early embryonic disc that contributes to the yolk-sac lining and helps establish the embryonic body axes during gastrulation.

\medterm{Hypocalcaemia} Alternate British spelling; see \textit{Hypocalcemia}.

\medterm{Hypocalcemia} An abnormally low level of calcium in the blood (total calcium below $8.5\,\text{mg/dL}$ or $2.1\,\text{mmol/L}$). Calcium is essential for nerve transmission, muscle contraction, and cardiac electrical stability. Low levels cause neuromuscular irritability, presenting with perioral numbness, fingertip tingling, painful muscle cramps, tetany, and hallmark physical signs: Chvostek's sign (facial twitching on tapping the facial nerve) and Trousseau's sign (carpal spasm during blood pressure cuff inflation). Severe cases cause seizures, airway spasm (laryngospasm), and a prolonged QT interval on ECG. Common causes include hypoparathyroidism (especially after thyroid surgery), severe vitamin D deficiency, chronic kidney disease, low magnesium, and pancreatitis. Severe cases require urgent intravenous calcium gluconate.
\synonyms{Low Blood Calcium; Calcium Deficiency; Hypocalcaemia}

\medterm{Hypocapnia} An abnormally low level of carbon dioxide in the blood, usually because a person is breathing too fast or deeply. Causes include anxiety, pain, low oxygen, lung disease, pregnancy, or excessive ventilator support.

\medterm{Hypochlorite} A chemical form of chlorine used in disinfectants and bleach. It can kill many germs, but swallowing it or getting concentrated solutions in the eyes, lungs, or skin can cause serious injury.

\medterm{Hypocholesterolemia} An unusually low blood cholesterol concentration caused by malnutrition, malabsorption, hyperthyroidism, chronic illness, liver disease, or inherited lipid disorders.

\medterm{Hypochondria} An older term for excessive fear of illness. Current clinical language generally uses illness anxiety disorder or health anxiety, depending on whether prominent physical symptoms are present.

\medterm{Hypochondriasis} Alternate historical term; see \textit{Illness Anxiety Disorder}.

\medterm{Hypochondroplasia} An inherited disorder of bone growth that causes short stature and shortened arms and legs, usually less severely than achondroplasia. People may have broad hands, increased inward curvature of the lower back, or limited elbow movement, while facial appearance and intelligence are often typical. It is usually caused by a change in the \textit{FGFR3} gene.

\synonyms
HCH; Mild Achondroplasia Variant

\medterm{Hypochromia} Reduced hemoglobin color in red blood cells, seen as increased central pallor on a blood smear and commonly associated with iron deficiency or thalassemia.

\medterm{Hypochromic Anaemia} Hypochromic anaemia is anemia in which red blood cells contain too little hemoglobin and appear unusually pale, commonly because of iron deficiency or thalassemia.

\medterm{Hypochromic Anemia} Anemia in which red cells contain too little hemoglobin and appear pale, most often because of iron deficiency, thalassemia, chronic disease, or sideroblastic disorders.

\medterm{Hypocitraturia} Low urinary citrate excretion, which reduces inhibition of calcium-crystal formation and increases risk of calcium kidney stones; potassium citrate may be used when clinically appropriate.

\medterm{Hypocupremia} An abnormally low copper concentration in blood or tissues, caused by inadequate intake, malabsorption, excess zinc, or inherited disorders, and potentially producing anemia, neutropenia, or neurologic disease.

\medterm{Hypodermic Needle} A hollow needle designed to inject medicines or withdraw fluid through the skin, with gauge and length selected according to the substance, route, body site, and patient.

\medterm{Hypodermis} The layer of fat and loose tissue beneath the skin. It cushions the body, stores energy, helps conserve heat, and contains blood vessels, nerves, and lymph channels.

\medterm{Hypodontia} A condition present from birth in which one or more permanent teeth fail to develop. It most often affects wisdom teeth, upper side incisors, or second premolars and may affect chewing, appearance, or tooth alignment.

\medterm{Hypoesthesia} Reduced sensitivity to touch, temperature, pain, or other sensations. It may result from nerve compression,
neuropathy, spinal disease, stroke, injury, or other neurologic conditions; sudden or spreading
numbness requires prompt assessment.

\medterm{Hypogalactia} Insufficient breast-milk production for an infant’s needs, caused by ineffective milk removal, hormonal or breast disorders, medications, illness, or other feeding difficulties.

\medterm{Hypogammaglobulinemia} Abnormally low levels of antibodies in the blood, which can lead to repeated infections, especially of the ears, sinuses, lungs, or bloodstream. It may be inherited or develop because of another disease or treatment.

\medterm{Hypogastric Region} The lower central area of the abdomen, just above the pubic bone. It overlies the bladder and parts of the bowel and reproductive organs; pain or a lump there may arise from urinary, intestinal, or pelvic disease.

\medterm{Hypoglossal Nerve} Cranial nerve XII, which supplies the muscles that move the tongue; a lesion causes tongue weakness, wasting, and deviation toward the affected side on protrusion.

\medterm{Hypoglossal Nerve Injury} Damage to the nerve that controls tongue movement. It may follow skull-base injury, surgery, a tumor, or infection and can cause tongue weakness, wasting, or deviation when the tongue is extended.

\medterm{Hypoglossal Nerve Palsy} Weakness or paralysis of the hypoglossal nerve (CN XII), leading to ipsilateral tongue
weakness, atrophy, and fasciculations. Causes include tumors at the skull base or in the
hypoglossal canal, surgical injury, or central lesions. Unilateral palsy causes the
tongue to deviate toward the affected side on protrusion.

\medterm{Hypoglycaemia} An abnormally low blood-glucose level. It may cause sweating, shaking, hunger, palpitations, confusion, seizures, or loss of consciousness and is especially associated with diabetes medicines, missed meals, or heavy alcohol use.

\medterm{Hypoglycemia} An abnormally low blood-glucose concentration that may cause sweating, tremor, palpitations, hunger,
confusion, seizures, or loss of consciousness. Causes include glucose-lowering medicines,
missed meals, alcohol, and critical illness. Treatment depends on the person’s ability to
swallow and the cause; severe or recurrent episodes require urgent assessment.

\medterm{Hypoglycemia In Diabetes} Blood glucose that falls too low in a person with diabetes, usually below 70 mg/dL. It may cause sweating, shaking, hunger, confusion, or fainting; severe episodes prevent self-treatment and require another person’s help.

\medterm{Hypoglycemia Unawareness} Failure to feel the usual warning signs of low blood glucose, such as shaking, sweating, or a racing heartbeat. It can occur after repeated low readings or long-standing diabetes and increases the risk of severe hypoglycemia.

\medterm{Hypoglycemic} Relating to an abnormally low blood glucose level, or describing a medicine that lowers blood glucose.

\medterm{Hypogonadism} Reduced function of the testes or ovaries, causing low sex-hormone levels and impaired production of sperm or eggs. It may lead to delayed puberty, absent periods, infertility, low sexual desire, hot flashes, reduced muscle or bone strength, or other symptoms.

\medterm{Hypogonadism In Men} Low testosterone caused by reduced testicular function or by a problem in the brain signals that control the testes. It may cause low sexual desire, erectile problems, infertility, low energy, reduced muscle, or weak bones.

\medterm{Hypogonadism In Women} Reduced ovarian hormone production, which may cause absent or irregular periods, hot flashes, vaginal dryness, infertility, low sexual desire, or loss of bone strength. Causes include ovarian failure, pituitary disease, some medicines, severe illness, or surgery.

\medterm{Hypogonadotropic Hypogonadism} Low production of sex hormones because the brain or pituitary gland does not send enough stimulating signals to the testes or ovaries. It may cause delayed puberty, absent periods, infertility, low sexual desire, or reduced testosterone. Causes include genetic conditions, pituitary disease, obesity, medicines, or severe illness.

\medterm{Hypohidrosis} Reduced or absent sweating. It can make it difficult for the body to cool itself and increase the risk of overheating. Causes include nerve disorders, skin damage, medicines, dehydration, and inherited conditions.

\medterm{Hypokalaemia} British spelling of Hypokalemia; see Hypokalemia.

\medterm{Hypokalemia} An abnormally low level of potassium in the blood, often caused by vomiting, diarrhea, diuretics, hormone disorders, or movement of potassium into cells. It may cause weakness, cramps, constipation, or dangerous heart-rhythm changes.

\medterm{Hypokalemic Periodic Paralysis} An inherited disorder causing repeated attacks of muscle weakness or paralysis when blood potassium falls. Attacks often begin after rest following exercise or after a high-carbohydrate meal, while strength may be normal between attacks.

\medterm{Hypokinesia} Reduced or slowed movement. It may cause small steps, less facial expression, reduced arm swing, or small handwriting and commonly occurs in Parkinson disease or other disorders affecting movement control.

\medterm{Hypolipidemic Agents} Medicines that lower cholesterol or triglycerides, including statins, ezetimibe, fibrates, bile-acid sequestrants, niacin, and selected injectable therapies.

\medterm{Hypolipoproteinemia} An abnormally low concentration of one or more blood lipoproteins, caused by inherited disorders, malabsorption, malnutrition, liver disease, or other systemic illness.

\medterm{Hypomagnesaemia} An abnormally low level of magnesium in the blood. It may cause tremor, muscle cramps, weakness, seizures, or abnormal heart rhythms and can result from diarrhea, poor intake, alcohol use, or certain medicines.

\medterm{Hypomagnesemia} A low level of magnesium in the blood, often caused by diarrhea, poor nutrition, alcohol use, kidney loss, or some medicines. It may cause tremor, cramps, weakness, seizures, or abnormal heart rhythms.

\medterm{Hypomania} A distinct period of unusually elevated, expansive, or irritable mood with increased energy or activity lasting at least four consecutive days. It may involve reduced need for sleep, rapid speech, racing thoughts, inflated self-esteem, or impulsive behavior; the change is observable by others but is not severe enough to cause marked impairment or require hospitalization and has no psychotic features.

\medterm{Hypomanic Episode} Alternate name; see \textit{Hypomania}.

\medterm{Hypomelanosis Of Ito} A rare condition causing swirled or streaky areas of lighter skin, often from birth. Some affected children also have seizures, developmental differences, eye problems, or abnormalities of the bones or nervous system.

\medterm{Hypomethylation} An abnormally low amount of methyl marking on DNA or other molecules. It can change which genes are active and may affect chromosome stability, development, inherited disorders, or cancer cells.

\medterm{Hypomnesia} Reduced ability to remember or recall information, which may occur with aging, neurologic disease, psychiatric conditions, medications, or other causes.

\medterm{Hyponatraemia} Alternate British spelling; see \textit{Hyponatremia}.

\medterm{Hyponatremia} An abnormally low level of sodium in the blood, often caused by excess water, fluid loss, medicines, heart or liver disease, kidney problems, or hormone disorders. It may cause nausea, headache, confusion, seizures, or coma; severe or rapidly developing cases are emergencies.

\medterm{Hypoparathyroidism} A disorder in which too little parathyroid hormone causes low blood calcium and high phosphate. It most often follows neck surgery but can be inherited or autoimmune. Tingling, muscle cramps, spasms, seizures, or abnormal heart rhythms may occur and need treatment.

\medterm{Hypopharyngeal Neoplasm} An abnormal growth in the lower part of the throat beside the voice box and esophagus. It may cause persistent sore throat, painful swallowing, difficulty swallowing, ear pain, or a neck lump and may be cancerous.

\medterm{Hypopharynx} The lower part of the pharynx behind and beside the larynx, extending from the oropharynx to the esophagus. It directs food toward the esophagus during swallowing.

\medterm{Hypophosphataemia} Alternate British spelling; see \textit{Hypophosphatemia}.

\medterm{Hypophosphatemia} An abnormally low level of phosphate in the blood (below $2.5\,\text{mg/dL}$ or $0.8\,\text{mmol/L}$). Because phosphate is essential for generating cellular energy (ATP) and oxygen delivery (2,3-DPG), its deficiency starves tissues of energy. It is classically triggered by refeeding syndrome (rapid carbohydrate refeeding in malnourished or alcoholic patients driving phosphate into cells), diabetic ketoacidosis recovery, severe malnutrition, or renal phosphate wasting. Symptoms include profound muscle weakness, respiratory failure from diaphragmatic weakness, acute muscle breakdown (rhabdomyolysis), hemolytic anemia, confusion, and seizures. Treatment consists of oral or intravenous phosphate replacement.
\synonyms{Low Blood Phosphate; Hypophosphataemia; Phosphate Deficiency}

\medterm{Hypophosphatemic Rickets} Rickets caused by low phosphate availability, often from renal phosphate wasting, leading to defective mineralization, bowed legs, bone pain, growth impairment, and dental abnormalities.

\medterm{Hypophysectomy} Surgical removal of all or part of the pituitary gland, usually performed through the nose and sphenoid sinus (transsphenoidal hypophysectomy) or rarely through an open craniotomy to treat pituitary tumors, relieve optic nerve compression, or control hormone hypersecretion.
\synonyms
Pituitary Resection; Transsphenoidal Adenomectomy

\medterm{Hypophysis} The pituitary gland, a small hormone-producing gland at the base of the brain. Its hormones influence growth, thyroid and adrenal activity, reproduction, milk production, and the body’s water balance.

\medterm{Hypophysis} An endocrine gland at the base of the brain that regulates growth, thyroid and adrenal function, reproduction, lactation, and water balance through its hormones.

\synonyms
Pituitary Gland; Master Endocrine Gland

\medterm{Hypophysitis} Inflammation of the pituitary gland that can cause headache, visual symptoms, and deficiency or excess of pituitary hormones; it may be autoimmune, related to pregnancy, or triggered by immune-checkpoint therapy.

\medterm{Hypopigmentation} Reduced skin pigmentation producing paler areas than the surrounding skin, caused by decreased melanin, reduced melanocyte function or number, inflammation, injury, infection, or inherited disorders.

\medterm{Hypopituitarism} Reduced production of one or more hormones from the pituitary gland because of disease, injury, surgery, radiation, or a nearby brain disorder. It may affect growth, thyroid or adrenal function, periods, fertility, milk production, or water balance.

\synonyms
Pituitary Insufficiency; Panhypopituitarism

\medterm{Hypoplasia} Incomplete or underdevelopment of a tissue, organ, or body structure, resulting in
reduced size or functional capacity. It may be congenital, developmental, genetic, or
acquired and ranges from a mild anatomical variant to severe organ dysfunction.
Assessment considers the timing of onset, associated malformations, tissue function, and whether the process is isolated or syndromic.

\medterm{Hypoplastic Left Heart Syndrome} A serious birth defect in which the left side of the heart and the aorta are too small to pump enough blood to the body. Newborns need urgent specialist treatment, often involving staged surgery or transplantation.

\medterm{Hypoploidy} A chromosome number below the normal diploid complement, caused by loss of one or more chromosomes; in tumor cytogenetics it may have prognostic significance.

\medterm{Hypopnea} A period of unusually shallow or reduced breathing, often during sleep. Repeated episodes can lower blood oxygen, disturb sleep, and contribute to obstructive sleep apnea or other breathing disorders.

\medterm{Hypoproteinaemia} Hypoproteinaemia is an abnormally low level of proteins in the blood, which may result from malnutrition, liver disease, kidney protein loss, or intestinal loss.

\medterm{Hypoproteinemia} An abnormally low amount of protein in the blood, often involving albumin, which can cause swelling and reflect poor nutrition, liver disease, kidney loss, or illness.

\medterm{Hypopyon} A visible white or yellow level of inflammatory material in the front chamber of the eye. It usually indicates severe eye inflammation or infection and may occur with pain, redness, light sensitivity, and reduced vision.

\medterm{Hyposmia} A reduced ability to smell, without complete loss of smell. It can follow a cold or COVID-19, nasal blockage, head injury, medicines, aging, or a neurologic disorder.

\medterm{Hypospadias} A birth condition in which the opening of the urethra is on the underside of the penis rather than at its tip. The penis may curve, and the foreskin may be incomplete; severity depends on the opening’s location and other abnormalities.

\medterm{Hypospadias Repair} Surgery reconstructs the urethral opening and, when needed, straightens the penis in a child born with hypospadias; timing and technique depend on anatomy and surgeon assessment.

\medterm{Hypotension} Abnormally low blood pressure that may reduce blood flow to the brain and other organs. It can result from bleeding, dehydration, heart problems, severe infection, allergy, or medicines and may cause dizziness, fainting, confusion, or shock.

\medterm{Hypotension In Anesthesia} Low blood pressure during or soon after anesthesia. It may result from anesthetic medicines, blood loss, dehydration, heart problems, or an allergic reaction; the anesthesia team monitors and treats it according to the cause and severity.

\medterm{Hypotensive} Having or causing abnormally low blood pressure; symptoms may include dizziness, fainting, weakness, confusion, or shock depending on the cause and severity.

\medterm{Hypothalamic Amenorrhea} Absence of menstrual periods because the hypothalamus reduces its normal release of reproductive signals. It is often linked to low energy availability, excessive exercise, stress, illness, or weight loss and may cause infertility and bone loss.

\medterm{Hypothalamic Dysfunction} A problem affecting the brain region that helps control hormones, temperature, appetite, thirst, sleep, stress, and reproduction. Causes include injury, tumors, inflammation, infection, surgery, radiation, or inherited disease.

\medterm{Hypothalamic Hamartoma} A congenital, noncancerous mass of disorganized hypothalamic tissue that can cause gelastic seizures, precocious puberty, behavioral symptoms, or cognitive decline.

\medterm{Hypothalamic Hormone} A releasing or inhibiting hormone produced by the hypothalamus that regulates pituitary secretion, including signals controlling thyroid, adrenal, reproductive, growth, and lactation axes.

\medterm{Hypothalamic Neoplasm} An abnormal growth in or near the hypothalamus, a brain region controlling hormones, temperature, appetite, sleep, and water balance. Effects may include vision changes, hormone problems, headaches, seizures, or altered growth.

\medterm{Hypothalamic Tumor} Alternate name; see \textit{Hypothalamic Neoplasm}.

\medterm{Hypothalamic-Pituitary-Adrenal Axis} A network linking the hypothalamus, pituitary gland, and adrenal glands to regulate stress hormones and other functions.

\synonyms
HPA Axis; Corticotrope-Adrenal System

\medterm{Hypothalamic-Pituitary-Adrenal Axis} The hormone system linking the brain’s hypothalamus and pituitary gland with the adrenal glands. It controls cortisol release, helping regulate the body’s response to stress, illness, blood pressure, metabolism, and immune activity.

\synonyms
HPA Axis; Neuroendocrine Stress Axis

\medterm{Hypothalamic-Pituitary-Adrenal Axis Physiology} The brain signals the pituitary gland, which signals the adrenal glands to release cortisol. Cortisol normally follows a daily pattern and feeds back to reduce further signals when enough is present.

\medterm{Hypothalamic-Pituitary-Gonadal Axis Physiology} The hormone system linking the hypothalamus, pituitary gland, and ovaries or testes. Signals from the hypothalamus prompt the pituitary to release hormones that control ovulation and menstrual cycles, or testosterone production and sperm development.

\medterm{Hypothalamic-Pituitary-Thyroid Axis Physiology} The hormone system linking the hypothalamus, pituitary gland, and thyroid. The hypothalamus signals the pituitary to release thyroid-stimulating hormone, which prompts the thyroid to make hormones that regulate metabolism, temperature, and energy use.

\medterm{Hypothalamus} A small region at the base of the brain that helps control automatic body functions and hormone release through the pituitary gland. It regulates body temperature, hunger, thirst, sleep, stress, and parts of emotional behavior.

\medterm{Hypothermia} A dangerously low internal body temperature, usually below 35 degrees Celsius, caused by cold exposure or impaired heat production. It may cause shivering, confusion, drowsiness, slow breathing, or an abnormal heart rhythm. Severe cases require urgent medical rewarming.

\synonyms
Accidental Hypothermia; Low Core Body Temperature

\medterm{Hypothermia In Anesthesia} An abnormally low body temperature during or after anesthesia. Anesthetic medicines reduce the body’s ability to control heat and may allow heat loss during surgery; warming devices and temperature monitoring help prevent complications.

\medterm{Hypothyroidism} A clinical endocrine syndrome resulting from deficient production or action of thyroid hormones (thyroxine [$T_4$] and triiodothyronine [$T_3$]), causing a generalized slowing of basal metabolism. Primary hypothyroidism (over 95\% of cases, characterized by elevated serum TSH and low free $T_4$) most commonly results from autoimmune destruction (Hashimoto thyroiditis), surgical thyroidectomy, radioactive iodine ablation, or severe iodine deficiency; central hypothyroidism stems from pituitary or hypothalamic failure. Clinical hallmarks include profound fatigue, cold intolerance, weight gain, constipation, sinus bradycardia, dry skin, non-pitting dermal edema (myxedema), delayed relaxation of deep tendon reflexes, and cognitive slowing.
\synonyms{Underactive Thyroid; Hypothyroid State; Thyroid Deficiency}

\medterm{Hypotonia} Decreased muscle tone presenting as a
floppy or poorly supported posture. It may result from muscle, peripheral nerve,
neuromuscular-junction, spinal-cord, brain, genetic, or metabolic disorders and requires
evaluation when persistent or associated with weakness or developmental delay.

\medterm{Hypotonic Solution} A liquid containing fewer dissolved particles than the fluid inside cells. Water tends to move into cells when they are exposed to it, so medical use requires careful selection and monitoring.

\medterm{Hypotrichosis} Abnormally sparse or absent hair caused by inherited follicle disorders, autoimmune disease, endocrine or nutritional problems, medications, or scarring alopecia.

\medterm{Hypotrophy} Underdevelopment or reduced size of a tissue or organ that is present but smaller or less developed than usual. It may begin before birth or result from poor nutrition, reduced use, impaired blood flow, or disease and may affect function.

\medterm{Hypotropia} Vertical eye misalignment in which one eye deviates downward relative to the other, potentially causing diplopia, abnormal head posture, or suppression of vision.

\medterm{Hypoventilation} Breathing too slowly or shallowly to remove enough carbon dioxide. It can lower oxygen levels and cause sleepiness, headache, confusion, or respiratory failure; causes include medicines, lung disease, obesity, and muscle or nerve disorders.

\medterm{Hypovolaemia} Hypovolaemia is an abnormally low volume of circulating blood or body fluid, commonly caused by bleeding, dehydration, or major fluid loss.

\medterm{Hypovolemia} Hypovolemia is reduced effective circulating blood volume, usually from bleeding, fluid loss, or inadequate intake; it can cause tachycardia, hypotension, poor perfusion, kidney injury, and hypovolaemic shock.

\medterm{Hypovolemic Shock} Circulatory failure caused by insufficient circulating blood or fluid volume, leading to reduced
tissue perfusion. Causes include hemorrhage, severe dehydration, and fluid loss.

\medterm{Hypoxaemia} An abnormally low oxygen level in arterial blood. It can result from lung disease, breathing problems, poor blood flow, or other conditions and may cause shortness of breath, confusion, bluish skin, or organ injury.

\medterm{Hypoxanthine} A naturally occurring substance made when the body breaks down certain genetic-material building blocks. It is normally reused or changed into other compounds during purine metabolism.

\medterm{Hypoxemia} An abnormally low oxygen level in the blood. It may result from lung or heart disease, blocked airways, shallow breathing, anemia, poor circulation, or high altitude and can cause breathlessness, confusion, or bluish skin.

\medterm{Hypoxia} Inadequate oxygen availability at the tissue level. It may result from low arterial oxygen, anemia,
reduced blood flow, or impaired cellular oxygen use.

\medterm{Hypoxia And Hypoxemia} Hypoxemia is low oxygen in arterial blood, while hypoxia is inadequate oxygen at the tissue level. Causes include lung disease, airway obstruction, heart disease, anemia, shock, or high altitude; severe breathlessness, cyanosis, confusion, or collapse is urgent.

\medterm{Hypoxic Cell Injury} Damage to cells caused by too little oxygen. It may result from poor blood flow, low blood oxygen, anemia, or poisoning that prevents oxygen use; the outcome depends on how severe and prolonged the oxygen shortage is.

\medterm{Hypoxic Ischemic Encephalopathy} Brain dysfunction caused by insufficient oxygen and blood flow, usually around birth in newborns but also after cardiac arrest or severe respiratory failure. It can cause seizures, altered consciousness, and long-term neurologic disability.

\medterm{Hypoxic Pulmonary Vasoconstriction} Narrowing of small blood vessels in an area of the lung that receives too little oxygen. This normally redirects blood toward better-ventilated lung areas, but widespread or prolonged narrowing can strain the heart.

\medterm{Hypoxic-Ischemic Brain Injury} Brain damage caused by too little oxygen or blood flow. It may follow cardiac arrest, suffocation, shock, severe breathing failure, or very low blood glucose and can cause seizures, coma, memory problems, movement problems, or disability.

\medterm{Hypsarrhythmia} A very abnormal, disorganized pattern on an electroencephalogram, with large slow waves and spikes arising from many areas of the brain. It is strongly associated with infantile spasms, a serious seizure disorder in babies, and requires prompt specialist evaluation and treatment.

\synonyms
Hypsarrhythmic EEG Pattern; Chaotic Dysrhythmic EEG

\medterm{Hysterectomy} Surgical removal of the uterus. The cervix, fallopian tubes, and ovaries may or may not also be removed, depending on the procedure and clinical indication. Hysterectomy permanently ends menstruation and the ability to carry a pregnancy, but it does not cause immediate menopause when functioning ovaries remain. If the cervix remains, cervical cancer screening may still be needed. Indications include cancer, fibroids, severe bleeding, adenomyosis, endometriosis, chronic pelvic pain, and uterine prolapse.

\medterm{Hysterectomy And Surgical Menopause} Hysterectomy removes the uterus and ends menstruation and the ability to carry a pregnancy. Surgical menopause occurs when both ovaries are removed or ovarian function otherwise stops abruptly; hysterectomy alone does not necessarily cause menopause if the ovaries remain.

\medterm{Hysterosalpingo-Contrast Sonography} An ultrasound test using fluid or foam contrast to assess the uterine cavity and whether the fallopian tubes are open. It is often used during fertility evaluation without X-ray radiation.

\synonyms
HyCoSy; Fallopian Tubal Contrast Ultrasound

\medterm{Hysterosalpingo-Contrast-Sonography} An ultrasound-based technique for evaluating the uterine cavity and fallopian tube
patency. Sterile saline or contrast medium is injected through the cervix while
transvaginal ultrasound is performed.

\medterm{Hysterosalpingogram} An X-ray procedure used to evaluate the uterine cavity and fallopian tube patency. A
radiopaque contrast medium is injected through the cervix, and serial X-ray images are
taken as the contrast fills the uterine cavity and spills from the fallopian tubes.

\medterm{Hysterosalpingography} An X-ray test in which contrast liquid is placed through the cervix to outline the inside of the uterus and fallopian tubes. It can show blockage, an unusual uterine shape, fibroids, polyps, or scar tissue during infertility evaluation. Cramping, bleeding, infection, or contrast reaction can occur.

\medterm{Hysteroscope} A thin camera-equipped instrument passed through the cervix to inspect the uterine cavity and, when designed for operative use, remove polyps, treat bleeding, or perform other procedures.

\medterm{Hysteroscopic Metroplasty} A procedure that uses a small camera and instruments passed through the cervix to divide a wall of tissue separating part of the uterus. It may improve the uterine cavity in selected people with fertility problems or repeated pregnancy loss.

\medterm{Hysteroscopic Myomectomy} A procedure that removes selected fibroids bulging into the inside of the uterus. A camera and instruments pass through the cervix, avoiding an abdominal incision; the procedure may reduce heavy bleeding and improve fertility.

\medterm{Hysteroscopy} A procedure that passes a small camera through the cervix to inspect the inside of the uterus and, when needed, obtain tissue or treat conditions such as polyps, fibroids, or intrauterine adhesions.

\medterm{Hysterotomy} A surgical incision into the uterus performed during selected obstetric or gynecologic procedures. It may be used to deliver a baby or reach the uterus, with risks including bleeding and infection.

\medterm{Infant Hearing Loss} Congenital or early-acquired auditory impairment in infants, resulting from genetic mutations (such as connexin 26 mutations), congenital infections (cytomegalovirus, rubella), prematurity, hyperbilirubinemia, or ototoxic drug exposure. Universal newborn hearing screening enables early amplification and speech-language intervention.

\synonyms
Congenital Hearing Loss; Pediatric Auditory Impairment; Neonatal Deafness

\medterm{Juvenile Hypermobile Spectrum Disorder} A hypermobility spectrum disorder beginning in childhood or adolescence, with excessive joint movement plus pain, instability, soft-tissue injury, fatigue, or related symptoms.

\medterm{Lanugo Hair} Lanugo: the fine, soft, usually unpigmented hair covering a fetus and sometimes present on a premature newborn; it is normally shed before or shortly after birth.

\medterm{Mevalonate Kinase Deficiency} An inherited autoinflammatory disorder caused by changes in the MVK gene. It can cause repeated fever attacks, swollen lymph nodes, abdominal pain, rash, joint pain, or mouth ulcers; the milder form is often called hyperimmunoglobulin D syndrome.

\medterm{Minimally Invasive Heart Bypass Surgery} Coronary artery bypass graft (CABG) surgery performed through a small thoracotomy incision between the ribs, frequently on a beating heart without cardiopulmonary bypass, avoiding a traditional median sternotomy.

\synonyms
MIDCAB; Minimally Invasive Direct Coronary Artery Bypass; Less Invasive CABG

\medterm{Neonatal Hyperglycemia} An abnormally elevated blood glucose level in infants and neonates (typically defined as blood glucose $>145\text{ mg/dL}$ or $>8.0\text{ mmol/L}$). Common causes include extreme prematurity, intravenous dextrose overinfusion, sepsis, severe physiological stress, corticosteroid therapy, or transient neonatal diabetes mellitus.

\synonyms
Infantile Hyperglycemia; High Blood Glucose in Infants; Newborn Hyperglycemia

\medterm{Neonatal Hypertension} Sustained blood pressure elevation in infants and neonates exceeding age- and birthweight-adjusted norms. Frequently secondary to umbilical artery catheter-related renal artery thrombosis, congenital renal anomalies, bronchopulmonary dysplasia, or congenital heart defects.

\synonyms
Infant Hypertension; Neonatal High Blood Pressure; Infantile Hypertension

\medterm{Oral Herpes} A common viral infection of the oral mucosa and lips caused primarily by herpes simplex virus type 1 (HSV-1). Characterized by recurrent outbreaks of fluid-filled vesicles (cold sores or fever blisters) that ulcerate and crust, accompanied by prodromal tingling or burning.

\synonyms
Herpes Simplex Labialis; Cold Sores; Fever Blisters; Oral HSV-1 Infection

\medterm{Pediatric Hepatitis A} An acute viral infection of the liver caused by the hepatitis A virus (HAV) in children, transmitted predominantly through the fecal-oral route via contaminated food or water. While often self-limiting, supportive hydration is critical and routine pediatric vaccination provides durable protection.

\synonyms
Childhood Hepatitis A; Pediatric HAV Infection; Acute Viral Hepatitis A in Children

\medterm{Pediatric Hepatitis B} A viral liver infection caused by the hepatitis B virus (HBV) in pediatric patients, most commonly acquired perinatally from an infected mother or during early childhood. Chronic carrier states are highest when acquired in infancy; universal infant vaccination and post-exposure immunoprophylaxis prevent transmission.

\synonyms
Childhood Hepatitis B; Perinatal HBV Infection; Chronic Pediatric Hepatitis B

\medterm{Pediatric Hepatitis C} Infection with the hepatitis C virus (HCV) in infants and children, primarily acquired vertically through perinatal maternal-fetal transmission or through parenteral exposure. Modern direct-acting antiviral (DAA) regimens provide high cure rates in pediatric cohorts.

\synonyms
Childhood Hepatitis C; Pediatric HCV Infection; Perinatal Hepatitis C

\medterm{Pediatric Hypercholesterolemia} Elevated circulating total cholesterol or low-density lipoprotein (LDL) levels in children and adolescents. Causes include familial hypercholesterolemia, obesity, insulin resistance, endocrine disorders, or high-saturated-fat diets.

\synonyms
Childhood Hyperlipidemia; Pediatric Dyslipidemia; Familial Hypercholesterolemia in Children

\medterm{Pediatric Hypertension} Blood pressure in children and adolescents persistently at or above the 95th percentile for age, sex, and height. Secondary etiologies (such as renal parenchymal disease or aortic coarctation) are common, especially in young children.

\synonyms
Childhood Hypertension; Pediatric High Blood Pressure; Adolescent Hypertension

\medterm{Perineal Hemostat} A clamp used during procedures to compress a blood vessel or hold tissue and control bleeding.

\synonyms
Perineal Hemostatic Clamp; Pelvic Bleeding Forceps

\medterm{Physical Half-Life} The time it takes for half of a radioactive substance to lose its radioactivity. It helps determine the timing of scans and the amount of radiation exposure.

\medterm{Piles} Alternate name; see \textit{Hemorrhoids}.

\medterm{Psychological Homeostasis} The ongoing adjustment of thoughts, emotions, and behavior that helps a person maintain psychological stability while responding to stress or change. Severe disruption can contribute to anxiety, depression, or impaired daily functioning.

\medterm{Qualitative HCG Blood Test} A diagnostic serological laboratory test that detects the presence or absence of human chorionic gonadotropin (hCG) in maternal blood. Primarily utilized for early confirmation of pregnancy.

\synonyms
Qualitative Serum Pregnancy Test; Blood hCG Screen; Qualitative Beta-hCG

\medterm{Quantitative HCG Blood Test} A precise diagnostic laboratory assay that measures the exact serum concentration of beta-human chorionic gonadotropin. Essential for monitoring early pregnancy viability, evaluating suspected ectopic pregnancy or miscarriage, and tracking gestational trophoblastic disease or germ-cell neoplasms.

\synonyms
Quantitative Beta-hCG; Serial hCG Test; Serum hCG Level

\medterm{Sensorineural Hearing Loss} Hearing loss caused by damage to the sound-sensing cells or nerve fibers of the inner ear, the auditory nerve, or related hearing pathways. Sounds may seem muffled, and understanding speech---especially in noisy places---may be difficult; ringing in the ears (tinnitus) may also occur. Common causes include ageing, loud noise, inherited conditions, infections, head injury, and some medicines. It differs from conductive hearing loss, which results from a problem in the outer or middle ear, and from mixed hearing loss, which includes both types.

\medterm{Singultus} Alternate medical name; see \textit{Hiccups}.

\medterm{Social History} Information about a person's living situation, relationships, occupation, education, tobacco or substance use, diet, safety, and other social factors relevant to health and care.

\medterm{Surgical Hemostat} A clamping surgical instrument used to control bleeding by compressing a blood vessel or tissue.
\synonyms
Surgical Hemostatic Forceps; Mosquito Clamp

\medterm{Swine Flu} Alternate name; see \textit{H1N1 Influenza}.

\medterm{Urticaria} Alternate medical name; see \textit{Hives}.

\medterm{Vascular Hemostat} A surgical clamp used to temporarily occlude a blood vessel and control bleeding.
\synonyms
Vascular Clamp; Atraumatic Hemostasis Forceps

\medterm{Viral Hepatitis} Inflammation of the liver caused by one of the primary hepatitis viruses (A, B, C, D, or E) or other viral agents. Manifestations range from asymptomatic elevation of aminotransferases to acute jaundice, fulminant hepatic failure, or chronic hepatitis leading to cirrhosis.

\medterm{Viral Hepatitis Types} Five viruses that can inflame the liver. They differ in how they spread, whether infection can become long-term, the complications they cause, and whether vaccination or curative treatment is available.

\medterm{Vitamin K Deficiency Bleeding} Bleeding in a newborn or young infant caused by inadequate vitamin K-dependent clotting activity. It may occur early, classically, or late and can involve the brain or gastrointestinal tract.

\synonyms
VKDB; Neonatal Hypoprothrombinemia
